% ================================
% \section*{\centering \bf Part IV: Applications}
% \addcontentsline{toc}{section}{Part IV: Applications}
% ================================




\begin{figure}
    \centering
    % \includegraphics[width=\linewidth, height=12cm]{example-image-golden}
    % \includegraphics[width=0.99\linewidth]{figures/part4.png}
    \includegraphics[width=\linewidth]{figures/part4_opener_v2.pdf}
    
    \caption{Part IV at a glance. (1) Leave the lab: deployment tests spatial intelligence under latency, safety, sensor failure, and the cost of error. (2) Domain by domain: embodied agents, extended reality, geo-spatial and scientific settings, assistants, and further fields, each stressing a different axis. (3) A horizontal foundation: one shared layer beneath science, industry, and daily life, with gaps found in the field feeding back into it.}
    \label{fig:p4_main}
\end{figure}


The preceding parts built spatial intelligence as a set of capabilities; this part asks what those capabilities are worth once they leave the laboratory. Applications are not a victory lap over Parts II and III but a stress test: a method that reconstructs a scene cleanly on a benchmark must, in the field, do so under a latency budget, a safety constraint, a sensor that fails, or a cost of error measured in more than accuracy. What a domain demands of spatial intelligence, and what it reveals about the gaps that remain, is visible only at deployment.

This part is therefore organized by domain rather than by capability, because the interesting content lives in the delta each domain adds over a general treatment. The tour runs from embodied agents that must act in the physical world, through the extended-reality systems that render it back to a user in real time, to the geo-spatial and scientific settings that model it at the scale of planets and of cells, to the assistants that inhabit a person's own environment, and out into a long tail of further fields. Each stresses spatial intelligence along a different axis: real-time action, on-device perception, metric fidelity, human-centered interaction. Read together, they make the case that spatial intelligence is not one more application but a horizontal foundation on which applications across science, industry, and daily life are increasingly built.




\chapter{Spatial Intelligence for Embodied Systems }
\label{sec:embodied}

\FramedInitial{E}{mbodied} systems are the arena in which spatial intelligence stops being a matter of analysis and becomes a matter of consequence. An embodied agent does not merely describe a scene; it occupies one, and every perception it forms and every inference it draws is ultimately answerable to an action it must take in the physical world, under real time budgets, sensor noise, and the cost of being wrong. This chapter turns from the capabilities of the previous part to their deployment, examining how perception, reasoning, and generation are brought to bear inside agents that move, manipulate, and assist.

What makes this more than a restatement is that embodiment changes the problem. The same spatial competence looks different in a manipulator closing a grasp, a car threading traffic, a wearable inferring its user's intent, a drone holding position in wind, or a fleet of robots coordinating on a warehouse floor: each fixes its own sensors, latencies, safety margins, and notion of success. The chapter is organized by these domains rather than by capability, because the interesting content lives in the domain-specific delta, what each setting demands of spatial intelligence that a general treatment does not reveal.


\section{Robotic Manipulation}
\label{sec:p4_s9_1}



 % \RespondAuthor{ }
\subsection{Introduction}
\label{subsec:p4-sec9-1-introduction}

Robotic manipulation has advanced rapidly, driven by strides in perception, learning, and control.
Yet a significant gap remains between competent laboratory systems and robust, general-purpose robots capable of operating in unstructured settings.
A central bottleneck is that manipulation is fundamentally spatial: every decision (where to reach, how to establish contact, how to route a trajectory through clutter, or how to use a tool) hinges on perceiving, representing, and reasoning about three-dimensional structure under uncertainty.
Manipulation is thus a direct test of the \emph{spatial intelligence} developed throughout Parts~I--III: it demands that systems go beyond semantic recognition to build actionable 3D world models, ground tasks in geometric constraints, predict contact outcomes, and plan behaviors that respect the physical laws of the environment.
The perception, reasoning, and generation capabilities that such systems presuppose are surveyed in Chapter~\ref{chapter:3:sec:6}, Chapter~\ref{sec:reasoning}, and Chapter~\ref{sec:generation}; this section focuses not on re-deriving those capabilities but on how they are composed, grounded, and constrained for physical manipulation.

\subsection{System Overview}
\label{subsec:p4-sec9-1-system-overview}
A manipulation system maps raw multimodal observations and a task instruction into executable, constraint-satisfying actions, and this section organizes that mapping into four strategic \textbf{pivots}, each a fundamental shift in how systems bridge perception and physical action.
\textbf{(i) Grounding Action in Spatial Anchors} moves beyond raw pixels to vision priors and sparse spatial tokens (keypoints, dense correspondences) that stay consistent across viewpoints and instances;
\textbf{(ii) Inferring Functional Action Possibilities} shifts from object classification to extracting affordances and tool-use capabilities, bridging perception with contact-rich execution;
\textbf{(iii) Encoding Spatial Reasoning in Generalist Models} transitions from single-task policies to foundation models (VLMs/VLAs) that integrate explicit 3D representations and reasoning traces for cross-embodiment generalization; and
\textbf{(iv) Satisfying Geometric and Logical Constraints} advances from simple trajectory generation to constrained optimization and neuro-symbolic planning (TAMP), ensuring actions respect strict safety limits and long-horizon logical structure.
The four pivots are complementary and mutually reinforcing rather than strictly sequential: anchors and affordances structure the perceptual input, generalist models supply broad priors, and constraint-based planners check feasibility under the modeled geometric, physical, and logical assumptions. Figure~\ref{fig:p4_s9_1_overview} situates them as complementary capabilities within a closed-loop perception-to-action system, and the remainder of this section treats them in this order.

\begin{figure}[t]
    \centering
    \includegraphics[width=\linewidth]{figures/part4/section9/manipulation_overview.pdf}
    \caption{Spatial intelligence for robotic manipulation, illustrated by placing a mug in an open drawer: four complementary spatial capabilities link multimodal inputs to 6-DoF actions, and execution returns new observations for replanning. Evaluation bars and depth/point-cloud graphics are schematic, not measured results.}
    \label{fig:p4_s9_1_overview}
\end{figure}


\subsection{Pivotal Task-1: Grounding Action in Spatial Anchors}
\label{subsec:p4-sec9-1-pivotal-task-1-grounding-action-in-spatial-anchors}

This pivot reframes visuomotor learning from operating over raw pixels to reasoning through explicit \textbf{spatial anchors} (sparse semantic keypoints, dense correspondences, point tracks, or affordance-based contact points) that compress high-dimensional observations into spatial tokens tied to object- or task-centric reference frames and designed to remain comparatively stable across viewpoint, instance, and clutter. As Figure~\ref{fig:p4_s9_1_anchors} indicates, the approach follows a three-step logic: large-scale vision priors first \emph{produce} the anchors, policies are then \emph{conditioned} on them, and the payoff is read off as a \emph{generalization gap} between seen and unseen instances.

\begin{figure}
    \centering
    \includegraphics[width=0.98\linewidth]{figures/part4/section9/robotic_anchors.pdf}
    \caption{Grounding action in spatial anchors: (1) anchors $P_t$ produced or transferred from vision priors across views, (2) an anchor-conditioned policy, and (3) generalization measured as the seen--unseen success gap $\Delta_{\mathrm{gen}}$. Scenes, correspondences, and score bars are schematic, not measured results.}
    \label{fig:p4_s9_1_anchors}
\end{figure}

The dominant methodology leverages large-scale vision priors to generate, propagate, or retrieve these anchors, then trains standard visuomotor learners on the resulting structured inputs. A common strategy prescribes anchor points in a few reference frames and propagates them across large datasets via off-the-shelf correspondence models: P3-PO~\cite{levy2024p3poprescriptivepointpriors} propagates manually specified keypoints using dense features from DINO~\cite{dino}, DINOv2~\cite{dinov2}, or diffusion-based representations~\cite{tang2023emergentcorrespondenceimagediffusion,zhao2025high,zhang2023talefeaturesstablediffusion} (these visual representations are surveyed in Chapter~\ref{chapter:single_modality}). Relatedly, Robo-ABC~\cite{ju2024roboabcaffordancegeneralizationcategories} retrieves affordance contact points from an external memory and transfers them to novel scenes via dense semantic correspondence, while Track2Act~\cite{bharadhwaj2024track2act} learns long-term point tracks from Internet-scale video and uses them to infer object-centric transformations that serve as open-loop manipulation plans. Once anchors are obtained, policy learning conditions action prediction directly on them, through keypoint-conditioned behavior cloning (KALM~\cite{fang2024keypointabstractionusinglarge}), diffusion policies in object-relative frames (AffordDP~\cite{wu2025afforddpgeneralizablediffusionpolicy}), and shared human--robot 3D keypoint spaces (Point Policy~\cite{haldar2025pointpolicyunifyingobservations}). Across these variants, anchoring actions to spatial structure reduces perceptual ambiguity and improves robustness under distribution shift.

Because the central claim is out-of-distribution generalization, evaluation pairs high-fidelity simulation with real-robot deployment rather than relying on either alone. In simulation, affordance-centric methods such as AffordDP~\cite{wu2025afforddpgeneralizablediffusionpolicy} run on IsaacGym~\cite{makoviychuk2021isaac} over articulated objects from GAPartnet~\cite{geng2023gapartnet} whose functional parts (handles, knobs, levers) vary widely in geometry, while imitation suites such as Robomimic~\cite{robomimic2021} and RoboTwin~\cite{mu2024robotwin} isolate trajectory tracking and dual-arm coordination, with P3-PO~\cite{levy2024p3poprescriptivepointpriors} using them for controlled spatial-anchoring ablations. Real-world suites then stress in-the-wild generalization on articulated (drawers, microwaves), deformable (towel folding), and tool-use (sweeping) tasks, as in Point Policy~\cite{haldar2025pointpolicyunifyingobservations} and Track2Act~\cite{bharadhwaj2024track2act}, and a zero-shot cross-category protocol (train on one knife, test on spatulas or novel knives) is formalized by Robo-ABC~\cite{ju2024roboabcaffordancegeneralizationcategories} to measure how well retrieved anchors transfer without extra supervision. Performance is reported as Task Success Rate and, to isolate the contribution of anchoring, a Generalization Gap between seen and unseen instances. Spatially grounded policies outperform pixel baselines under shift in these studies: P3-PO reports large gains on cluttered, novel-object scenes~\cite{levy2024p3poprescriptivepointpriors}, and Robo-ABC improves both affordance-retrieval accuracy and real-robot success without heavy robot-data collection, binding demonstrations to reusable spatial structure more effectively than classic one-shot imitation~\cite{duan2017oneshotimitationlearning}.


\subsection{Pivotal Task-2: Inferring Functional Action Possibilities}
\label{subsec:p4-sec9-1-pivotal-task-2-inferring-functional-action-possibilities}

This pivot shifts from general object recognition to extracting \textbf{object affordances}, \ie, the specific action possibilities implied by an object's functional properties, and uses them as the intermediate representation that bridges raw perception and contact-rich execution. Whereas Pivotal Task-1 treats contact points as viewpoint-stable \emph{anchors} for conditioning a policy, affordances here are studied as \emph{functional possibilities} that specify what action an object supports. As Figure~\ref{fig:p4_s9_1_affordances} shows, the thread runs from how affordances are \emph{represented}, through the \emph{actions} they drive (grasping, placement, tool use), to \emph{outcome-oriented} benchmarks that score achieved function rather than imitated trajectory.

\begin{figure}
    \centering
    \includegraphics[width=0.98\linewidth]{figures/part4/section9/robotic_affordances.pdf}
    \caption{Affordances as functional action possibilities: a task- and context-conditioned affordance representation (left) supports three parallel task examples (6-DoF grasp, stable placement, hook-and-pull tool use), each judged by its functional outcome. Green checks mark intended outcomes in these illustrations, not measured results.}
    \label{fig:p4_s9_1_affordances}
\end{figure}

To capture functional properties beyond whole-object categorization, prior work builds object-centric representations (part-level segmentation~\cite{do2018affordancenet}, dense visual descriptors~\cite{florencemanuelli2018dense}, and keypoint-based encodings~\cite{manuelli2019kpam,qin2020keto,mo2021where2act}) that directly support action synthesis, and more recent vision--language models predict spatial affordance points directly from open-vocabulary prompts (RoboPoint~\cite{yuan2024robopoint}, MOKA~\cite{liu2024moka}). These representations drive learning-based 6-DoF grasping~\cite{sundermeyer2021contact,zhao2025high,jiang2021synergies,zhao2024efficient,yuan2023m2t2,deshpande2025graspmolmo} and stable object placement~\cite{zeng2020transporter,liu2022structformer,zhao2025towards,yuan2023m2t2,zhao2025anyplace} by grounding generation in functional structure rather than raw geometry, and a specialized strand further models tool--behavior relationships, either learning tool--action--effect pairs~\cite{stoytchev2008learning} or predicting specific affordances such as pulling~\cite{tikhanoff2013exploring}, so as to infer end-effector poses that achieve the intended functional outcome.

Benchmarks for affordance methods emphasize understanding object functionality over geometric matching: the property under test is \emph{functional} generalization, not trajectory imitation. For grasping, YCB-based real-world evaluations~\cite{Calli2015YCB}, GraspNet-1Billion~\cite{Fang2020GraspNet}, and simulation suites such as Dex-Net~\cite{Mahler2017DexNet} score 6-DoF stability and contact validity across novel instances and viewpoints; placement benchmarks such as RLBench~\cite{james2020rlbench} and ManiSkill~\cite{mu2021maniskill} test stable stacking and insertion in cluttered, articulated scenes, and M2T2~\cite{yuan2023m2t2} explicitly evaluates valid placement poses conditioned on object geometry and scene context. For articulated objects and tools, Where2Act~\cite{mo2021where2act} and PartNet-Mobility~\cite{xiang2020sapien} score the prediction of actionable parts and interaction parameters for pulling, pushing, and rotating, while tool-use tasks judge whether the inferred end-effector pose produces the intended causal effect (drawing, hooking, pulling) rather than reproducing a specific trajectory. Evaluation therefore centers on outcome-oriented success: stable grasps that lift and hold unseen objects under pose, scale, and appearance variation; placement success, pose accuracy, and post-release stability under clutter and occlusion; and, for articulated and tool interactions, whether the predicted parameters yield the desired functional result.



\subsection{Pivotal Task-3: Encoding Spatial Reasoning in Generalist Models}
\label{subsec:p4-sec9-1-pivotal-task-3-encoding-spatial-reasoning-in-generalist-models}

This pivot moves from specialized, single-task policies toward \textbf{generalist robotic foundation models}. Existing Vision--Language--Action (VLA) models such as RT-2~\cite{zitkovich2023rt2} and OpenVLA~\cite{kim2024openvla} show strong semantic understanding but often lack the spatial common sense needed for precise manipulation or robust generalization to novel viewpoints, so the central challenge is to strengthen spatial grounding during pre-training, enabling explicit reasoning about 3D geometry and generalization across environments and embodiments~\cite{chen2024spatialvlm}. The general foundation-model and spatial-reasoning machinery these policies inherit is treated in Chapter~\ref{sec:reasoning}; here we examine only its manipulation-specific adaptation, which Figure~\ref{fig:p4_s9_1_vla} groups into three families: explicit 3D representations, view-invariant representations, and grounded reasoning traces.

\begin{figure}
    \centering
    \includegraphics[width=0.98\linewidth]{figures/part4/section9/robotic_vla.pdf}
    \caption{Spatial grounding mechanisms for generalist robot policies: (1) explicit 3D features fused with task tokens (3D-LLM, 3D-VLA), (2) view canonicalization (OG-VLA) and spatial action discretization (SpatialVLA), two separate methods rather than stages of one, and (3) grounded reasoning traces that condition action prediction. The bottom row shows evaluation of viewpoint robustness and cross-robot transfer; all maps and overlays are schematic.}
    \label{fig:p4_s9_1_vla}
\end{figure}

The first family builds explicit 3D representations: models such as 3D-LLM~\cite{hong20233d} and 3D-VLA~\cite{zhen20243d} bridge point clouds and large language models, with 3D-VLA adding a 3D generative pre-training objective that lets the model reason about action outcomes in a consistent geometric space, effectively ``imagining'' future states. The second groups view canonicalization and spatial action encoding as distinct mechanisms: OG-VLA~\cite{singh2025og} projects multi-view observations into a point cloud and re-renders canonical orthographic views to reduce sensitivity to camera placement, while SpatialVLA~\cite{qu2025spatialvla} uses Ego3D position encoding for observations and \textbf{Adaptive Action Grids} to discretize spatial movement actions; adapting to new setups can involve re-discretizing the grids. The third injects spatial reasoning through grounded reasoning traces: ``thinking'' models such as Emma-X~\cite{sun2024emma}, MolmoAct series~\cite{lee2025molmoact, fang2026molmoact2}, closely related to TraceVLA~\cite{zheng2024tracevla}, replace direct pixel-to-action regression with intermediate, grounded outputs (2D motion trajectories or relative depth maps) that explicitly condition the final action prediction.

These approaches are evaluated under method-specific zero-shot and adaptation protocols across embodiments, viewpoints, and environments, with key metrics including multimodal goal-generation quality, robustness to extrinsic camera perturbations, and skill transfer across robot morphologies such as from a fixed-base Franka arm to a mobile manipulator~\cite{team2024octo}. The reported gains point the same way: canonical orthographic views keep grounding consistent under viewpoint change (OG-VLA), spatially aligned action grids share 3D structure across setups through discretized movement (SpatialVLA), and intermediate visual or depth-aware reasoning traces act as a regularizer on action prediction (Emma-X, MolmoAct).

\subsection{Pivotal Task-4: Satisfying Geometric and Logical Constraints}
\label{subsec:p4-sec9-1-pivotal-task-4-satisfying-geometric-and-logical-constraints}

This pivot moves planning from heuristic trajectory generation toward formal \textbf{constrained optimization} that enforces geometric feasibility, fundamental physical laws (non-penetration, friction cones, force limits), and high-level logical task structure. By casting manipulation as multi-objective optimization, these methods seek to satisfy modeled hard constraints while optimizing soft task objectives, spanning both high-frequency control and long-horizon decision-making. Figure~\ref{fig:p4_s9_1_constraints} organizes the thread along three fronts (symbolic--geometric planning, contact and trajectory optimization, and constraint-aware generative planning) before synthesizing the roles constraints play in learning-based systems.

\begin{figure}
    \centering
    \includegraphics[width=0.98\linewidth]{figures/part4/section9/robotic_constraints.pdf}
    \caption{Geometric and logical constraints in manipulation: (1) symbolic--geometric planning with revision after infeasibility, (2) contact and trajectory optimization (inset: normal force $f_n$ into the object, tangential force $f_t$ along its surface), and (3) constraint-aware generative planning. The bottom row shows three roles constraints play for learned policies.}
    \label{fig:p4_s9_1_constraints}
\end{figure}

Methodological progress has moved from purely symbolic planners toward neuro-symbolic, learning-augmented optimization frameworks that unify discrete logic with continuous physics. Classical Task and Motion Planning (TAMP) and Logic-Geometric Programming (LGP)~\cite{toussaint2015lgp} provide the foundation by jointly reasoning over symbolic task structure and geometric feasibility, and recent systems augment them with Large Language Models and Vision-Language Models for scalability in open-ended domains; for example, VIZ-COAST~\cite{yan2025vizcoast} uses the spatial reasoning of VLMs to predict refinement failures \emph{before} expensive geometric optimization, substantially reducing search complexity on long-horizon problems. At the level of physical interaction, contact-implicit optimization discovers contact schedules without explicit mode enumeration but is traditionally too slow for real-time control; GPU-parallelized systems such as SPaSM~\cite{chen2025spasm} compile constraint evaluation and gradient computation into CUDA kernels, accelerating joint optimization of object placements and robot trajectories. Beyond deterministic constraints, learning-based and generative models capture uncertainty and complex structure: differentiable physics engines~\cite{freeman2021brax} let gradients flow through physical interactions, stochastic implicit neural signed distance functions~\cite{quintero2024stochastic} model geometric uncertainty directly within the planning manifold, and diffusion-based planners such as JM2D~\cite{jung2025joint} couple model-free diffusion with model-based constraint optimization during sampling, rather than relying only on a sequential plan-then-filter interface.

Such systems are benchmarked on multi-stage re-grasping under uncertainty, forceful tool use, and tight-tolerance sequential interaction, since what they promise is reliability over long, contact-rich horizons; suites such as the Functional Manipulation Benchmark (FMB)~\cite{fmb2023} scoring not only task completion but constraint satisfaction under partial observability, external disturbance, and compounding execution error. In these systems constrained optimization plays three roles: it acts as a \textit{teacher}, generating physically valid trajectories that supervise imitation and offline reinforcement learning; as a \textit{safety shield}, where model-based projections map unconstrained neural or diffusion policy outputs onto constraint-satisfying sets at execution time~\cite{römer2025diffusionpredictivecontrolconstraints}; and as a \textit{structural prior} that prunes physically infeasible regions of the action space and curbs compounding error in long-horizon tasks. Together these roles keep learned and generative planners grounded in geometric and physical reality, reducing physically implausible outputs and supporting interpretable, feasible plans under the modeled assumptions; real-world execution still depends on state estimation, model accuracy, solver convergence, and disturbance handling.

Table~\ref{tab:p4_s9_1_benchmarks} consolidates the principal benchmarks, representative methods, and headline evaluation focus across the four pivots, situating the spatially grounded methods surveyed above within a common evaluation landscape.

\begin{table}[t]
\centering
\fontsize{8}{10}\selectfont
\setlength{\tabcolsep}{1.0mm}
\caption{Representative benchmarks, methods, and headline metrics for the four manipulation pivots, drawn from simulation and real-robot settings; metrics are reported on out-of-distribution object instances and configurations where the underlying work permits.}
\label{tab:p4_s9_1_benchmarks}
\begin{tabular}{L{0.15\textwidth} L{0.32\textwidth} L{0.25\textwidth} L{0.19\textwidth}}
\toprule
\textbf{Pivot} & \textbf{Representative Benchmark(s)} & \textbf{Representative Method(s)} & \textbf{Headline Metric} \\
\midrule
P1: Spatial anchors & IsaacGym~\cite{makoviychuk2021isaac} with GAPartNet~\cite{geng2023gapartnet}, Robomimic~\cite{robomimic2021}, RoboTwin~\cite{mu2024robotwin} + real-world & P3-PO~\cite{levy2024p3poprescriptivepointpriors}, Point Policy~\cite{haldar2025pointpolicyunifyingobservations} & Task SR / Generalization Gap \\
P2: Affordances & GraspNet-1B~\cite{Fang2020GraspNet}, Dex-Net~\cite{Mahler2017DexNet}, Where2Act~\cite{mo2021where2act}, PartNet-Mobility~\cite{xiang2020sapien} & M2T2~\cite{yuan2023m2t2} & Grasp / Placement SR \\
P3: Generalist VLA & ManiSkill~\cite{mu2021maniskill}, Octo eval~\cite{team2024octo} & SpatialVLA~\cite{qu2025spatialvla}, 3D-VLA~\cite{zhen20243d} & Zero-shot cross-embodiment SR \\
P4: Constraints & FMB~\cite{fmb2023} & SPaSM~\cite{chen2025spasm}, JM2D~\cite{jung2025joint} & Constraint satisfaction / long-horizon SR \\
\bottomrule
\end{tabular}
\end{table}

\subsection{Established Tools and Systems}
\label{subsec:p4-sec9-1-established-tools-and-systems}

Progress in spatially grounded manipulation has been inseparable from a maturing ecosystem of simulators, open software stacks, shared datasets, and, increasingly, commercial robot platforms that serve as common development and deployment substrates.

\textbf{Simulation and physics engines.} GPU-parallel engines such as NVIDIA IsaacGym and the newer Isaac Sim / Isaac Lab\footnote{\url{https://developer.nvidia.com/isaac/sim}}~\cite{makoviychuk2021isaac}, articulated-object simulators such as SAPIEN~\cite{xiang2020sapien} (with its PartNet-Mobility assets), and general-purpose physics engines including MuJoCo~\cite{todorov2012mujoco}, PyBullet\footnote{\url{https://pybullet.org}}, and the recent Genesis\footnote{\url{https://github.com/Genesis-Embodied-AI/Genesis}} provide high-throughput, contact-rich environments, while ManiSkill~\cite{mu2021maniskill}, RLBench~\cite{james2020rlbench}, Robomimic~\cite{robomimic2021}, and RoboTwin~\cite{mu2024robotwin} package standardized task suites, demonstrations, and evaluation protocols that have become de facto testbeds for imitation and policy learning.

\textbf{Open software stacks and datasets.} On the software side, the Robot Operating System (ROS\,2) with the MoveIt motion-planning framework\footnote{\url{https://moveit.ai}} remains the de facto middleware for real-robot integration, and Hugging Face's LeRobot\footnote{\url{https://github.com/huggingface/lerobot}} packages datasets, pretrained policies, and training pipelines for low-cost reproduction. Large cross-embodiment efforts such as Open X-Embodiment\footnote{\url{https://robotics-transformer-x.github.io}}, which aggregates demonstrations and the RT-X policies across dozens of robot platforms, together with open generalist checkpoints such as Octo~\cite{team2024octo} and OpenVLA~\cite{kim2024openvla} and NVIDIA's GR00T humanoid foundation models\footnote{\url{https://developer.nvidia.com/isaac/gr00t}}, expose pretrained weights and fine-tuning interfaces that lower the barrier to transferring foundation-model policies onto new morphologies.

\textbf{Commercial and humanoid platforms.} On the deployment side, a rapidly growing set of companies is moving spatially grounded manipulation out of the laboratory. Boston Dynamics\footnote{\url{https://bostondynamics.com}} (the electric Atlas and the Stretch warehouse robot), Tesla Optimus\footnote{\url{https://www.tesla.com/we-robot}}, Figure AI\footnote{\url{https://www.figure.ai}}, 1X Technologies\footnote{\url{https://www.1x.tech}}, Agility Robotics\footnote{\url{https://agilityrobotics.com}} (the Digit humanoid, piloted in logistics), Physical Intelligence\footnote{\url{https://www.physicalintelligence.company}} (the $\pi_0$ generalist policy), and Sanctuary AI\footnote{\url{https://sanctuary.ai}} each pair learned manipulation with humanoid or mobile-manipulator hardware, while firms such as Covariant\footnote{\url{https://covariant.ai}} target warehouse pick-and-place at commercial scale. These platforms are the clearest evidence that the spatial-intelligence capabilities surveyed above are crossing from benchmark tasks into economically deployed systems.

Together these tools convert the representation and reasoning machinery of Chapter~\ref{sec:reasoning} into reproducible, physically grounded manipulation pipelines; their principal limitation remains the fidelity gap between simulated and real contact dynamics, which the following discussion revisits.


\subsection{Challenges}
\label{subsec:p4-sec9-1-challenges}

Despite rapid progress across all four pivots, bringing \emph{spatial intelligence} to robust, general-purpose manipulation remains bounded by a set of intertwined difficulties that cut across perception, policy, and planning. They cluster into four threads: keeping a persistent and calibrated internal model of the scene; closing the loop between contact physics and active sensing; transferring skills across embodiments while sourcing the data to do so; and meeting the reliability budget of a deployed system.

\textbf{Persistent, calibrated world models.} The first requirement is a \emph{persistent 3D belief}: an actionable, temporally consistent world model that survives heavy occlusion, articulation, deformation, and sensor drift, and that is maintained at control rates rather than reconstructed from scratch each frame. Persistence alone is not enough if the model cannot say how much it trusts itself. \emph{Calibrated uncertainty} asks that distributional beliefs over geometry, contact state, friction, and articulation propagate through perception, policy, and planner, so that a decision reflects confidence rather than a bare point prediction and a low-confidence grasp can trigger caution instead of failure.

\textbf{Contact modeling, sensing, and active perception.} Manipulation is ultimately about contact, yet most spatial models still reason about collision-free geometry rather than the physics of touch. Progress on \emph{contact modeling and sensing} hinges on tightly integrating vision with force and tactile feedback to reason about stick--slip, compliance, and incipient failure. Because contact is often unobservable until it happens, this couples naturally to \emph{information-seeking control}: interleaving manipulation with active perception and probing (touch-to-disambiguate, look-while-acting) inside tractable belief-space planners. And because contact is precisely where simulators diverge from reality, \emph{sim-to-real for contact} stays open, demanding better material, friction, and compliance models together with fast system identification for on-robot adaptation.

\textbf{Cross-embodiment transfer and data at scale.} Skills learned on one platform rarely move cleanly to another. \emph{Cross-embodiment and cross-scene transfer} requires aligning action spaces and spatial tokens across grippers, cameras, and environments without hand-tuned frames or retraining. Sustaining that transfer is a data problem as much as an algorithmic one: \emph{data and evaluation at scale} calls for ``long-tail'' data engines that cover rare geometries, materials, and tools with counterfactual and safety labels, and for benchmarks that measure calibration, persistence, recovery, and long-horizon competence rather than success rate alone.

\textbf{Systems reliability.} Finally, a deployed manipulator must hold together as a \emph{system}: meeting latency budgets, absorbing calibration drift, detecting out-of-distribution geometry, and enabling safe fallback and continual learning without catastrophic forgetting. Closing these gaps will require tighter couplings among 3D world models, affordance predictors, language-grounded objectives, and constraint-aware planners, so that spatial representations are not only expressive but \emph{actionable, calibrated, and safe} in unstructured settings.

\newpage
\section{Autonomous Driving}
\label{sec:p4_s9_2}

% NuScenes-SpatialQA: A Spatial Understanding and Reasoning Benchmark for Vision-Language Models in Autonomous Driving
% Gpt-driver: Learning to drive with gpt


\subsection{Introduction}
\label{subsec:p4-sec9-2-introduction}
Autonomous driving sits at the intersection of spatial intelligence, transportation engineering, and human--machine interaction, and it is one of the most demanding real-world testbeds for the spatial capabilities surveyed in Part~\ref{part:3}.
A self-driving vehicle must continuously close a \emph{perceive--localize--decide--act} loop: it has to recognize and track the surrounding traffic scene, estimate its own pose within a metric map, reason about how nearby agents will move, and commit to safe control actions, all under tight latency and safety budgets.
Advances in deep learning, multi-sensor fusion, and vehicle-to-everything (V2X) communication have steadily pushed the boundaries of what vehicles can perceive, predict, and decide, while fundamental questions around robustness, trust, ethics, and regulation remain open.
Perception and reasoning themselves are treated in Chapter~\ref{chapter:3:sec:6} and Chapter~\ref{sec:reasoning}; what is specific to driving is that each capability must hold up at speed, in all weather, and with lives at stake.

Figure~\ref{fig:p4_s9_2_levels} summarizes the SAE J3016 automation levels that frame this domain.
The SAE J3016 standard \footnote{\url{https://www.sae.org/news/2019/01/sae-updates-j3016-automated-driving-graphic}} provides a widely adopted framework for defining automation levels, with most current production systems still operating at Level 2.
Progress toward higher levels, especially full Level 5 autonomy, faces major technical \cite{YurtseverLCT20}, ethical \cite{StahlTM16}, and financial \cite{fisher2020governmental} challenges.



Deep learning has driven most of this progress, from Transformers~\cite{WangLLLXLCLZH22} and graph networks~\cite{ZhangFWH22} that model multi-agent interaction to lifelong learning for open-world adaptation~\cite{ZhuBZZW24}, but its opacity raises safety concerns compounded by sensor noise and the limited reach of vehicle-to-vehicle communication, and reaching Levels 4 and 5 depends as much on redundant hardware, large-scale data pipelines, and verifiable safety assurance as on model architecture.


\begin{figure}[!ht]
    \centering
    \includegraphics[width=0.98\linewidth]{figures/part4/section9/driving_levels.pdf}
    \caption{Functional distinctions among SAE J3016 driving-automation levels. Levels 0--2 retain a human driver; Level 3 requires a human to resume driving when necessary, whereas Levels 4--5 do not rely on human fallback within their operating scope. The stair heights indicate ordered categories, not measured safety, model complexity, or technology readiness.}
    \label{fig:p4_s9_2_levels}
\end{figure}

\subsection{System Overview}
\label{subsec:p4-sec9-2-system-overview}
Modern self-driving stacks realize the \emph{perceive--localize--decide--act} loop as a modular pipeline, sketched in Figure~\ref{fig:p4_s9_2_overview}. 
A front-end \emph{perception} module converts raw, multi-modal sensor streams (camera, LiDAR, radar) into a structured model of the scene: what objects are present, where the drivable surface and lane geometry lie, and which agents are most safety-critical. 
A \emph{localization and mapping} module places the ego-vehicle within a metric, typically HD-map-anchored, representation of that scene, so that perception outputs can be expressed in a stable world frame. 
A \emph{decision and control} module then predicts how the scene will evolve and selects the steering, acceleration, and braking commands that are sent to the actuators and fed back into the environment, closing the loop. Each of the three stages is a distinct spatial-intelligence problem with its own representations, benchmarks, and failure modes, and they are the pivotal tasks below, in pipeline order.

\begin{figure}[!ht]
    \centering
    \includegraphics[width=0.98\linewidth]{figures/part4/section9/autonomous_driving_overview.pdf}
    \caption{Closed-loop spatial intelligence for autonomous driving: complementary sensor streams feed perception and localization and mapping, whose scene and ego/map state drive decision, planning, and trajectory tracking; vehicle--environment interaction returns new observations. The controller also receives ego state; the curved wire crossing is not a junction. Dashed execution logs feed evaluation separately from control. Scenes, maps, and evaluation plots are illustrative, not measured results.}
    \label{fig:p4_s9_2_overview}
\end{figure}

\subsection{Pivotal Task-1: Perception and Detection}
\label{subsec:p4-sec9-2-pivotal-task-1-perception-and-detection}
Perception is the front-end of the pipeline: it turns raw camera, LiDAR, and radar streams into a structured, layered model of the scene that every downstream module consumes, reusing the detection, segmentation, and 3D scene-understanding machinery of Chapter~\ref{chapter:3:sec:6} but specializing it to the driving Operational Design Domain (ODD). It is best read as a set of complementary structured representations, sketched in Figure~\ref{fig:p4_s9_2_perception}: object detection localizes traffic participants, semantic and occupancy prediction labels the surrounding space, lane detection recovers the road geometry that constrains motion, and pedestrian recognition isolates the most safety-critical agents.

\begin{figure}
    \centering
    \includegraphics[width=0.98\linewidth]{figures/part4/section9/driving_perception_tasks.pdf}
    \caption{Four complementary driving perception tasks feeding a shared scene model: (1) object detection, (2) semantics and 3D occupancy, (3) lane geometry, and (4) pedestrian/VRU tracking and motion prediction. Annotations and the dashed pedestrian path are illustrative, not measured predictions.}
    \label{fig:p4_s9_2_perception}
\end{figure}

The base layer is \textbf{object detection}, which identifies and localizes obstacles and is the most direct determinant of collision safety \cite{FengHWD22}. Its central practical constraint is the ODD: environmental factors such as poor weather degrade sensor visibility and model accuracy \cite{KoopmanF19}, so recent work formalizes and iteratively validates the ODD through quantifiable environmental models and scenario catalogues, trusting a detector only within conditions it has demonstrably handled, though the absence of a standardized ODD specification still makes cross-system safety claims hard to compare. Robustness to adverse weather is therefore a first-class concern \cite{zhang2023perception,TahirZACE24,HnewaR21}: rain, snow, fog, and glare corrupt LiDAR, radar, and camera signals alike, and because driving datasets under-represent these conditions, models trained on fair-weather data generalize poorly. Responses range from dedicated robust pipelines, such as multi-stage visibility enhancement coupled with deep CNN tracking \cite{HassaballahKMM21}, to new sensing hardware, \eg, frequency-modulated continuous-wave (FMCW) LiDAR and high-dynamic-range cameras that recover usable signal in glare at higher cost. Moving from the image plane into metric space, \textbf{3D object detection} localizes and tracks objects in the vehicle's surroundings \cite{LiWLXSLQD22}, where the bird's-eye-view (BEV) representation has become dominant because it yields a unified top-down map well suited to planning. BEVFormer \cite{YangCTTZZHLQLZD23} fuses multi-camera and temporal cues with spatiotemporal transformers to approach LiDAR-level accuracy on nuScenes, while Fast-BEV \cite{LiHCCLSLXSOS24} and SparseBEV \cite{LiuTLWW23} trade marginal accuracy for the lightweight, sparse, query-based computation needed for real-time on-vehicle deployment, together tracing the field's balancing of accuracy, latency, and deployability.

Detections tell the planner \emph{where objects are} but not \emph{which space is free}; that dense layer is supplied by \textbf{semantic segmentation} and its 3D extension, occupancy prediction \cite{ZhouBWN20,NestiRNBB22}. The progression runs from 2D pixel labeling with FCN~\cite{long2015fcn} and DeepLab~\cite{chen2017deeplabv2} to volumetric reasoning: the Tri-Perspective View \cite{HuangZZ0L23} lifts RGB into a 3D representation rivalling LiDAR, SurroundOcc \cite{WeiZZZZL23} uses multi-camera 2D--3D attention to capture arbitrary shapes, and occupancy world models such as OccWorld~\cite{zheng2024occworld} extend this to forecasting how the 3D scene evolves, moving beyond the limits of bounding boxes. Because these dense predictors are data-hungry and domain-sensitive, robustness techniques such as the SePiCo pixel-contrast framework \cite{XieLLLHW23} and GAN-based synthetic data \cite{AbdollahiPSMA21} are used to survive scene variability, occlusion, and weather shift. On top of the free-space layer, \textbf{lane detection} recovers the road geometry that defines the lateral reference for lane-keeping and planning; early CNN- and YOLO-style detectors trace markings well in clear conditions but are brittle on curved, worn, or occluded lanes, motivating successors that lift lanes into 3D and use attention over global context to infer occluded segments. 
Finally, \textbf{pedestrian recognition} focuses on vulnerable road users whose behavior can be difficult to predict: CNN-, YOLO-, and SSD-based detectors deliver the required precision and latency, color--thermal fusion handles low-light scenes, and spatial--temporal models couple detection with motion prediction to anticipate crossing behavior.

Across all four layers, perception is benchmarked on large-scale driving datasets, KITTI~\cite{KITTI-GeigerLU12}, the Waymo Open Dataset~\cite{WaymoOpen-SunKDCPTGZCCVHN20}, and nuScenes~\cite{caesar2020nuscenes} for 2D/3D detection, Cityscapes~\cite{cordts2016cityscapes} and BDD100K~\cite{yu2020bdd100k} for segmentation, dedicated lane benchmarks for lane detection, and CityPersons~\cite{zhang2017citypersons} for pedestrians, and scored with mean Average Precision (mAP) and the nuScenes Detection Score (NDS), mean Intersection-over-Union (mIoU), and F1/accuracy, respectively (Table~\ref{tab:p4_s9_2_benchmarks}). Increasingly these metrics are reported on adverse-weather and out-of-distribution splits, precisely to measure behavior within and beyond the declared ODD.

\subsection{Pivotal Task-2: Localization and Mapping}
\label{subsec:p4-sec9-2-pivotal-task-2-localization-and-mapping}
If perception answers \emph{what is around the vehicle}, localization and mapping answer \emph{where the vehicle is} and \emph{what the persistent structure of the world is}. The two are tightly coupled: a metric, HD-map-anchored pose lets perception outputs be expressed in a stable world frame that planning can rely on, and the map itself is built and maintained from the same sensor streams. As Figure~\ref{fig:p4_s9_2_localization} shows, the problem decomposes into pose estimation against a known map (localization), the construction of that map (3D mapping and reconstruction), and the joint solution of both when no map exists (SLAM), with multi-sensor fusion the common thread that holds accuracy together when any single modality degrades.

\begin{figure}
    \centering
    \includegraphics[width=0.98\linewidth]{figures/part4/section9/driving_localization_and_mapping.pdf}
    \caption{Localization and mapping through uncertainty-aware sensor fusion: camera, LiDAR, radar, and GNSS/IMU measurements update the state posterior, yielding a map-relative pose (evaluated by ATE/RPE), an incrementally updated 3D map, and loop closures that feed back only after verification. Pose and map are parallel products of joint estimation; only verified loop constraints return to the estimator. Recall@$N$ evaluates place retrieval, not pose accuracy.}
    \label{fig:p4_s9_2_localization}
\end{figure}

\textbf{Localization} has moved from GPS-dependent positioning to learned matching of live camera or LiDAR observations against a pre-built map, reaching centimeter-level accuracy even in GPS-denied urban canyons; semantic cues such as roads and buildings further let a vehicle re-localize in changed environments without full remapping. Because no single sensor is reliable across occlusion, weather, and GPS dropout, \textbf{sensor-fusion localization} is the dominant paradigm: probabilistic estimators (particle filters, Kalman filters, and NDT scan matching) reconcile GNSS, IMU, and LiDAR, while learned vision--radar fusion compensates for modality-specific weaknesses in rain, fog, and glare. The recurring bottlenecks here are cross-sensor spatial/temporal synchronization and fault tolerance, which motivate Bayesian uncertainty weighting and online-adaptive fusion that recalibrates as conditions change.

\textbf{3D mapping and reconstruction} supply the persistent representation that localization matches against, turning LiDAR and stereo streams into metric models; for the place-recognition and loop-closure that keep maps globally consistent, transformer descriptors such as SeqOT~\cite{ma2022seqot} match sequential LiDAR scans faster than the sensor frame rate, trained on richly annotated corpora such as ApolloScape~\cite{wang2019apolloscape}. The open problems are the cost of processing dense point clouds at city scale and the handling of occluded or incomplete returns, which favor incremental map updating over full re-mapping. When the map is unavailable, \textbf{SLAM} solves localization and mapping jointly, which is critical in GPS-denied corridors; modern driving-oriented vSLAM hybridizes classical geometry with learned features (for example, extending ORB-SLAM~\cite{mur2017orb} with learned features, semantic labels, and learned depth to favor stable landmarks) and recovers from drift through loop closure with bundle adjustment.

These capabilities are evaluated on KITTI Odometry~\cite{KITTI-GeigerLU12}, Oxford RobotCar~\cite{Maddern2017RobotCar}, and ApolloScape~\cite{wang2019apolloscape} for long-horizon SLAM and place recognition, on nuScenes~\cite{caesar2020nuscenes} and Waymo~\cite{WaymoOpen-SunKDCPTGZCCVHN20} for map-relative pose, and on Argoverse~\cite{Argoverse-ChangLSSBHW0LRH19} for the trajectory forecasting that localization feeds, \eg, with the Motion Transformer~\cite{shi2022mtr} (Table~\ref{tab:p4_s9_2_benchmarks}). The standard metrics, \ie, Absolute and Relative Pose Error (ATE/RPE), recall@$N$ for loop closure, and minADE/minFDE for forecasting, are reported under GPS-denied, adverse-weather, and sensor-degradation conditions to stress exactly the robustness the fusion machinery is built to provide.

\subsection{Pivotal Task-3: Decision Making and Control}
\label{subsec:p4-sec9-2-pivotal-task-3-decision-making-and-control}
Decision making and control is the action end of the pipeline: it consumes the scene model from perception and the pose from localization, predicts how the scene will evolve, and commits to the steering, acceleration, and braking that close the loop, drawing on the spatial-reasoning and planning paradigms of Chapter~\ref{sec:reasoning}. As Figure~\ref{fig:p4_s9_2_decision} indicates, two architectural philosophies dominate, differing in how much of the perceive-to-act mapping is learned end-to-end versus decomposed into interpretable modules, and they are increasingly converging.

\begin{figure}
    \centering
    \includegraphics[width=0.98\linewidth]{figures/part4/section9/driving_decision_and_control.pdf}
    \caption{Three alternative architectures for driving decision and control: (1) modular pipelines with an explicit BEV scene, (2) end-to-end sensor-to-control learning, and (3) hybrids with inspectable scene or trajectory intermediates and explicit constraint checks.}
    \label{fig:p4_s9_2_decision}
\end{figure}

The first philosophy is \textbf{learning-based decision making}, centered on (deep) reinforcement learning, which acquires driving policies through trial-and-error interaction and is well suited to the sequential, interactive nature of driving. 
Representative methods illustrate its design axes: the stochastic latent actor-critic (SLAC)~\cite{lee2020stochastic} learns compact latent states from high-dimensional images, PPO~\cite{schulman2017proximalpo} within a nonlinear control framework handles sharp turns and obstacle avoidance, and the DeepSet-Q architecture~\cite{huegle2019dynamic}, building on permutation-invariant set encoders~\cite{zaheer2017deep}, ingests a variable number of surrounding agents for high-level lane-change decisions.
Hybrids with explicit planning, such as Monte Carlo tree search guided by DRL or scene-representation transformers that forecast neighbor intentions, improve sample efficiency and situational awareness in dense traffic. 
The binding difficulty is that safety-critical episodes are rare in logged data, which pushes the field toward high-fidelity simulation for rare-event coverage and toward encoding hard safety constraints directly in the policy rather than trusting the reward signal alone. 
The second philosophy, \textbf{end-to-end decision making}, learns the mapping from sensors to trajectory or control directly, removing hand-designed interfaces between modules. Early instances learned state-to-action policies with DDPG~\cite{lillicrap2015continuous} in simulators such as TORCS, and multi-task backbones such as HybridNets~\cite{vu2022hybridnets} share computation across detection and segmentation without yet closing the loop to control. Current systems are trained by imitation on large driving logs: TransFuser~\cite{chitta2023transfuser} fuses camera and LiDAR features with transformer attention and predicts waypoints for a closed-loop agent, and UniAD~\cite{hu2023uniad} chains tracking, mapping, motion forecasting, and occupancy prediction through query interfaces inside one network optimized for planning. Driving world models such as GAIA-1~\cite{hu2023gaia} generate future multi-view scenes conditioned on ego actions, giving such policies a learned simulator for training and evaluation.
Opacity is the price of end-to-end training: attribution methods such as layer-wise relevance propagation (LRP)~\cite{bach2015pixel} support interpretation but do not by themselves establish safety or regulatory compliance, which is why hybrid designs keep inspectable scene or trajectory intermediates together with explicit constraint checks, whose guarantees extend only as far as the modeled constraints.

Closed-loop driving policies are evaluated chiefly in simulation, \ie, the CARLA urban simulator~\cite{dosovitskiy2017carla} (NoCrash, Longest6), the nuPlan planning benchmark~\cite{caesar2021nuplan}, and highway-interaction simulators built on HighD/INTERACTION (Table~\ref{tab:p4_s9_2_benchmarks}), because rare, safety-critical events cannot be exercised safely on public roads. The headline metrics are the CARLA Driving Score (route completion penalized by infractions), collision and intervention rates, and route completion, complemented by comfort proxies such as jerk and, for the open-loop forecasting sub-problem, minADE/minFDE.

\begin{table}[!th]
\centering
\fontsize{8}{10}\selectfont
\setlength{\tabcolsep}{1.0mm}
\caption{Representative benchmarks for spatial intelligence in autonomous driving, covering perception, localization and mapping, motion prediction, and decision making.}
\label{tab:p4_s9_2_benchmarks}
\begin{tabular}{L{0.22\textwidth} L{0.22\textwidth} L{0.28\textwidth} L{0.19\textwidth}}
\toprule
\textbf{Pivotal task} & \textbf{Benchmark} & \textbf{Representative method} & \textbf{Headline metric} \\
\midrule
Perception (3D detection) & nuScenes & BEVFormer~\cite{YangCTTZZHLQLZD23} & NDS / mAP \\
Perception (3D detection) & Waymo Open & SparseBEV~\cite{LiuTLWW23} & mAP / latency \\
Perception (segmentation) & Cityscapes & SePiCo~\cite{XieLLLHW23} & mIoU \\
Perception (occupancy) & nuScenes Occ3D & SurroundOcc~\cite{WeiZZZZL23} & mIoU \\
Localization \& mapping & KITTI Odometry & Deep / semantic vSLAM & ATE / RPE \\
Prediction (forecasting) & Argoverse & Motion Transformer (MTR)~\cite{shi2022mtr} & minADE / minFDE \\
Decision \& control & CARLA (Longest6) & TransFuser~\cite{chitta2023transfuser} & Driving Score \\
Decision \& control & nuScenes (open-loop planning) & UniAD~\cite{hu2023uniad} & L2 error / collision rate \\
Decision \& control & nuPlan & Learned planners & Closed-loop score \\
\bottomrule
\end{tabular}
\end{table}


\subsection{Established Tools and Systems}
\label{subsec:p4-sec9-2-established-tools-and-systems}
The systems on the road today populate a deployment ladder whose rungs are the SAE levels of Figure~\ref{fig:p4_s9_2_levels}, but what separates one rung from the next is rarely a single breakthrough in perception accuracy.
It is, instead, two coupled engineering choices: how wide an Operational Design Domain (ODD) the system claims, and where legal responsibility for the driving task sits.
Three tiers are in genuine production, \ie, supervised Level~2 assistance, certified Level~3 ``eyes-off'' automation, and driverless Level~4 robotaxis, and, as the following progression shows, moving up the ladder consistently means adding sensing and compute redundancy while \emph{narrowing} the domain in which the autonomy guarantee holds.

\subsubsection{Large-scale Production ADAS (SAE Level 2 / 2+)}
The defining property of Level~2 is that the system performs continuous lateral and longitudinal control, yet the human remains the legal driver and the immediate fallback; consequently, every Level~2 product is fundamentally an exercise in bounding the ODD and enforcing driver attention rather than in achieving autonomy, and the deployed systems differ mainly in how they strike that balance.
One design philosophy gates hands-free operation to pre-validated, HD-mapped limited-access highways and polices it with in-cabin driver monitoring: \textbf{GM Super Cruise}\footnote{\url{https://www.gm.com/innovation/autonomous-driving}}, \textbf{Ford/Lincoln BlueCruise}\footnote{\url{https://www.lincoln.com/technology/bluecruise/}}, and \textbf{BMW Driving Assistant Professional}\footnote{\url{https://www.bmwgroup.com/en/news/general/2023/driver-assistance-systems.html}} all keep the vehicle centered, regulate following distance, and offer assisted lane changes, but escalate alerts and disengage the moment attention drifts or the car leaves a supported segment; their reliability is a direct product of this deliberately narrow envelope.
A second philosophy trades map dependence for scale: \textbf{Tesla Autopilot}\footnote{\url{https://www.tesla.com/en_gb/support/autopilot}} and \textbf{FSD (Supervised)}\footnote{\url{https://www.tesla.com/en_gb/fsd}} rely on a camera-first vision stack refined through fleet-wide over-the-air training, which extends assistance to a far broader and less-bounded set of scenarios but, lacking comparable geofencing, still demands constant supervision and degrades on poor markings, construction zones, and adverse weather. At the conservative end, \textbf{Toyota Safety Sense} declines hands-free operation altogether, functioning as an always-on crash-avoidance layer (automatic emergency braking, lane-keeping, and adaptive cruise) that intervenes only to prevent or mitigate collisions. Despite these differing emphases, the binding limitation is identical across all of them: a bounded ODD coupled with a mandatory human fallback, and it is precisely that fallback which the next tier begins to remove.

\subsubsection{Certified ``Eyes-off'' Conditional Automation (SAE Level 3)}
Level~3 is the first tier at which, inside a defined envelope, the driver may legally divert attention and the manufacturer assumes liability for the driving task. This shift in responsibility, rather than any new perception capability, is what dictates the engineering.
\textbf{Mercedes-Benz Drive Pilot}\footnote{\url{https://www.mercedes-benz.co.uk/passengercars/technology/autonomous-driving.html}} and \textbf{Honda Traffic Jam Pilot}\footnote{\url{https://global.honda/en/tech/Traffic_Jam_Pilot_TJP/}} purchase the ``eyes-off'' guarantee by pairing complementary sensing (cameras, radar, and LiDAR with high-precision, HD-map-anchored localization) with redundant braking, steering, and power paths that retain control through a single-component failure.
Both, tellingly, keep the envelope deliberately small (low-speed or traffic-jam flow on mapped, divided highways), and both center on the same safety mechanism: as conditions approach the ODD boundary, the system issues a timed takeover request and, absent a driver response, executes a minimal-risk maneuver such as a controlled stop with hazard signals.
The narrowness of the Level~3 envelope is therefore not a marketing concession but the direct cost of guaranteeing safe behavior once the human is no longer continuously in the loop: a cost the next tier pays in full by removing the driver entirely.

\subsubsection{Driverless Robotaxis and Autonomous Ride-hailing (SAE Level 4)}
Level~4 removes the human from the loop entirely within a geofence: there is no driver to hand back to, so the vehicle must resolve every in-domain situation itself or bring itself to a safe state, with remote operators supplying only high-level guidance rather than moment-to-moment control.
To meet this far stricter bar, \textbf{Waymo One}\footnote{\url{https://waymo.com/}} and \textbf{Baidu Apollo Go}\footnote{\url{https://www.apollogo.com/}} converge on essentially the same recipe: a dense, overlapping multi-sensor suite (lidar, radar, and camera vision) for redundant coverage, HD-map-anchored localization, and conservative prediction-and-planning that favors lawful, cautious maneuvers through dense urban traffic.
Each makes the problem tractable not by claiming universal competence but by tightly restricting where and when it operates and validating performance city by city before extending the geofence.
Across the three tiers, every step up the ladder buys wider autonomy by narrowing the operating domain and adding redundancy, until at Level~4 the system finally internalizes the fallback that Levels~2 and~3 delegated to the driver.



\subsection{Challenges}
\label{subsec:p4-sec9-2-challenges}
Every tier of the deployment ladder works because it limits where it operates, and general, trustworthy autonomy requires lifting those limits. Four threads, each cutting across perception, localization, and decision, stand in the way: justifying life-critical decisions (\emph{safety and explainability}), dependable behavior across the long tail of road, weather, and traffic conditions (\emph{resilience and robustness}), adaptation to genuinely novel situations from little prior data (\emph{unfamiliar road scenarios}), and continual learning on deployed vehicles without eroding mastered behavior (\emph{enhanced autonomy in complex scenarios}).

\paragraph{Safety and Explainability.}
Because driving decisions are life-critical, the inability to explain why a model selected a given maneuver is a first-order obstacle to certification and public trust.
Post-hoc attribution methods such as layer-wise relevance propagation~\cite{bach2015pixel} and Shapley-value attribution expose which inputs drove a decision, supporting both engineering debugging and regulatory audit, while causal-inference models go further by separating mere correlation from the actual influence of factors such as road state and a maneuver's effect on surrounding agents.
The unresolved difficulty is that these explanations remain largely retrospective and local to individual decisions; converting them into guarantees that provably hold across the full operating domain is still an open problem.

\paragraph{Resilience and Robustness.}
Robustness to the long tail of driving conditions is commonly pursued through curriculum learning, where a policy is trained on a graded sequence from simple navigation to dense mixed-traffic interaction, which improves sample efficiency and training stability relative to confronting hard cases directly.
Its effectiveness, however, hinges on a scenario curriculum that faithfully spans the deployment distribution and on objectives that keep safety and comfort constraints satisfied at every difficulty level.
The practical risk is that, absent a benchmark which deliberately surfaces rare but critical events, strong curriculum performance on common scenes can mask failures precisely on the cases that determine real-world safety.

\paragraph{Unfamiliar Road Scenarios.}
A distinct difficulty is generalization to genuinely novel situations (an unseen sign, an irregular work zone, an atypical vehicle) where data-hungry models degrade but human drivers adapt from a single exposure.
Approaches that target this gap include probabilistic program induction in a Bayesian framework, which represents a new concept as a compact program inferred from few examples, and the injection of compositional, causal, and meta-learning priors so that, for instance, a system can infer that a fallen obstacle will trigger sudden lane changes in adjacent traffic.
The standing obstacle is validating such few-shot behavior to a safety standard, since the very rarity of these events makes their failure modes hard to measure and certify.

\paragraph{Enhanced Autonomy in Complex Scenarios.}
Sustained operation in changing environments motivates lifelong learning~\cite{ZhuBZZW24}, in which a deployed system keeps acquiring new knowledge (a new traffic pattern or road layout) without retraining from scratch and without overwriting prior competence.
The central technical barrier is catastrophic forgetting, which is addressed by memory-augmented architectures and continual-learning regularizers that protect previously learned parameters, together with transfer mechanisms that reuse knowledge across tasks and cities.
The harder, still-open requirement is to perform such updates on-vehicle under real-time and verification constraints, so that an update improving one scenario cannot silently regress a safety-critical behavior elsewhere.

\newpage

\section{Wearable AI System}
\label{sec:p4_s9_3}


\subsection{Introduction}
\label{subsec:p4-sec9-3-introduction}

Wearable devices have become a central sensing modality for embodied and spatial intelligence, delivering continuous, fine-grained measurements of human motion, physiology, and physical interaction. Because they are worn on or integrated into the body, they access proprioceptive and kinesthetic cues, joint rotations, muscle activations, contact forces, and locomotion patterns, which exteroceptive sensors such as RGB-D cameras or LiDAR cannot capture directly. We use the term broadly, covering inertial measurement units (IMUs), surface electromyography (EMG) sensors, smart textiles, haptic interfaces, and wearable cameras: body-mounted technologies that together model human dynamics and embodied interaction in real-world settings. It is this egocentric, body-centric vantage point that makes wearables valuable for spatial intelligence, letting a system reason about how a person moves through, reads, and acts upon their surroundings, and it has driven rapid adoption across health monitoring, sports analytics, rehabilitation~\cite{patel2012review}, industrial safety, and human--robot collaboration.

The signals wearables provide are especially valuable because they are anticipatory. Joint kinematics and muscle activity reveal what a person is doing and, crucially, what they are about to do next, which is what underpins intent recognition, action forecasting, and safe coordination in shared spaces. Wearable IMU arrays, for example, recover whole-body skeletal motion in cluttered or occluded settings where optical capture breaks down~\cite{huang2018dip, vonmarcard2017sparse}, and EMG surfaces pre-activation patterns that precede any visible movement~\cite{atzori2014ninapro}.

These devices are often bracketed with head-mounted XR hardware, yet their role in spatial intelligence differs in kind. XR headsets integrate sensing and display for immersive visualization and mixed-reality interaction, leaning on outward-facing cameras for SLAM and environment mapping. The wearables treated here look inward instead, supplying the biomechanical and physiological state, high-frequency, continuous, and on-body, that an XR headset does not observe.

The spatial-intelligence capabilities that wearables ultimately serve, pose estimation, action recognition, localization, and trajectory prediction, are formulated in Part~III (spatial perception in Chapter~\ref{chapter:3:sec:6}, spatial reasoning in Chapter~\ref{sec:reasoning}). This section does not re-derive them; it concentrates on the body-centric sensing layer, the devices that acquire embodied signals, the systems assembled around them, and the tasks, methods, and failure modes that surface once spatial intelligence must be grounded in continuous on-body measurement, spanning long-horizon activity and gait analysis, intention prediction for human--robot collaboration, and real-time safety in industrial and logistics settings.


\subsection{Wearable Devices}
\label{subsec:p4-sec9-3-wearable-devices}

Body-centric sensing spans a diverse device ecosystem, summarized in Table~\ref{tab:p4_s9_3_devices}.

\paragraph{Motion and Body-State Sensing.}
Inertial measurement units (accelerometers, gyroscopes, and often magnetometers) are the most widely deployed body-worn sensors. Consumer watches and bands such as the Apple Watch\footnote{\url{https://www.apple.com/watch/}}, Samsung Galaxy Watch\footnote{\url{https://www.samsung.com/global/galaxy/galaxy-watch/}}, Fitbit\footnote{\url{https://www.fitbit.com/}}, and Xiaomi Band stream inertial data for activity and gait monitoring, while research and industrial systems place IMUs at specific anatomical sites and fuse them into full-body kinematics~\cite{vonmarcard2017sparse, yi2021transpose}; because they need no line of sight, they keep working under occlusion and poor lighting. Full-body motion-capture suits such as Xsens\footnote{\url{https://www.movella.com/products/xsens}}, Rokoko Smartsuit Pro\footnote{\url{https://www.rokoko.com/}}, and Noitom\footnote{\url{https://www.noitom.com/}} combine many IMUs, sometimes with pressure or stretch sensors, into accurate skeletal pose outside the studio, which has made them standard in biomechanics, sports, robotics, and character animation~\cite{trumble2017totalcapture}. Smart clothing embeds conductive fabrics, stretch sensors, or distributed pressure sensors in garments for unobtrusive, long-duration monitoring of posture, joint angles, and whole-body contact~\cite{luo2021learning}, and physiological sensors (heart rate, galvanic skin response, respiration), foot-mounted inertial or pressure sensors, and ear-worn IMUs add workload, ground-contact timing, and head orientation~\cite{patel2012review}.

\paragraph{Neuromuscular, Egocentric, and Interactive Wearables.}
Electromyography (EMG) armbands and high-density arrays sense the muscle activations that precede visible movement, giving early cues for gesture recognition and intent estimation~\cite{atzori2014ninapro}, and combining EMG with inertial data makes them more robust to arm posture and fatigue~\cite{atzori2016deep}. Wearable cameras add the exteroceptive, first-person view: egocentric RGB, depth, or fisheye streams supply scene geometry and activity context that can be fused with on-body signals~\cite{grauman2022ego4d, adt}. Two device families also act on the body. Haptic gloves such as HaptX\footnote{\url{https://haptx.com/}} and SenseGlove\footnote{\url{https://www.senseglove.com/}} and haptic vests such as bHaptics\footnote{\url{https://www.bhaptics.com/}} return tactile, vibration, or force feedback during manipulation, teleoperation, and training~\cite{pacchierotti2017wearable}, and exoskeletons combine IMUs, torque and force sensors, and joint encoders with actuation, estimating user intent in real time to assist or rehabilitate movement~\cite{young2017state}, which makes them a direct test of embodied spatial intelligence.

\begin{table}[!th]
\centering
\caption{Wearable body-centric sensing devices, by sensed signal and recovered spatial or biomechanical quantity. Haptic and exoskeletal wearables are partly \emph{actuators} that return force and tactile signals to the wearer.}
\label{tab:p4_s9_3_devices}
\fontsize{8}{10}\selectfont
\setlength{\tabcolsep}{1.0mm}
\begin{tabular}{L{0.22\textwidth} L{0.28\textwidth} L{0.42\textwidth}}
\toprule
\textbf{Device family} & \textbf{Primary signal} & \textbf{Spatial / biomechanical quantity recovered} \\
\midrule
IMU wearables            & accel., gyro., (magnet.)          & orientation, acceleration, gait, sparse-to-full body kinematics \\
Motion-capture suits     & multi-IMU (+ pressure/stretch)    & full-body skeletal pose in unconstrained settings \\
EMG / gesture bands      & neuromuscular activation          & muscle intent, hand/arm gestures, pre-movement cues \\
Smart clothing / textile & distributed pressure, strain      & posture, joint angles, whole-body contact patterns \\
Egocentric cameras       & first-person RGB / depth          & scene geometry, activity, environmental context \\
Haptic / force-feedback  & tactile, vibration, force (I/O)   & contact rendering, spatial alerts (feedback channel) \\
Exoskeletons             & IMU, torque, force, joint enc.\   & joint torque, user effort and movement intent \\
Physiological / other    & HR, GSR, respiration; foot press. & workload, exertion, ground contact, head orientation \\
\bottomrule
\end{tabular}
\end{table}


\subsection{System Overview}
\label{subsec:p4-sec9-3-system-overview}

Taken together, the devices above form the front end of a single body-centric perception-to-action loop that defines the wearable AI system as a whole. 
Raw on-body signals are first \emph{sensed} (inertial, neuromuscular, tactile, and egocentric-visual streams), then \emph{modeled} into geometric and biomechanical state (joint kinematics, body pose, contact forces), \emph{interpreted and predicted} into higher-level behavioral and intentional descriptions (activities, gestures, short-horizon motion forecasts), and finally fed back to the wearer or to a collaborating agent through \emph{action} channels (haptic cues, exoskeleton assistance, safety alerts, or shared-control commands). 
The loop then closes: the consequences of feedback alter the wearer's motion, which the sensors observe anew.

Figure~\ref{fig:p4_s9_3_overview} schematizes this loop and situates the three pivotal tasks that organize the remainder of this section as complementary stages of it. 
\textbf{Pivotal Task-1 (Embodied Motion Sensing and Pose Reconstruction)} occupies the sense-to-model stage, turning body-centric signals into stable geometric descriptions of motion. 
\textbf{Pivotal Task-2 (Behavioral Understanding and Intention Modeling)} occupies the interpret-and-predict stage, mapping reconstructed motion and neuromuscular cues to recognized activities and anticipated actions. 
\textbf{Pivotal Task-3 (Safe Collaboration and Embodied Interaction)} closes the loop at the action stage, converting interpreted human state into coordinated, safe, and assistive responses in shared physical environments. 
Because all three stages must run under the power, compute, and reliability constraints of body-worn hardware, lightweight on-device inference is a cross-cutting concern that recurs throughout.


\begin{figure*}[!t]
    \centering
    \includegraphics[width=\linewidth]{figures/part4/section9/wearable_overview.pdf}
    \caption{The wearable AI system as a body-centric loop: on-body signals are modeled into kinematic and biomechanical state (T1), interpreted into activity and intent (T2), and acted on through haptic, exoskeleton, and safety channels (T3), used separately or together. Device and prediction illustrations are schematic.}
    \label{fig:p4_s9_3_overview}
\end{figure*}


\subsection{Pivotal Task-1: Embodied Motion Sensing and Pose Reconstruction}
\label{subsec:p4-sec9-3-pivotal-task-1-embodied-motion-sensing-and-pose-reconstruction}

The first stage of the wearable loop turns raw body-centric signals into a geometric description of how the wearer moves. 
It is a bottom-up reconstruction problem: per-segment inertial and biomechanical measurements must yield joint angles, segment velocities, and ultimately a full-body pose that remains stable over long horizons and in settings where exteroceptive capture fails. 
This task subsumes full-body skeletal pose estimation, fine-grained limb tracking, real-world locomotion and gait analysis, and body-centric localization in GPS-denied spaces; the underlying pose-estimation formulations are developed in Chapter~\ref{chapter:3:sec:6}, so the emphasis here is on reconstruction from on-body sensing. 
Figure~\ref{fig:p4_s9_3_motion} illustrates this reconstruction stack from segment-level signals to full-body pose.

\begin{figure}[!t]
    \centering
    \includegraphics[width=\linewidth]{figures/part4/section9/wearable_motion.pdf}
    \caption{Motion sensing and pose reconstruction from body-worn IMUs: estimators supported by a parametric body model recover root-relative full-body pose, while visual--inertial fusion anchors a world-frame trajectory. Contact or zero-velocity updates apply only when their assumptions hold, and metric scale requires sufficient observability; signals and trajectories are illustrative.}
    \label{fig:p4_s9_3_motion}
\end{figure}

\textbf{From signals to kinematics.} Reconstruction begins with kinematic and dynamic modeling, which provides the physical grounding for interpreting inertial and pressure-based signals. 
Forward and inverse kinematics infer joint angles and segment orientations from inertial measurements, compensating for drift and sensor noise through optimization-based filtering and biomechanical constraints. Classical recursive estimators such as extended and unscented Kalman filters incorporate kinematic priors and temporal smoothness to improve stability in real-world settings~\cite{welch1995kalman, madgwick2011estimation}, while articulated-body dynamics and parametric body models such as SMPL~\cite{loper2015smpl} enforce anatomical plausibility for complex motions such as rapid locomotion or limb coordination. 
Multi-IMU fusion against musculoskeletal models extends these ideas to full-body pose estimation in challenging environments where optical motion-capture systems fail~\cite{huang2018dip, yi2021transpose, yi2022pip}, and recent systems fold in physics-based dynamics and inter-sensor ultra-wideband ranging to tighten accuracy and suppress jitter~\cite{armani2024ultra}.

\textbf{Learning-based reconstruction and sensor fusion.} Where analytic models struggle with sparse or noisy sensing, data-driven approaches learn the mapping from raw inertial streams to pose directly~\cite{huang2018dip, yi2021transpose}, and sensor fusion integrates heterogeneous modalities for drift-compensated estimates. IMU-and-vision fusion resolves ambiguities in inertial tracking by leveraging environmental context from egocentric cameras~\cite{grauman2022ego4d, adt}, while drift-compensation techniques, including zero-velocity updates and constraint-based corrections, stabilize long-horizon wear. 
For body-centric localization, sparse-sensor systems couple inertial pose estimation with simultaneous localization and mapping, so that the wearer's global trajectory is recovered together with body pose rather than left to drift freely~\cite{yi2023egolocate}; visual-inertial odometry that fuses the same inertial signals with first-person imagery offers a complementary route through GPS-denied indoor spaces, assessed on dedicated benchmarks~\cite{Schubert2018TUMVI}.

\textbf{Benchmarks and evaluation.} The kinematic, learning-based, and fusion methods above are measured against the same yardsticks: large-scale motion-capture corpora such as AMASS~\cite{mahmood2019amass}, which provides a canonical surface-and-skeleton reference for full-body pose, alongside dedicated inertial-pose datasets for sparse-IMU reconstruction~\cite{trumble2017totalcapture, huang2018dip}, while body-centric localization is assessed separately on visual-inertial odometry benchmarks~\cite{Schubert2018TUMVI}. Reconstruction quality is reported as mean per-joint position error (MPJPE) and joint-angle error, and localization as absolute trajectory error and accumulated drift (Table~\ref{tab:p4_s9_3_benchmarks}). 
Read against the needs of real deployment, these numbers reveal a clear gap. On curated benchmarks, sparse-IMU systems now recover \emph{root-relative} full-body pose to within roughly five centimeters of optical ground truth~\cite{yi2022pip}, but global translation still drifts over long sequences unless it is anchored by SLAM or ranging~\cite{yi2023egolocate}, and even the root-relative accuracy falls off in the wild, where fast or unseen motions and loose sensor placement push errors past the tolerance that downstream control or clinical gait analysis can absorb. The pose the next stage consumes is thus dependable on average yet brittle at exactly the fast, atypical movements that matter most.


\subsection{Pivotal Task-2: Behavioral Understanding and Intention Modeling}
\label{subsec:p4-sec9-3-pivotal-task-2-behavioral-understanding-and-intention-modeling}

Once motion has been reconstructed, the second stage of the loop interprets it: recognizing what the wearer is doing and anticipating what they will do next. 
Continuous neuromuscular and inertial cues feed activity recognition and segmentation and then short-horizon intention prediction, even in daily-life settings where visual systems face severe occlusion or field-of-view limits. 
Because the recognition and reasoning formulations themselves are treated in Chapter~\ref{sec:reasoning}, the focus here is on how high-frequency wearable streams are modeled to support behavioral inference. 
Figure~\ref{fig:p4_s9_3_behavior} sketches the path from multimodal time-series and body-graph representations to recognized activity and predicted intent.

\begin{figure}[!t]
    \centering
    \includegraphics[width=\linewidth]{figures/part4/section9/wearable_behavior.pdf}
    \caption{Behavioral understanding and intention modeling: IMU and EMG sequences and a kinematic body graph pass through temporal and graph encoders, fused by attention to recognize current activity and forecast short-horizon intent. The EMG pre-activation inset draws anticipatory cues from the observed stream, not from future measurements.}
    \label{fig:p4_s9_3_behavior}
\end{figure}

\textbf{Sequence models for wearable time-series.} The availability of large-scale wearable recordings has driven deep models tailored to high-frequency temporal streams from IMU, EMG, and multimodal sensors. Recurrent networks, temporal convolutional networks, and Transformers are widely applied to gait classification, gesture recognition, and activity segmentation~\cite{ordonez2016deep}: temporal convolutions capture multi-scale dependencies for long-horizon pattern detection~\cite{bai2018tcn}, while attention-based architectures learn complex correlations across sensor channels~\cite{attention}. Because labeled wearable data are scarce and expensive, self-supervised pretraining on unlabeled motion recordings has become the dominant route to transferable representations~\cite{haresamudram2021contrastive}: masked-reconstruction encoders such as LIMU-BERT learn compact IMU features from raw streams~\cite{xu2021limubert}, and multimodal contrastive objectives such as IMU2CLIP align inertial signals with video and language, pointing toward wearable foundation models that generalize across tasks and users~\cite{moon2023imu2clip}.

\textbf{Graph-based body representations.} Because human anatomy is structured as a kinematic tree, graph neural networks offer a natural representation for wearable-derived motion~\cite{kipf2017semisupervised}, encoding relational dependencies between joints and limbs directly in graph space. 
Such formulations fuse IMU placements across body segments, propagate temporal features along skeletal edges, and remain robust to missing or corrupted sensor channels; combining graph convolution with temporal attention jointly models spatial structure and temporal evolution for activity recognition and biomechanical analysis~\cite{yan2018stgcn}.

\textbf{Neuromuscular cues and multimodal fusion for intent.} Beyond recognizing observed motion, EMG-based wearables detect subtle neuromuscular activations that precede visible movement~\cite{atzori2014ninapro}, providing predictive cues for early intention estimation and responsive human-robot collaboration.
Fusing EMG with inertial data sharpens fine-grained gesture and anticipatory-command decoding~\cite{atzori2016deep}, while fusing IMU with egocentric video grounds activity recognition in environmental context~\cite{stein2013combining, grauman2022ego4d}.
A recent milestone shows how far this can scale: a surface-EMG wristband trained across more than a thousand users decodes gestures, clicks, and handwriting from neuromuscular signals alone and, critically, generalizes to new wearers without per-user calibration~\cite{kaifosh2025neuromotor}, addressing the inter-user variability that has long limited EMG interfaces.
Physiological and pressure cues additionally support inference of fatigue, balance, and abnormal gait over extended periods.

\textbf{Benchmarks and evaluation.} Taking the reconstructed motion of the previous stage as input, the sequence, graph, and fusion models above are evaluated on wearable human-activity-recognition corpora~\cite{reiss2012pamap2, chavarriaga2013opportunity}, EMG gesture datasets~\cite{atzori2014ninapro}, and egocentric activity datasets such as Ego4D~\cite{grauman2022ego4d} and EPIC-KITCHENS~\cite{EPIC-KITCHENS}, with recognition reported as accuracy and macro-F1 and intention prediction as anticipation accuracy and prediction lead time (Table~\ref{tab:p4_s9_3_benchmarks}). 
Here the gap to maturity is one of generalization rather than peak score. Balanced locomotion corpora such as PAMAP2 are close to saturated, yet the same architectures score markedly lower on cluttered, null-dominated sets such as Opportunity and lose accuracy again under cross-user and cross-environment evaluation; reliable intention forecasting, meanwhile, still extends only to short, well-rehearsed horizons. This falls well short of the open-world, anticipatory understanding that the safe-collaboration stage below depends on, where the system must read unfamiliar users in unfamiliar settings.


\subsection{Pivotal Task-3: Safe Collaboration and Embodied Interaction}
\label{subsec:p4-sec9-3-pivotal-task-3-safe-collaboration-and-embodied-interaction}

The third stage closes the loop, converting interpreted human state into coordinated, safe, and assistive responses in shared physical environments. 
Contact- and force-level sensing of human--environment interaction supports real-time inference of effort and hazard, which in turn drives feedback for collaboration and shared autonomy. 
The collaboration and shared-control reasoning that this task builds on is developed in Chapter~\ref{sec:reasoning} and, for the manipulation setting specifically, in Section~\ref{sec:p4_s9_1}; here the emphasis is on the wearable feedback and safety layer. 
Figure~\ref{fig:p4_s9_3_safety} depicts this closing stage, from contact sensing to coordinated safety response.

\begin{figure}[!t]
    \centering
    \includegraphics[width=\linewidth]{figures/part4/section9/wearable_safety.pdf}
    \caption{Safe human--robot collaboration with wearables: fused contact, force, and proximity signals, together with intent forecasts (T2) and robot motion context, drive motion prediction and risk evaluation, which select glove, exoskeleton, or cobot responses. Prediction envelopes and risk thresholds are schematic, not measured safety guarantees.}
    \label{fig:p4_s9_3_safety}
\end{figure}

\textbf{Sensing interaction and effort.} Haptic and tactile wearables capture contact forces, surface interactions, and manipulative behaviors, supplying the signals needed to model human--environment interaction and object handling~\cite{HOI4D, luo2021learning}. Exoskeletons equipped with torque and force sensors use these signals to infer user effort and detect hazardous conditions, adjusting assistive strategies accordingly~\cite{young2017state}; probabilistic fusion that integrates biomechanical priors and uncertainty estimates is central to robustness in such safety-critical settings.

\textbf{Coordination, safety, and feedback.} In industrial and shared-autonomy scenarios, on-body inertial and proximity sensors contribute to safety monitoring by detecting collision risks, unsafe postures, and sudden motion anomalies, while haptic channels return spatial alerts and force feedback that close the sensorimotor loop~\cite{pacchierotti2017wearable}. 
Combined with the intention forecasts of Pivotal Task-2, these signals form the basis for predictive safety mechanisms and shared-control frameworks that anticipate human motion rather than merely reacting to it.

\textbf{On-device constraints.} Because feedback must be timely, the action stage is where the wearable system's power, memory, and compute limits bite hardest. Model compression (pruning, quantization, and distillation)~\cite{hinton2015distilling} and edge-optimized architectures~\cite{howard2017mobilenets} enable real-time inference on wearable microcontrollers or smartphone-class processors, and on-device adaptation mitigates inter-user variability without relying on remote servers, keeping safety-critical computation local and low-latency.

\textbf{Benchmarks and evaluation.} Closing the loop, the interaction-sensing, safety, and on-device methods above are assessed on human-object-interaction datasets such as HOI4D~\cite{HOI4D} and, for assistive control, on exoskeleton and teleoperation testbeds, with metrics spanning contact- and force-estimation error, collision-avoidance and safety-violation rates, and intention-anticipation lead time (Table~\ref{tab:p4_s9_3_benchmarks}). 
Evaluation here is the least mature of the three stages: it remains fragmented across hardware-specific testbeds with few shared protocols, so reported gains are hard to compare and seldom stress real safety guarantees. 
Current systems perform convincingly in instrumented demonstrations, but closed-loop, certified safety under the full variability of human partners, the bar that genuine shared autonomy must clear, is still an open target rather than a measured result. The robustness, generalization, and safety gaps that recur at each stage of the loop are what the challenges below take up directly.


\begin{table}[!th]
\centering
\caption{Representative benchmarks for wearable spatial-intelligence tasks, grouped by pivotal task.}
\label{tab:p4_s9_3_benchmarks}
\fontsize{8}{10}\selectfont
\setlength{\tabcolsep}{1.0mm}
\begin{tabular}{L{0.16\textwidth} L{0.20\textwidth} L{0.34\textwidth} L{0.21\textwidth}}
\toprule
\textbf{Pivotal Task} & \textbf{Benchmark / Dataset} & \textbf{Representative method / use} & \textbf{Headline metric} \\
\midrule
T1: Motion \& Pose & AMASS~\cite{mahmood2019amass} & Full-body pose reference (SMPL surface) & MPJPE \\
T1: Motion \& Pose & TotalCapture~\cite{trumble2017totalcapture} & Video+inertial full-body pose & MPJPE \\
T1: Motion \& Pose & DIP-IMU~\cite{huang2018dip} & Sparse-IMU pose reconstruction (DIP, TransPose, PIP) & Joint-angle / position error \\
T1: Localization & TUM-VI~\cite{Schubert2018TUMVI} & Visual-inertial odometry (body-centric) & ATE / drift \\
T2: Activity & PAMAP2~\cite{reiss2012pamap2} & Wearable activity recognition (deep recurrent / TCN) & Accuracy / macro-F1 \\
T2: Activity & Opportunity~\cite{chavarriaga2013opportunity} & On-body activity recognition (DeepConvLSTM) & Weighted / macro-F1 \\
T2: Activity & Ego4D~\cite{grauman2022ego4d} & Egocentric activity recognition & Accuracy / F1 \\
T2: Gesture & Ninapro~\cite{atzori2014ninapro} & EMG gesture / movement classification & Accuracy \\
T2: Gesture & emg2qwerty~\cite{sivakumar2024emg2qwerty} & Large-scale surface-EMG typing & Character error rate \\
T3: Interaction & HOI4D~\cite{HOI4D} & Human-object interaction sensing & Contact / interaction acc. \\
\bottomrule
\end{tabular}
\end{table}


\subsection{Established Tools and Systems}
\label{subsec:p4-sec9-3-established-tools-and-systems}
Commercial wearables have converged on the head and the wrist, and what separates current products is less their sensors than where perception runs and whether anything is shown back to the wearer. Full augmented-reality headsets are covered with the XR platforms in Section~\ref{sec:p4_s10_1}; here the focus is on everyday-wearable glasses, neural input bands, and the research platforms and chips behind them.

\textbf{AI glasses without a display.} By volume, the dominant wearable AI device is a pair of ordinary-looking glasses carrying a camera, microphones, open-ear speakers, and an IMU, with a multimodal assistant running on the paired phone or in the cloud. Meta's line with EssilorLuxottica\footnote{\url{https://www.meta.com/ai-glasses/}} defines the category: the second-generation Ray-Ban Meta (2025) records 3K video from a 12-megapixel camera and runs up to eight hours per charge, and the Oakley Meta Vanguard adapts the design to sport with a 122-degree camera. Display-less models made up 96\% of AI-glasses shipments in the first half of 2026, almost all of them Meta's\footnote{\href{https://counterpointresearch.com/en/insights/global-ai-glasses-shipments-surge-263-percent-yoy-with-display-less-ai-glasses}{Counterpoint Research, global AI glasses shipments, H1 2026}.}, while Xiaomi's AI glasses (2025) and Alibaba's Qwen-based Quark AI glasses\footnote{\url{https://www.quark.cn/}} (November 2025) lead a fast-growing Chinese segment. Google and Samsung previewed Gemini-based Android XR glasses\footnote{\url{https://www.android.com/xr/}} with Gentle Monster and Warby Parker in May 2026, with the first display-less models due in the fall. Spatially, these devices answer questions about what the wearer is looking at one camera frame at a time; most keep no persistent map of the surroundings, so the continuous body and scene tracking of Pivotal Tasks~1 and~2 is still largely absent from shipping products.

\textbf{Display glasses and neural input.} Adding a display makes the glasses an interface, and input then becomes the hard problem. Meta Ray-Ban Display\footnote{\url{https://www.meta.com/ai-glasses/meta-ray-ban-display/}} (September 2025) puts a small monocular display, with a 20-degree field of view, in the right lens and ships with the Meta Neural Band, a surface-EMG wristband that decodes finger gestures and handwriting; it is the commercial form of the generic neuromotor interface described in Pivotal Task-2~\cite{kaifosh2025neuromotor}. In waveguide display glasses, Rokid\footnote{\url{https://global.rokid.com/}} led shipments in the first half of 2026, ahead of Meta and Even Realities\footnote{\url{https://www.evenrealities.com/}}, whose glasses favor a minimal text display over cameras.

\textbf{Research platforms and on-glass compute.} Much of the egocentric data behind the methods of this section was recorded on research glasses rather than consumer products. Meta's Project Aria~\cite{engel2023projectaria}\footnote{\url{https://www.projectaria.com/}} provides glasses with RGB, tracking, and eye-tracking cameras, IMUs, and spatial microphones, together with machine-perception services for trajectories, eye gaze, and hand tracking; its second generation (2025) adds a heart-rate sensor and a contact microphone in the nosepad, alongside on-device perception. The platform underlies Aria Digital Twin~\cite{adt}; Ego-Exo4D~\cite{grauman2024egoexo4d}, which pairs egocentric Aria recordings of skilled activities with synchronized exocentric cameras over 1{,}422 hours; and Nymeria~\cite{ma2024nymeria}, which adds full-body ground-truth motion from a motion-capture suit to 300 hours of everyday activity, the kind of supervision the pose reconstruction of Pivotal Task-1 needs. Ego4D~\cite{grauman2022ego4d} gathered its 3{,}670 hours from several head-mounted cameras, and remains the main benchmark for egocentric activity understanding. On the compute side, Qualcomm's Snapdragon AR1+ Gen~1\footnote{\url{https://www.qualcomm.com/products/mobile/snapdragon/xr-vr-ar}} (2025) runs a one-billion-parameter language model on the glasses without a phone, a first step toward the on-device inference that the latency and power constraints of Pivotal Task-3 call for.

\subsection{Challenges}
\label{subsec:p4-sec9-3-challenges}

Wearable sensing is coupled to the dynamics and variability of the human body in a way that external sensing is not, and that coupling is the source of its hardest problems. The four below track the loop of the System Overview: signal quality at the sense stage, robustness under partial observation, fusion across heterogeneous modalities, and generalization across users.

\paragraph{Signal Instability and Drift.}
On-body signals are inherently unstable. IMU drift, EMG noise, and calibration shifts arise continuously from motion artifacts, sensor displacement, and soft-tissue deformation, and they strike precisely the quantities that pose estimation, trajectory tracking, and dynamics reconstruction integrate over time, so small per-step errors compound into large ones across a long recording.

\paragraph{Incomplete and Intermittent Sensing.}
Sensing is rarely continuous in practice. Sweat, motion, and wireless dropouts cause brief losses of contact and stretches of degraded signal, leaving the system to reconstruct reliable motion from partial and sometimes corrupted segment-level data, a demand that clean-benchmark evaluation rarely exercises.

\paragraph{Heterogeneous Multimodal Fusion.}
The modalities that make wearables powerful are also mismatched. IMU, EMG, pressure, and egocentric-visual streams differ in sampling rate, latency, noise, and coordinate frame, and fusing them under real-world variability, without letting one noisy channel corrupt the whole estimate, remains an open integration problem rather than a solved preprocessing step.

\paragraph{Inter-User Variability.}
Finally, the body itself is the confound. Anatomy, muscle-activation patterns, and movement style vary enough between people that a model tuned on a small or homogeneous population transfers poorly, and the effect is sharpest exactly where wearables matter most, in motion reconstruction, intent prediction, and gait analysis, where individual biomechanics dominate the sensor signature. Personalization, domain adaptation, and larger, more diverse training corpora are the levers for closing this gap, and how far they succeed will decide whether wearable spatial intelligence moves from convincing demonstrations to dependable deployment.

\newpage
\section{Aerial Agents}
\label{sec:p4_s9_4}


% AeroVerse: UAV-Agent Benchmark Suite for Simulating, Pre-training, Finetuning, and Evaluating Aerospace Embodied World Models
% Airvista-ii: An agentic system for embodied uavs toward dynamic scene semantic understanding
% Towards Autonomous UAV Visual Object Search in City Space: Benchmark and Agentic Methodology
% Next-Generation LLM for UAV: From Natural Language to Autonomous Flight

\subsection{Introduction}
\label{subsec:p4-sec9-4-introduction}


Aerial agents (satellites, high-altitude platform stations (HAPS), and unmanned aerial vehicles (UAVs)) acquire rich, dynamic, large-scale observations from the sky and near space and, increasingly, perceive, decide, and act on them autonomously through onboard processing and adaptive control (Table~\ref{tab:air-agents}).
The aerial domain is a demanding testbed for the spatial capabilities surveyed in Part~\ref{part:3}: it stresses perception across extreme scale and viewpoint changes (Chapter~\ref{chapter:3:sec:6}) and decision-making under partial observability and tight latency budgets (Chapter~\ref{sec:reasoning}), across platform tiers that differ in altitude, resolution, endurance, and latency by orders of magnitude.

Satellites supply the global tier. Navigation constellations, namely GPS~\cite{enge1994global,leick2015gps} within the broader Global Navigation Satellite System (GNSS), Russia's GLONASS~\cite{ivanov1992glonass}, the European Union's Galileo, and China's BeiDou~\cite{yang2019introduction,yang2020basic}, provide the positioning, navigation, and timing on which vehicles~\cite{li2022real}, smartphones~\cite{zangenehnejad2021gnss}, network synchronization~\cite{hasan2018gnss}, and every lower aerial tier rely.
Earth-observation satellites operate continuously over years, which suits long-term, wide-area monitoring that neither ground sensors nor short-endurance UAVs can provide: seasonal vegetation dynamics~\cite{malingreau1986global}, wildlife migration~\cite{dobrynin2017integration}, river-basin ecosystem health~\cite{ferreira2013estimating}, deforestation~\cite{tucker2000strategies,finer2018combating}, desertification~\cite{del1998status,collado2002satellite}, biodiversity loss~\cite{turner2015free}, and climate indicators such as land surface temperature, ocean color, and atmospheric composition~\cite{yang2013role,kaufman2002satellite}, together with ecological conservation~\cite{pettorelli2014satellite,boyle2014high}, multispectral and hyperspectral resource discovery~\cite{pettorelli2019satellite}, precision agriculture~\cite{guo2018multi} and crop-yield estimation~\cite{maas1988using,ferencz2004crop}, and the detection of algal blooms~\cite{fang2022global} and pollutant leakage~\cite{pandey2019satellite}.

\begin{table}[!ht]
\centering
\caption{Comparison of satellite, HAPS, and UAV platforms.} 
\label{tab:air-agents}
\small
\setlength{\tabcolsep}{2.5mm}
\begin{tabular}{@{}llll@{}}
\toprule
\textbf{Parameter} & \textbf{Satellite} & \textbf{HAPS} & \textbf{UAV} \\
\midrule
Altitude & 500--36,000 km & 20--50 km & < 10 km \\
Spatial Resolution & m to sub-m & dm to cm & cm-level \\
Temporal Coverage & Hours--days revisit & Persistent (regional) & Real-time (local) \\
Coverage Area & Global/regional & Regional & Localized \\
Endurance & Years & Months & Hours \\
Maneuverability & Minimal & Limited & High \\
Data Latency & High & Medium & Low \\
Payload Capacity & High & Moderate & Low \\
Onboard Processing & Limited & Moderate & High \\
Atmospheric Impact & Minimal & Low & High \\
Deployment Speed & Slow & Medium & Fast \\
Cost Structure & High (low per area) & Medium & Low (high per area) \\
%AI Suitability & Large-scale datasets & Regional monitoring & Real-time inference \\
\bottomrule
\end{tabular}
\end{table}


Differently, UAVs, commonly known as drones, have gained widespread popularity due to their agility, low operational cost, and capability for autonomous operation. Equipped with onboard sensors and processing units, UAVs can perceive their surroundings, plan flight paths, and adapt to dynamic conditions in real time. This unique perspective allows UAVs to capture high-resolution imagery and video from flexible vantage points, enabling frequent revisits and reduced occlusions in complex environments. Compared to satellites, UAVs excel in capturing fine-grained, localized data, and can be rapidly deployed to areas of interest, including hazardous or otherwise inaccessible regions. They now serve aerial photography~\cite{zhao2023ms,shao2022efficient}, cinematography~\cite{pitas2022autonomous,dang2022path}, environmental mapping~\cite{zheng2020university,wang2021each,lin2022joint}, precision agriculture~\cite{goodrich2023placement}, geolocalization~\cite{liu2019lending,shi2019spatial,shi2020optimal,shi2020looking,shi2022beyond,shi2022accurate,shi2022cvlnet,shi2023boosting,song2023learning}, and last-mile delivery~\cite{sorbelli2023wind,dissanayaka2023review}.

Filling the altitude gap between satellite systems and UAVs, emerging platforms such as stratospheric balloons, solar-powered high-altitude platform stations, and small satellite constellations (\eg, CubeSats~\cite{lemmer2017propulsion}) are increasingly being integrated into aerial observation networks.
These high-altitude platform stations (HAPS) offer complementary advantages: HAPS~\cite{gonzalo2018capabilities} provide persistent regional coverage with higher spatial resolution than satellites, while orbital systems~\cite{cunningham2004national} enable global, long-term monitoring at varying temporal and spectral scales.
Together, they form a multi-tiered aerial ecosystem capable of capturing heterogeneous multimedia data, from visible and thermal imagery to hyperspectral and LiDAR scans, across spatiotemporal scales unattainable by any single platform alone.
For instance, solar-powered HAPS such as Airbus Zephyr~\cite{kothe2019flight} or HAWK30 by AeroVironment/HAPSMobile~\cite{fladeland2019demonstrating} can operate autonomously in the stratosphere (about 20 km altitude) for weeks, offering persistent observation and communication relay services with sub-meter resolution, which is superior to most Earth observation satellites and at lower cost than maintaining UAV fleets.

A growing trend is the coordination of heterogeneous aerial agents into cooperative networks, where satellites offer global context, HAPS serve as communication relays or edge computing nodes, and UAVs execute localized missions, enabling real-time, closed-loop intelligence. The utility of these systems is further amplified by advances in artificial intelligence, including machine learning, computer vision, big data analytics, and transfer learning~\cite{wang2024multiple,wang2022learning,zheng2022adaptive}. These technologies empower intelligent processing of massive aerial datasets, enabling automated image and video analysis~\cite{liu2023robust,li2021uav,9854892,zheng2019unsupervised,wang2023context,ju2025video2bev}, real-time object detection and tracking~\cite{zhang2019eye,wen2021detection}, predictive modeling for urban dynamics~\cite{fang2022cognitive,shi2022geometry}, natural language-based control interfaces~\cite{chu2024towards}, and city-scale digital twin simulations~\cite{zheng2022citysim}. Data privacy, airspace integration, energy efficiency, and operational robustness remain open, and are taken up in the Challenges below.


\subsection{System Overview}
\label{subsec:p4_s9_4_overview}
An aerial agent realizes a sensing-to-action loop that spans the multi-tier ecosystem of Table~\ref{tab:air-agents}: satellites supply global context and positioning, HAPS act as persistent regional relays and edge-computing nodes, and UAVs execute localized, low-latency missions at the point of interest. Across this stack, spatial intelligence is exercised through three complementary stages, which constitute the pivotal tasks of this section. Cross-view perception anchors aerial and ground observations to a common geographic frame, so that what a drone sees can be matched to a map or a street-level query (Pivotal Task-1); remote-sensing detection and environmental understanding turn that imagery into structured scene knowledge about objects, layouts, and human activity (Pivotal Task-2); and autonomous navigation and decision making convert perception into flight plans and control actions in unseen environments (Pivotal Task-3). Figure~\ref{fig:p4_s9_4_overview} sketches this pipeline, with a feedback path from action back onto sensing that ties the three stages into one closed loop rather than a feed-forward chain.

\begin{figure}[!ht]
    \centering
    \includegraphics[width=0.98\linewidth]{figures/part4/section9/aerial_overview.pdf}
    \caption{Spatial intelligence for aerial agents as a closed loop from multi-tier aerial sensing (satellite, HAPS, UAV, GNSS, onboard sensors), through three complementary core capabilities, to flight plans with 6-DoF pose targets and downstream flight control. The controller executes the plan; subsequent observations close the sensing loop. The core comprises cross-view geolocalization, detection and environmental understanding, and navigation and decision making. Land-use, water, and infrastructure layers are parallel examples. Platform and scene imagery is illustrative.}
    \label{fig:p4_s9_4_overview}
\end{figure}

\subsection{Pivotal Task-1: Ground-Air Collaborative Perception and Visual Geolocalization}
\label{subsec:p4_s9_4_task1}

A foundational requirement of any aerial system is to place its observations on the map: the integration of ground-level and aerial views for cross-modal understanding enables visual geolocalization, determining the geographic location of a query image by matching it against a reference database. This task anchors the perception stage of the system and is particularly challenging due to the large modality gap between ground street-view images and overhead (aerial or satellite) imagery, an instance of the cross-view correspondence problem treated more generally in Chapter~\ref{chapter:3:sec:6}.
For instance, the \textit{University-1652} dataset~\cite{zheng2020university} has become a benchmark for this task, offering ground, drone and satellite image triplets of campus buildings, thus facilitating research on cross-view matching and localization. Such capabilities are critical for urban search-and-rescue, augmented reality, and infrastructure monitoring, where precise alignment between ground observations and aerial maps is essential.
Progress has come from bridging the modality gap, encoding spatial structure into a viewpoint-invariant representation, and exploiting multiple views per location (Figure~\ref{fig:p4_s9_4_geoloc}).



\textbf{Bridging the modality gap.} Most extant research frames geolocalization as a cross-view image retrieval problem, where the core challenge is to learn a viewpoint-invariant representation that effectively bridges the disparity between images captured from different perspectives. Recent advances in deep learning have led to the widespread application of Convolutional Neural Networks (CNNs) and Vision Transformers (ViTs) for extracting robust visual features. One prominent line of inquiry employs metric learning to construct a shared latent space for images acquired from diverse platforms. Workman~\etal~\cite{workman2015location} show that a CNN pre-trained on the Places dataset~\cite{zhou2014learning} can, by itself, yield highly discriminative features without explicit fine-tuning for the task. The integration of contrastive loss, which minimizes the distance between positive pairs in the feature space, further enhances geolocalization performance~\cite{workman2015wide,lin2015learning}. Liu~\etal~\cite{liu2019stochastic} introduce the Stochastic Attraction and Repulsion Embedding (SARE) loss, which minimizes the Kullback-Leibler divergence between the learned and actual data distributions, and more recently Deuser~\etal~\cite{deuser2023sample4geo} show that hard negative sampling strategies also significantly influence the final discriminative power.

\textbf{Encoding spatial structure.} Another significant body of work addresses the problem of spatial misalignment in ground-to-aerial matching, injecting geometric priors into the representation. Vo~\etal evaluate various network architectures and propose an orientation regression loss to train an orientation-aware network~\cite{vo2016localizing}. Zhai~\etal~\cite{zhai2017predicting} utilize semantic segmentation maps to facilitate semantic alignment, while Hu~\etal~\cite{hu2018cvm} incorporate a NetVLAD layer to extract more discriminative features. Subsequently, Liu~\etal~\cite{liu2019lending} propose a Siamese Network to explicitly integrate spatial cues, such as orientation maps, into the training process, and Shi~\etal~\cite{shi2020optimal} introduce a spatial-aware layer to improve localization performance. Given that each location has multiple training images from different viewpoints, Zheng~\etal~\cite{zheng2020university} train a multi-branch CNN as a base model. Building upon this multi-branch architecture, Wang~\etal~\cite{wang2021each} further partition the feature map into multiple segments to improve matching discriminability, while Lin~\etal~\cite{lin2022joint} incorporate keypoint detection into the learning process. More recently, researchers resort to ViT structures: L2LTR~\cite{yang2021cross} adopts a hybrid CNN-ViT architecture, leveraging CNNs to extract high-level semantic features and ViTs to encode them into global representations; GeoDTR~\cite{zhang2023cross} further enhances geometric reasoning by introducing a ViT-based layout extractor that generates geometry-aware attention maps; and TransGeo~\cite{zhu2022transgeo} advances the paradigm with a two-branch ViT model and a two-stage coarse-to-fine learning strategy.

\begin{figure}[!ht]
    \centering
    \includegraphics[width=0.98\linewidth]{figures/part4/section9/aerial_geoloc.pdf}
    \caption{Cross-view visual geolocalization as viewpoint-invariant metric learning: the illustrated shared encoder embeds ground, drone, and satellite views so that same-place (positive) pairs cluster and different-place negatives separate. Architecture variants may instead use view-specific branches. Arrows within the feature-space sketch indicate attraction and repulsion during representation learning, not geographic motion. Synthetic views illustrate a common site rather than a calibrated image triplet.}
    \label{fig:p4_s9_4_geoloc}
\end{figure}

\begin{table}[t]
\centering
\setlength{\tabcolsep}{3pt}
\caption{Comparison of visual geolocalization datasets. R@K: recall at top $K$; AP: average precision; PR: precision-recall.}
\label{table:geo-dataset}
\small
\resizebox{\linewidth}{!}{
\begin{tabular}{lccccccc}
\toprule
Dataset & Univ.-1652~\cite{zheng2020university} & SUES-200~\cite{zhu2023sues} & CVUSA~\cite{zhai2017predicting} & CVACT~\cite{liu2019lending} & Lin~\etal~\cite{lin2015learning} & Tian~\etal~\cite{tian2017cross} & Vo~\etal~\cite{vo2016localizing} \\
\midrule
\#Train imgs. & 701$\times$71.6 & 120 loc. & 35.5k$\times$2 & 35.5k$\times$2 & 37.5k$\times$2 & 15.7k$\times$2 & 900k$\times$2 \\
Platform & Drone, Ground, Sat. & Drone, Sat. & Ground, Sat. & Ground, Sat. & Ground, 45$^{\circ}$ Aerial & Ground, 45$^{\circ}$ Aerial & Ground, Sat. \\
Imgs./loc. & 54+16.6+1 & drone + 1 & 1+1 & 1+1 & 1+1 & 1+1 & 1+1 \\
Target & Building & Diverse & User & User & Building & Building & User \\
Eval. & R@K, AP & R@K, AP & R@K & R@K & PR, AP & PR, AP & R@K \\
\bottomrule
\end{tabular}}
\end{table}

\textbf{From image pairs to multi-view benchmarks.} The maturation of these methods is closely tied to the datasets that score them (see Table~\ref{table:geo-dataset}). The majority of prior geolocalization datasets are based on image pairs, designed for matching images from two distinct platforms such as terrestrial cameras and satellites. One of the earliest contributions~\cite{lin2015learning} leverages public sources to create 78,000 image pairs spanning a 45-degree bird's-eye perspective and the ground. In a similar vein, Tian~\etal~\cite{tian2017cross} compile image pairs for urban localization, positing that buildings are crucial for this task and thus incorporating building detection into their pipeline. The CVUSA~\cite{zhai2017predicting} and CVACT~\cite{liu2019lending} datasets instead address matching panoramic ground-view images with satellite-view images, which is particularly useful for user localization when GPS data is unavailable. A key distinction is one of objective: the latter two focus on localizing the user who captures the photograph, while the former focus on localizing a specific target within the image, so that multiple views of the target are more beneficial for the target-localization datasets (Lin~\etal, Tian~\etal) because they encourage the model to learn the target's structure. However, most existing datasets typically provide only two views of a given location. The University-1652 dataset~\cite{zheng2020university} addresses this by including a greater number of views per target, thereby fostering more robust viewpoint-invariant features; SUES-200~\cite{zhu2023sues} follows a similar multi-view strategy with drone-captured real-world data. Wang~\etal~\cite{wang2024multiple} further extend University-1652 to multi-weather evaluation (rain, fog, and other complex conditions), while Chu~\etal~\cite{chu2024towards} introduce linguistic annotations and a language-based query setting for a more intuitive user experience.

\textbf{Metrics.} These benchmarks are scored by retrieval metrics that reflect their gallery structure. Several datasets~\cite{liu2019lending,zhai2017predicting,vo2016localizing} adopt Recall@K, defined as 1 if a correct match appears within the top K retrieved images; Recall@K is highly sensitive to the rank of the first correct match and is best suited for test sets containing only a single ground-truth match in the gallery. In contrast, datasets such as University-1652~\cite{zheng2020university} and SUES-200~\cite{zhu2023sues} include multiple ground-truth matches from various viewpoints, where Recall@K fails to comprehensively reflect retrieval performance; they therefore also employ Average Precision (AP), as used by Lin~\etal~\cite{lin2015learning} and Tian~\etal~\cite{tian2017cross}. AP, defined as the area under the Precision-Recall (PR) curve, considers all ground-truth images in the gallery, and recent evaluations usually report both Recall@K and the mean Average Precision (mAP) across all queries.


\subsection{Pivotal Task-2: Remote Sensing Detection and Environmental Understanding}
\label{subsec:p4_s9_4_task2}

Once observations are geolocalized, the next stage turns aerial imagery into structured scene knowledge. Aerial platforms are integral to large-scale environmental monitoring and public safety, facilitating fine-grained object detection in complex scenes, a specialization of the detection and segmentation machinery of Chapter~\ref{chapter:3:sec:6} to the aerial viewpoint. Prominent datasets such as \textit{DOTA}~\cite{xia2018dota}, \textit{DOTA-v2}~\cite{xia2018dota,ding2021object}, and \textit{VisDrone}~\cite{du2019visdrone} provide comprehensive annotations for detecting various objects, including vehicles and pedestrians, under a wide range of weather, illumination, and scale conditions. These benchmarks have spurred the development of robust detection models specifically engineered to address challenges inherent to aerial perspectives, such as extreme aspect ratios, dense object clustering, and diminutive object sizes. A significant application is found in UAV-based wildlife monitoring, where systems are deployed for animal census in conservation areas~\cite{vermeulen2013unmanned}, highlighting the utility of aerial technologies in promoting ecological sustainability. Figure~\ref{fig:p4_s9_4_detection} relates the aerial-specific difficulties of orientation, scale, and density to the detectors built for them and to the human-centric understanding that aerial video enables.

\begin{figure}
    \centering
    \includegraphics[width=0.98\linewidth]{figures/part4/section9/aerial_detection.pdf}
    \caption{Remote-sensing detection from the aerial viewpoint: arbitrary orientation, extreme scale variation, and dense small objects (top) motivate (1) oriented-object detection and (2) multi-scale feature learning, which provide complementary representations for human-centric understanding (bottom). Person re-identification matches identities across views; action examples are alternative classes, not a prescribed sequence; tracking maintains identities across frames. Synthetic scenes illustrate these tasks.}
    \label{fig:p4_s9_4_detection}
\end{figure}

Beyond object detection, as shown in Table~\ref{tab:uav_reid_datasets}, with the proliferation of UAVs in smart-city and other human-centric applications, human behavior understanding~\cite{li2021uav} has emerged as a critical and formidable research challenge, including person re-identification and action recognition. For instance, the UAV person re-identification (re-ID) task~\cite{albaluchi2025uav} is designed to address cross-drone identity matching amid substantial viewpoint and illumination variations, thereby motivating the creation of view-invariant feature representations. In conjunction with resources like \textit{VisDrone}, which contains multi-target tracking sequences from mobile UAVs, these datasets facilitate the design of intelligent video analytics systems capable of performing long-term human behavior analysis and anomaly detection in dynamic environments.


\begin{table}[t]
\centering
\caption{Comparison of UAV-based human and object datasets.}
\label{tab:uav_reid_datasets}
\fontsize{8}{10}\selectfont
\setlength{\tabcolsep}{1.0mm}
\begin{tabular}{L{0.15\textwidth} L{0.26\textwidth} L{0.25\textwidth} L{0.25\textwidth}}
    \toprule
    \textbf{Feature} & \textbf{UAV-Human}~\cite{li2021uav} & \textbf{DOTA-v2}~\cite{xia2018dota,ding2021object} & \textbf{VisDrone}~\cite{du2019visdrone} \\
    \midrule

    \textbf{Primary Focus} & Human Behavior Understanding & Object Detection & Object Detection / Tracking / Crowd Counting \\

    \textbf{Data Type} &  Videos \& Images & Images & Videos \& Images \\

    \textbf{\#Images} & 22,476 frames for pose estimation, 41,290 frames for re-ID, 22,263 frames for attribute recognition & 11,268 & 10,209 \\

    \textbf{\#Videos} & 67,428  & N/A & 288 \\

    \textbf{\#Categories} & 155 action categories & 18 object categories & 10 object categories \\

    \textbf{\#Instances} & 119 subjects for action recognition, 1,144 identities for re-ID & 1,793,658 & > 2.6 million annotated bounding boxes \\

    \textbf{Sensor Modalities} & RGB, Depth, Infrared (IR), Fisheye, Night-Vision & RGB (from Google Earth, aerial) \& Grayscale (from satellites) & RGB \\

    \textbf{Image Resolution} & 1920$\times$1080 (RGB), 640$\times$576 (Depth,IR), 640$\times$480(Fisheye, Night-Vision), & 800$\times$800 to 20,000$\times$20,000 pixels & 2000$\times$1500 (image),  3840$\times$2160 (video) \\

    \textbf{Annotation Format} & Bounding boxes, 17 keypoints for pose & Oriented Bounding Boxes (Quadrilaterals) & Bounding boxes with attributes (\eg, occlusion) \\

    \textbf{Scene Diversity} & Urban and rural, daytime and nighttime & Wide area, diverse landscapes (ports, airports, cities, etc.) & 14 different cities, urban and rural, various weather and lighting \\

    \textbf{Platform} & Drone (DJI Matrice 100) & Satellite, Aerial Platforms & Drone \\

    %\textbf{Unique Features} & Rich multi-modal data for fine-grained human-centric analysis. & Very high-resolution images with oriented bounding boxes for rotated objects. & A comprehensive benchmark with multiple tailored tasks for drone vision. \\
    \bottomrule
  \end{tabular}
\end{table}


Remote sensing object detection is primarily challenged by the prevalence of arbitrarily oriented objects and substantial scale variations~\cite{xia2018dota, ding2021object, sun2022fair1m, liu2016ship, cheng2022anchor, yao2023automated}, and existing research can be broadly organized along two corresponding fronts.
\textbf{Detection of arbitrarily oriented objects.} A significant line of work focuses on developing specialized oriented object detection (OOD) methods. Early approaches often augment horizontal bounding boxes (HBBs) with an angle parameter, employing techniques such as feature refinement in the detector neck~\cite{yang2019scrdet,yang2021r3det}, extraction of rotated regions of interest (RoIs)~\cite{ding2019learning,xie2021oriented}, or designing specific OOD heads~\cite{hou2022shape,han2021redet,ming2021dynamic}. However, these methods frequently suffer from boundary discontinuity and representation ambiguity due to their reliance on angle parameters. To mitigate these issues, subsequent research has explored more flexible \textit{object representations}: some methods describe oriented objects using points~\cite{fu2020point, li2022oriented} or gliding vertexes~\cite{xu2020gliding}, while others model them with Gaussian distributions~\cite{yang2021rethinking, yang2021learning, cheng2022dual} and design corresponding loss functions~\cite{qian2021learning}.
\textbf{Handling substantial scale variations.} Another research direction designs specialized backbones and feature extraction strategies to cope with large scale variations in remote sensing images (RSIs). A common strategy is multi-scale feature learning, achieved through data augmentation~\cite{zhao2019multi, shamsolmoali2021rotation, chen2020stitcher}, multi-scale feature integration~\cite{zhang2019hierarchical,zheng2020hynet,9382268,liu2016ssd}, enhanced feature pyramid networks (FPNs)~\cite{lin2017feature,hou2022refined,zhang2021laplacian,guo2020rotational}, or multi-scale anchor generation~\cite{guo2018geospatial, hou2021self, qiu2019a2rmnet}. Other works focus on ensuring equivalent receptive fields for objects of different orientations~\cite{han2021redet,pu2023adaptive}. More recently, some methods~\cite{Li_2023_ICCV} have employed large kernel convolutions~\cite{liu2022convnet,ding2022scaling,liu2022more} to enlarge the spatial receptive field for large objects, though this can introduce background noise for smaller ones, while others harness multi-scale convolution kernels and attention to address scale variance~\cite{guo2022segnext,dai2023multi,zhang2022multiscale,cai2024poly}.

These method fronts are evaluated on a complementary set of benchmarks and metrics. The VisDrone dataset~\cite{du2019visdrone} provides a comprehensive benchmark for drone-captured imagery, featuring urban and rural scenes with challenges like varying altitudes and occlusions; DOTA-v2~\cite{xia2018dota,ding2021object} focuses on oriented object detection in high-resolution satellite images, emphasizing rotated bounding boxes for objects like planes and ships; and UAV-Human~\cite{li2021uav} targets human behavior understanding from UAVs, with multi-modal sequences for pose and action analysis. Complementary datasets, such as UAVDT~\cite{du2018uavdt} for traffic surveillance and iSAID~\cite{waqas2019isaid} for instance segmentation, broaden the scope of aerial detection tasks and collectively enable cross-dataset fine-tuning to mitigate domain gaps. Evaluation extends standard metrics to address orientation, scale, and density: the primary metric, mean Average Precision (mAP), is computed as the area under the precision-recall curve, typically at an Intersection over Union (IoU) threshold of 0.5 (mAP@0.5) or across 0.5:0.95 for stricter assessment~\cite{lin2014microsoft}, and for oriented detection, as in DOTA-v2, rotated IoU variants ensure alignment with object headings~\cite{xia2018dota}. Per-class and per-size Average Precision exposes imbalances, and small objects ($<32\times32$ pixels) remain the main source of error on drone imagery~\cite{du2019visdrone}. Additional metrics include precision, recall, and F1-score to assess false positives (\eg, background clutter) and misses (\eg, occlusions), while Frames Per Second (FPS) is critical for real-time UAV deployment, targeting 30+ FPS~\cite{yang2023laiyolo}. Video-specific metrics like Multiple Object Tracking Accuracy (MOTA) are used for dynamic scenes in datasets like UAV-Human~\cite{li2021uav}. Modern real-time detectors of the YOLO family~\cite{wang2024yolov9} are increasingly adapted to the aerial setting, though aggregated aerial mAP still trails their COCO-level performance owing to small-object and orientation hardness, keeping open gaps such as multi-modal fusion and zero-shot adaptation active.


\begin{figure}[!ht]
    \centering
    \includegraphics[width=0.98\linewidth]{figures/part4/section9/aerial_navigation.pdf}
    \caption{Autonomous aerial navigation from a language instruction: a UAV observation is grounded into a semantic-topometric graph; the instruction and graph condition an LLM planner that produces pose waypoints with velocity. The displayed route is a top-down projection, with attitude and altitude encoded in the waypoint state. Multi-drone planning is an extension with coordinated altitude separation. SR: success rate; SPL: success weighted by path length; NE: navigation error. Metric glyphs indicate evaluation criteria, not measured scores; scenes and graph geometry are schematic.}
    \label{fig:p4_s9_4_navigation}
\end{figure}


\subsection{Pivotal Task-3: Autonomous Aerial Navigation and Decision Making}
\label{subsec:p4_s9_4_task3}

The final stage closes the loop, converting perception into action. Traditional UAV geolocalization addresses basic aerial navigation by matching a query (\eg, a satellite image or textual description) to a database of historical flight paths or geographic points~\cite{zheng2020university,chu2024towards}. However, a critical distinction exists between such perception-oriented tasks and the overarching goal of autonomous navigation: the latter requires complex decision-making and control to bridge perception to action, drawing on the spatial-reasoning and planning paradigms of Chapter~\ref{sec:reasoning}. In practical scenarios, the objective is often for an aerial agent to autonomously locate a destination, particularly an unexplored target, through environmental exploration and, when available, collaboration with a map. This goal introduces significant challenges in generalization to unseen environments, including novel objects, unmapped regions, and interaction with the environment. For instance, recent benchmarks adapted from general Vision-Language Navigation (VLN) have begun to address this by requiring models to output a detailed flight plan with 6-DoF control signals~\cite{wang2024towards}. Figure~\ref{fig:p4_s9_4_navigation} follows the problem from the perception-to-action gap to foundation-model-driven VLN and multi-agent extensions.



\begin{table}[!ht]
\small
\centering
\caption{Comparison of recent aerial VLN datasets.}
\label{tab:aerial_datasets}
\begin{tabular}{@{}p{2.5cm}lp{2.5cm}p{2cm}p{3cm}p{3.5cm}@{}}
\toprule
\textbf{Dataset} & \textbf{Yr} & \textbf{Env.} & \textbf{Size} & \textbf{Modalities} & \textbf{Key Features} \\
\midrule
UAVA~\cite{albanis2020dronepose} & 2020 & Indoor simulated (MP3D) & 90 traj. & RGB, depth, flow, ego/exo views & Photorealistic UAV assist; vision/robotics tasks; gamified traj. \\
AVDN~\cite{fan2022aerial} & 2022 & Outdoor simulated (xView) & 3,064 traj. with dialog & RGB, dialog & Async human dialogs; attn. for interactive nav. \\
AerialVLN~\cite{liu2023aerialvln} & 2023 & Outdoor simulated (cities) & 8,446 traj. and 25,338 instructions & RGB, lang. instr. & Synthetic VLN; landmark nav in urban scenes. \\
CityNav~\cite{lee2024citynav} & 2024 & Real-world (UK cities) & 32,637 traj./instr. & 3D point clouds, RGB, geo-lang & Real 3D scans (4.65 km²); human demos; scalable VLN. \\
UAV-Need-Help~\cite{wang2024towards} & 2024 & Outdoor simulated (OpenUAV) & 12,149 traj. & RGB, lang. instr. & Continuous 6-DoF traj.; target-oriented. \\
OpenFly~\cite{gao2025openfly} & 2025 & Outdoor real/sim (18 scenes) & 100k traj. & 3D point clouds, RGB, lang. instr. & Large-scale; real-to-sim transfer. \\
AirV2X~\cite{gao2025airv2x} & 2025 & Outdoor simulated (CARLA+AirSim) & 6.73 hours of UAV-ground traj. & 3D point clouds, RGB & V2X collab; drone strategies; reduces occlusions. \\
\bottomrule
\end{tabular}
\end{table}


\textbf{From visual perception to language-guided planning.} Recent years have witnessed significant progress in UAV navigation, with a primary focus on visual perception~\cite{loquercio2018dronet,kang2019generalization,giusti2015machine,smolyanskiy2017toward,fan2020learn,bozcan2020air,majdik2017zurich} and collision avoidance~\cite{xu2025navrl,singla2019memory}. Building on this aerial-perception foundation, Vision-and-Language Navigation (VLN) for UAVs remains an emerging area of research. To address it, recent efforts have leveraged Large Multi-modal Models (LMMs) to harness general knowledge for navigation planning: STMR~\cite{gao2024aerial} builds upon AerialVLN~\cite{liu2023aerialvln} with a zero-shot LLM-based framework utilizing a Semantic-TopoMetric Representation for enhanced spatial reasoning. A parallel line reframes UAV autonomy as a full agentic loop rather than one-shot instruction following. AeroVerse~\cite{yao2024aeroverse} assembles a simulator, real-virtual pre-training corpora, and five downstream tasks, \ie, scene awareness, spatial reasoning, navigational exploration, task planning, and motion decision, to train aerospace embodied world models, and proposes SkyAgentX as a unified perception-reasoning-navigation-planning agent; AirVista-II~\cite{lin2025airvista} pursues zero-shot dynamic-scene semantic understanding via agent-based task scheduling and adaptive keyframe selection; and PRPSearcher~\cite{ji2025uavsearch} casts aerial visual object search over the CityAVOS benchmark as a perception-reasoning-planning problem solved by a multimodal LLM agent. Complementary systems such as NeLV~\cite{yuan2025nelv} operationalize the language-to-flight pipeline end to end, chaining an LLM parser, route and path planners, and a control platform across short-, medium-, and long-range missions. Beyond single-agent setups, some works have begun to explore multi-drone systems; for instance, Wu~\etal~\cite{wu2025aeroduo} deploy two drones flying at different altitudes to facilitate planning.

\textbf{Datasets and simulators.} As shown in Table~\ref{tab:aerial_datasets}, several foundational frameworks and datasets have been established. One of the early works, \ie, UAVA~\cite{albanis2020dronepose}, leverages the panorama images in the Matterport3D dataset~\cite{chang2017matterport3d} to render the indoor environment, while harnessing the camera viewpoints as physically available flight anchors to generate trajectories. AerialVLN~\cite{liu2023aerialvln} introduces a large-scale dataset with an effective baseline, and further infrastructure contributions include datasets such as CityNav~\cite{lee2024citynav}, AVDN~\cite{fan2022aerial}, and OpenFly~\cite{gao2025openfly}. Notably, UAV-Need-Help~\cite{wang2024towards} provides a realistic UAV dynamics simulator, \ie, OpenUAV, alongside its evaluation protocol; OpenUAV is further used in HaL-13k~\cite{wu2025aeroduo} to simulate two drones at different altitudes. Additional observation agents, such as road infrastructures and vehicles with depth sensing, have also been incorporated to enrich environmental understanding~\cite{gao2025airv2x}.

\textbf{Metrics.} Evaluation metrics for UAV aerial navigation, especially VLN, measure how well a model navigates a drone based on visual inputs and language instructions. Common metrics include Success Rate (SR), which checks if the UAV reaches the target location within a specified distance, and Success weighted by Path Length (SPL), which balances success with how efficiently the drone travels (shorter paths score higher). Path Length (PL) measures the total distance traveled, while Navigation Error (NE) calculates the distance between the final position and the goal. Together these metrics assess both accuracy (reaching the target) and efficiency (taking an optimal route), providing a clear way to compare model performance across datasets like UAVA, AerialVLN, and CityNav.


\subsection{Established Tools and Systems}
\label{subsec:p4-sec9-4-established-tools-and-systems}
Beyond the methods and benchmarks above, the aerial domain rests on a real-world layer of simulators, platforms, and infrastructure that turns research prototypes into deployable systems; we group it into three tiers.

\textbf{Simulation environments and benchmarks.} Because flight data is costly and safety-critical to collect, high-fidelity simulation underpins much of the field. The OpenUAV simulator released with UAV-Need-Help~\cite{wang2024towards} supplies realistic UAV dynamics and a continuous 6-DoF evaluation protocol, and is reused to simulate multi-altitude drone teams in HaL-13k~\cite{wu2025aeroduo}; OpenFly~\cite{gao2025openfly} targets large-scale real-to-sim transfer across diverse scenes; and AerialVLN~\cite{liu2023aerialvln} pairs a city-scale synthetic environment with a navigation baseline. Photorealistic indoor rendering for UAV tasks is provided through Matterport3D~\cite{chang2017matterport3d} in UAVA~\cite{albanis2020dronepose}, while V2X-style multi-agent simulation combines CARLA and Microsoft AirSim\footnote{\url{https://microsoft.github.io/AirSim/}} in AirV2X~\cite{gao2025airv2x}.

\textbf{Commercial platforms and vendors.} On the hardware side, consumer and enterprise drones from DJI\footnote{\url{https://www.dji.com/}} (whose Matrice series collects datasets such as UAV-Human) and autonomy-focused vendors such as Skydio\footnote{\url{https://www.skydio.com/}} dominate the UAV tier. In the stratospheric tier, solar-powered HAPS such as Airbus Zephyr~\cite{kothe2019flight} and AeroVironment/HAPSMobile's HAWK30~\cite{fladeland2019demonstrating} provide persistent regional coverage, complemented by small-satellite constellations built on CubeSat platforms~\cite{lemmer2017propulsion} and broader HAPS programs~\cite{gonzalo2018capabilities}.

\textbf{Connectivity and positioning infrastructure.} Aerial autonomy ultimately depends on shared positioning and communication backbones: the GNSS constellations introduced above (GPS~\cite{enge1994global}, GLONASS~\cite{ivanov1992glonass}, Galileo, and BeiDou~\cite{yang2019introduction}) supply global timing and localization, while space-based broadband relays such as Starlink~\cite{michel2022first} and orbital observation systems~\cite{cunningham2004national} increasingly serve as the connectivity layer for coordinated, real-time aerial operations.


\subsection{Challenges}
\label{subsec:p4-sec9-4-challenges}
Aerial autonomy is limited less by any single model than by the gaps between platforms, scales, and operating conditions. Four threads follow the pipeline in order: fusing heterogeneous platforms into a coherent situational model, perceiving small objects and complex activities reliably, navigating unseen and GPS-denied environments, and meeting the ethical and operational constraints under which any of this can be deployed.

\paragraph{Heterogeneous-Platform Fusion.}
A primary obstacle lies in integrating aerial agents with heterogeneous platforms for robust environmental perception. Coordinating a UAV with ground-based sensors or satellite systems to construct a cohesive situational-awareness model remains complex, particularly in dynamic or cluttered settings where the platforms differ in resolution, latency, and viewpoint (Table~\ref{tab:air-agents}). Reconciling these asynchronous, multi-resolution streams into a single consistent representation, without one modality's noise corrupting the fused estimate, is the standing difficulty, and it motivates the V2X-style cooperative-perception efforts that incorporate road infrastructure and ground vehicles as additional observers.

\paragraph{Small-Object and Complex-Activity Understanding.}
Aerial perception is defined by extremes of scale and density. The recognition of small objects, such as identifying debris in disaster zones or distinguishing wildlife in ecological surveys, demands high-resolution imaging and detectors robust to dense clustering, which often strain onboard computational resources. Understanding complex activities, such as discerning human behavior during search-and-rescue missions, compounds the problem by requiring nuanced spatio-temporal and contextual analysis that current systems struggle to deliver consistently. Closing this gap calls for detectors and activity models that remain accurate at extreme scale variation while fitting a UAV's tight compute and power budget.

\paragraph{Navigation in Unseen and GPS-denied Environments.}
Autonomous operation in unseen or GPS-denied environments, such as dense urban canyons or subterranean spaces, further complicates the action stage, as agents must rely on real-time mapping and obstacle avoidance without pre-existing data. The core requirement is generalization: a navigation policy must extend to novel objects, unmapped regions, and genuinely interactive scenarios, rather than replaying matched trajectories. This is precisely the gap that foundation-model-driven VLN and multi-altitude cooperative planning aim to narrow, but validating such behavior to a safety standard in the open world remains open.

\paragraph{Privacy, Airspace, and Operational Robustness.}
Beyond perception and navigation, several ethical and operational challenges persist. Data-privacy concerns arise when aerial agents capture high-fidelity multimedia data, potentially compromising individual rights without stringent safeguards, motivating privacy-preserving approaches such as federated learning. Airspace integration poses another hurdle, as regulatory frameworks struggle to accommodate the proliferation of aerial agents alongside traditional aviation, necessitating harmonized protocols to ensure safety. Energy efficiency remains a critical bottleneck, since prolonged missions demand lightweight yet powerful power sources and motivate renewable solutions such as solar-powered propulsion. Finally, operational robustness is challenged by environmental uncertainties, such as extreme weather or electromagnetic interference, that can disrupt system reliability. Advances in artificial intelligence and sensor fusion that address these constraints, rather than perception accuracy alone, are what will let aerial agents move from bounded demonstrations toward dependable, large-scale deployment.


\newpage

\section{Warehouse and Industrial Logistics}
\label{sec:p4_s9_5}



\subsection{Introduction}
\label{subsec:p4-sec9-5-introduction}

Warehouse and industrial logistics is one of the most mature real-world testbeds for the spatial capabilities surveyed in Part~\ref{part:3}: it demands perception-to-action pipelines that interpret complex, dynamic environments such as dense shelving systems populated with diverse stock keeping units (SKUs), pallets, forklifts, and human workers. 
These systems construct and maintain spatial memory representations (semantic maps, object and pose states) to support downstream planning and manipulation for tasks including inventory counting, order picking, pallet handling, and coordinated fleet operations. 
Perception and reasoning themselves are treated in Chapter~\ref{chapter:3:sec:6} and Chapter~\ref{sec:reasoning}; logistics adds throughput targets, engineered but crowded layouts, and people sharing the floor with fleets of robots. 
Performance is judged in terms of throughput, order accuracy, safety, and operational cost, especially under high SKU diversity, visual clutter, elongated aisles, and human-robot collaboration. 
Dense shelving and manipulation in clutter remain the hardest cases, as the Amazon Picking Challenge made plain~\cite{Correll2016AmazonChallenge}, and they demand that perception, pose estimation, mapping, and multi-agent planning be integrated rather than tuned in isolation.

\begin{figure}[!ht]
    \centering
    \includegraphics[width=0.98\linewidth]{figures/part4/section9/warehouse_overview.pdf}
    \caption{Spatial intelligence for warehouse and industrial logistics: RGB-D, LiDAR, and RFID observations support five complementary capabilities for picking and pallet transport, and execution outcomes update maps and plans. Scene images and diagrams are illustrative.}
    \label{fig:p4_s9_5_overview}
\end{figure}

\subsection{System Overview}
\label{subsec:p4-sec9-5-system-overview}

The complete warehouse automation pipeline, sketched in Figure~\ref{fig:p4_s9_5_overview}, integrates five complementary capabilities into a closed loop. 
\textbf{Perception} converts raw multi-sensor streams into semantic object information (what is present, where, and in what pose); 
\textbf{localization and mapping} grounds that information in a metric spatial map and keeps it consistent over time; 
\textbf{reasoning and planning} consume these representations to allocate tasks and generate optimized fleet schedules; 
\textbf{manipulation and control} physically execute the resulting picking and transport actions; 
and a \textbf{human-robot collaboration and safety} layer adapts behavior dynamically in mixed human-robot environments. 
Integration is achieved via middleware such as ROS\,2 and industrial protocols like Open Platform Communications Unified Architecture (OPC-UA), while digital twin platforms such as NVIDIA Omniverse\footnote{https://www.nvidia.com/en-gb/omniverse/} support simulation and validation before deployment. 
The five capabilities are the pivotal tasks below, taken in pipeline order.



\subsection{Pivotal Task-1: Perception and Recognition}
\label{subsec:p4-sec9-5-pivotal-task-1-perception-and-recognition}
Perception constitutes the sensory foundation of the pipeline: it transforms raw sensor data into actionable knowledge, reusing the detection, segmentation, and pose-estimation machinery of Chapter~\ref{chapter:3:sec:6} but specializing it to the warehouse Operational Design Domain. 
Core capabilities include detecting and classifying objects (packages, pallets, tools), estimating their 3D poses, and tracking their motion and state over time. 
Unlike outdoor environments, warehouses impose distinct perceptual challenges: dense clutter, heavy occlusion, visually similar items, reflective packaging, narrow aisles, and continually changing layouts, which together make robust perception a central and ongoing research focus \cite{Tubis2023IntelligentWarehouse,redmon2016yolo}. 
Formally, perception is a structured mapping from a synchronized multimodal stream $X_t = \{I_t^{\text{RGB}}, I_t^{\text{Depth}}, P_t^{\text{LiDAR}}, R_t^{\text{RFID}}\}$ (together with calibration parameters and a prior map $M$) to a scene graph $G_t = \{o_i, c_i, \mathbf{b}_i, \mathbf{p}_i, \Sigma_i\}$, in which each detected object $o_i$ carries a class label $c_i$, a bounding box or segmentation mask $\mathbf{b}_i$, an estimated pose $\mathbf{p}_i$, and an uncertainty measure $\Sigma_i$ that grounds subsequent manipulation, navigation, and decision making. 
Figure~\ref{fig:p4_s9_5_perception} sketches this representation stack.

\begin{figure}[!ht]
    \centering
    \includegraphics[width=0.98\linewidth]{figures/part4/section9/warehouse_perception.pdf}
    \caption{The warehouse perception stack: 2D detection and segmentation, 6D object pose $T_i$, and a structured scene graph whose nodes carry identity, class, box or mask, pose, and uncertainty $\Sigma_i$. RFID can support identity association; masks and axes are schematic, not calibrated estimates.}
    \label{fig:p4_s9_5_perception}
\end{figure}

\textbf{Detection and recognition} form the base layer. 
Convolutional and transformer-based architectures, \ie, YOLO \cite{redmon2016yolo}, RetinaNet \cite{lin2017focal}, and the transformer-based DETR with deformable attention \cite{Zhu2021DeformableDETR}, have steadily improved accuracy and efficiency under clutter, and are scored on COCO \cite{lin2014microsoft}, whose mean Average Precision (mAP) broken down by object size is particularly informative for the small-SKU and barcode-label cases that dominate warehouse failure modes. 
\textbf{Pose estimation} lifts these 2D detections into the metric space that manipulation requires: PoseCNN \cite{xiang2017posecnn}, DenseFusion \cite{Wang2019DenseFusion}, and CosyPose \cite{Labbe2020CosyPose} fuse visual and geometric cues to infer six-degree-of-freedom (6D) orientation and position even under occlusion, and are evaluated on YCB-Video \cite{xiang2017posecnn} for textured household and industrial objects and on T-LESS \cite{Hodan2017TLESS} for textureless, industry-relevant parts, with Average Distance (ADD) the standard accuracy measure and MOTA/MOTP reporting tracking quality in dynamic scenes. 
The Amazon Picking Challenge \cite{Correll2016AmazonChallenge} couples perception with manipulation in real warehouse scenarios, stressing generalization across novel object sets and thereby anticipating the manipulation task of §\ref{subsec:p4_s9_5_manip} below. 
Representative benchmarks across the section are consolidated in Table~\ref{tab:p4_s9_5_benchmarks}. In industrial practice, perception pipelines are increasingly standardized and modularized; the deployed software stacks that realize them are surveyed in the Established Tools and Systems subsection.

\subsection{Pivotal Task-2: Localization and Mapping}
\label{subsec:p4-sec9-5-pivotal-task-2-localization-and-mapping}
If perception answers \emph{what is around the robot}, localization and mapping answer \emph{where the robot is} and \emph{what the persistent structure of the warehouse is}, forming the spatial backbone on which fleet coordination rests and drawing on the same scene-construction techniques as Chapter~\ref{chapter:3:sec:6}. 
Automated guided vehicles (AGVs) and autonomous mobile robots (AMRs) require centimeter-level accuracy to navigate narrow aisles, dock precisely at shelves, and operate safely alongside workers. 
The problem is conceptually the SLAM problem of autonomous driving, yet it faces warehouse-specific constraints: highly repetitive structures (rows of identical shelves), dynamic obstacles (workers, pallets, forklifts), and low-light or GPS-denied indoor conditions. 
Formally, a sequence of odometry and sensor measurements $Z_{1:t} = \{\text{IMU}, \text{LiDAR}, \text{cameras}, \text{wheel encoders}, \text{UWB beacons}\}$, optionally supplemented by a prior map $M_0$ or fiducial anchors, is fused into a pose graph $\mathcal{X}_t = \{x_1, \ldots, x_t\}$ that estimates the trajectory and an evolving semantic map $M_t$ fusing geometric structure with storage zones and restricted areas. 
Beyond navigation, this map underpins the digital twins (spatial models synchronized with the facility whose simulations feed back into its operation; §\ref{sec:p4_s11_1}) that increasingly govern modern facilities \cite{Cadena2016SLAMSurvey}. Figure~\ref{fig:p4_s9_5_localization} casts the task as a multi-sensor fusion problem.

\begin{figure}
    \centering
    \includegraphics[width=0.98\linewidth]{figures/part4/section9/warehouse_localization.pdf}
    \caption{Localization and mapping in the warehouse as a factor-graph fusion problem: IMU, LiDAR, camera, wheel-encoder, and UWB measurements constrain pose states, optionally with a prior map or beacon anchors, producing poses over warehouse geometry, storage zones, and restricted areas. Repetitive shelving can cause false place matches, so only verified loop candidates add constraints.}
    \label{fig:p4_s9_5_localization}
\end{figure}

A family of \textbf{SLAM algorithms} has been tailored to these conditions. ORB-SLAM3 \cite{Campos2021ORBSLAM3} extends visual SLAM with inertial and multi-camera support for accurate visual-inertial localization; LIO-SAM \cite{Shan2020LIOSAM} provides real-time LiDAR--IMU fusion through factor-graph optimization that stays robust under fast motion and partial occlusion; and RTAB-Map \cite{Labbe2019RTABMap} supports online RGB-D mapping and loop-closure detection with memory management suited to long-term operation, all expressible within the unified factor-graph framework of GTSAM \cite{Dellaert2012GTSAM}. 
\textbf{Multi-modal fusion} is the common thread that holds accuracy together when any single modality degrades, and is also the dominant emerging research direction, alongside learning-based estimation (neural depth prediction, differentiable bundle adjustment, and learned place-recognition embeddings) and the migration of SLAM computation to the edge for fleet-scale latency. 
These capabilities are evaluated on KITTI \cite{Geiger2013KITTI} for outdoor visual and LiDAR odometry and loop closure, on the indoor visual-inertial EuRoC MAV dataset \cite{Burri2016EuRoC} whose challenging motion and illumination mirror warehouse conditions, and on the high-fidelity Replica reconstructions \cite{straub2019replica} for dense mapping and semantic labeling, with Absolute and Relative Pose Error (APE/RPE) and loop-closure precision/recall the standard metrics and, for warehouse use, semantic-zone IoU and map completeness (Table~\ref{tab:p4_s9_5_benchmarks}).


\begin{figure}[!ht]
    \centering
    \includegraphics[width=0.98\linewidth]{figures/part4/section9/warehouse_planning.pdf}
    \caption{Two families for warehouse reasoning and planning over shared task, fleet, and map inputs: combinatorial optimization (assignment and routing) and learning-based coordination (per-agent policies), optionally combined in hierarchical or hybrid designs. Allocations, paths, and schedules are illustrative.}
    \label{fig:p4_s9_5_planning}
\end{figure}



\subsection{Pivotal Task-3: Reasoning and Planning}
\label{subsec:p4-sec9-5-pivotal-task-3-reasoning-and-planning}
Reasoning and planning form the decision-making core of warehouse automation, transforming the perceptual and localization outputs into task assignments and motion plans, and instantiating the spatial-reasoning and planning paradigms of Chapter~\ref{sec:reasoning}. 
This layer determines which robot executes which task, how tasks are sequenced, and how routes and schedules are optimized for safe, efficient, coordinated operation. 
In contrast to autonomous driving, where planning typically focuses on individual vehicles, warehouse reasoning must address multi-agent coordination across fleets of dozens or hundreds of robots in constrained indoor space \cite{Veloso2015CMUCoBot}, maximizing throughput and minimizing latency while maintaining safety under dynamic obstacles, human interactions, and stochastic task arrivals. 
Formally, the perceived states $G_t$, a task list $\mathcal{T}$, the fleet state $F_t$ (positions, capabilities, battery levels), and the warehouse map $M_t$ are mapped to a task-allocation matrix $A$, optimized motion paths $\mathcal{R}$, and temporal schedules $\pi$ that jointly specify when and how tasks execute. Figure~\ref{fig:p4_s9_5_planning} contrasts the two solution families.


The first family is \textbf{combinatorial optimization}. Traditional methods (the Hungarian algorithm, branch-and-bound, and mixed-integer linear programming) remain widely used for task allocation and scheduling in small- to medium-scale deployments, while Vehicle Routing Problem (VRP) formulations \cite{Cordeau2007VRP} and their multi-depot and time-window variants are central to modeling fleet movement. The second family is \textbf{learning-based coordination}: reinforcement learning (RL) and multi-agent RL (MARL) handle dynamic environments and real-time uncertainty \cite{Sartoretti2019PRIMAL,samvelyan2019smac}, with toolkits such as Google OR-Tools \cite{Perron2019ORTools} and PyMARL enabling rapid experimentation across both families. 
Emerging research couples the two: hierarchical and hybrid frameworks combine symbolic reasoning with learned components, graph neural networks and differentiable optimization accelerate multi-agent coordination and transfer policies across layouts, and digital-twin integration enables simulation-informed, predictive re-planning, increasingly under multi-objective criteria that balance throughput, energy, and human safety. 
Planning systems are benchmarked on Solomon's VRP suite \cite{Solomon1987VRP} for static routing and scheduling and on the Open Motion Planning Library (OMPL) \cite{Sucan2012OMPL} for continuous-space motion planning, complemented by synthetic multi-robot testbeds; common metrics are throughput (items/hour), task-completion latency, total travel distance or time, collision rate, and re-planning latency (Table~\ref{tab:p4_s9_5_benchmarks}).

\subsection{Pivotal Task-4: Manipulation and Control}
\label{subsec:p4_s9_5_manip}
Manipulation constitutes the physical execution layer, enabling robots to interact directly with objects, containers, and infrastructure such as bins, shelves, and conveyors. 
In industrial logistics this typically involves robotic arms performing grasping, placing, sorting, or palletizing, as well as AGVs docking to pallets or workstations. Manipulation bridges high-level symbolic or geometric plans with continuous motor control, and therefore demands both precision and robustness to the inherent variability in packaging, geometry, and material properties \cite{Levine2016DeepGrasp}; the grasping and control techniques here connect directly to the robotic-manipulation treatment of §\ref{sec:p4_s9_1} and the planning paradigms of Chapter~\ref{sec:reasoning}. 
Formally, estimated object poses $\mathbf{p}_i$ and task descriptions (grasp targets, placement goals, trajectory constraints) are mapped to low-level motor commands $\mathbf{u}_t$ that ensure stable grasping, precise placement, and compliant, collision-free motion. Figure~\ref{fig:p4_s9_5_manipulation} contrasts modular grasp-and-plan systems with learned visuomotor control.

\begin{figure}
    \centering
    \includegraphics[width=0.98\linewidth]{figures/part4/section9/warehouse_manipulation.pdf}
    \caption{Warehouse manipulation through (a) a modular route, where grasp candidates scored by analytic metrics and learned quality (Dex-Net) feed sampling-based planning (RRT*) and trajectory tracking, and (b) learned visuomotor control from camera observations. Action spaces vary across systems, and the routes can be combined.}
    \label{fig:p4_s9_5_manipulation}
\end{figure}

\textbf{Grasp synthesis} anchors the methodology: Dex-Net combines synthetic grasp data and analytic robustness metrics to train a grasp-quality predictor for candidate selection \cite{Mahler2017DexNet}, while learned visual control can predict grasp success from camera images and candidate task-space motions, then select actions for adaptive manipulation \cite{Levine2016DeepGrasp}. 
\textbf{Motion planning} then realizes the chosen grasp: sampling-based algorithms such as RRT* \cite{Karaman2011RRTStar} remain fundamental for generating collision-free trajectories in high-dimensional configuration spaces, exposed through frameworks like MoveIt \cite{Sucan2010MoveIt}. 
These capabilities are benchmarked on the Amazon Picking Challenge \cite{Correll2016AmazonChallenge}, which established a standardized testbed for picking under clutter and uncertainty, and on the YCB Object and Model Set \cite{Calli2015YCB}, a diverse collection of everyday items with physical models for grasping and manipulation; grasp success rate, cycle time per pick, placement accuracy, and task-completion rate are the headline metrics (Table~\ref{tab:p4_s9_5_benchmarks}). 
Hybrid control frameworks that combine analytical models with data-driven adaptation are emerging as the practical route to generalizing across variable geometries and grasp contexts, complemented by tactile sensing and force control for slip detection and contact regulation.

\subsection{Pivotal Task-5: Human-Robot Collaboration and Safety}
\label{subsec:p4-sec9-5-pivotal-task-5-human-robot-collaboration-and-safety}
Even in highly automated warehouses, human workers remain indispensable for complex, dexterous, or exception-handling tasks that require contextual understanding and flexible decision making, so the final stage of the pipeline is safe and efficient human-robot collaboration (HRC). 
HRC systems must balance productivity with stringent safety guarantees, enabling mixed teams (robots handling repetitive or heavy-load work, humans managing fine-grained, judgment-based activities) to share workspaces without physical or operational interference \cite{Lasota2015HRC}. 
The task maps multimodal human-detection and intent signals (pose, gesture, gaze, and trajectory data from vision, LiDAR, wearables, or motion capture) to adaptive robot behaviors such as dynamic speed modulation, path re-routing, stop-and-wait actions, or task reallocation in response to predicted human movement. Figure~\ref{fig:p4_s9_5_hrc} sketches this predict-then-adapt loop.

\begin{figure}[!ht]
    \centering
    \includegraphics[width=0.98\linewidth]{figures/part4/section9/warehouse_human_robot_collaboration.pdf}
    \caption{The human--robot collaboration loop: multimodal human sensing and intent recognition inform proximity-aware planning, which selects slowing, rerouting, or stopping. Clearance regions and routes are illustrative, not certified safety zones.}
    \label{fig:p4_s9_5_hrc}
\end{figure}

\textbf{Intent recognition} lets robots anticipate human motion: sequence models based on hidden Markov models or recurrent networks predict trajectories and proactively adjust behavior \cite{Mainprice2016Intent}, while \textbf{proxemics-aware planning} maintains socially acceptable distances and re-routes trajectories under human proximity \cite{dragan2013legibility}, with multimodal sensing (vision-based detection fused with wearable or RFID localization) improving reliability. 
These principles are embodied in force- and speed-limited collaborative robots and safety-rated AMRs (see Established Tools and Systems), governed by international standards such as ISO 10218 and ISO/TS 15066. 
Evaluation combines quantitative safety metrics with qualitative interaction assessment: human-motion corpora such as Human3.6M~\cite{IonescuPOS14} supply sequences for intent prediction, and key indicators include intervention frequency (safety-triggered stops per hour), near-miss rate, human-detection accuracy, and compliance with ISO-defined safety zones.

\begin{table}[!th]
\centering
\fontsize{8}{10}\selectfont
\setlength{\tabcolsep}{1.0mm}
\caption{Representative benchmarks for spatial intelligence in warehouse and industrial logistics, spanning the five pivotal tasks.}
\label{tab:p4_s9_5_benchmarks}
\begin{tabular}{L{0.21\textwidth} L{0.21\textwidth} L{0.31\textwidth} L{0.18\textwidth}}
\toprule
\textbf{Pivotal task} & \textbf{Benchmark} & \textbf{Representative method} & \textbf{Headline metric} \\
\midrule
Perception (detection) & COCO~\cite{lin2014microsoft} & YOLO~\cite{redmon2016yolo} / DETR~\cite{Zhu2021DeformableDETR} & mAP \\
Perception (6D pose) & YCB-Video~\cite{xiang2017posecnn} & DenseFusion~\cite{Wang2019DenseFusion} & ADD \\
Perception (textureless pose) & T-LESS~\cite{Hodan2017TLESS} & CosyPose~\cite{Labbe2020CosyPose} & ADD / VSD \\
Localization \& mapping & KITTI~\cite{Geiger2013KITTI} & Visual / LiDAR SLAM & APE / RPE \\
Localization \& mapping & EuRoC MAV~\cite{Burri2016EuRoC} & ORB-SLAM3~\cite{Campos2021ORBSLAM3} (visual-inertial) & APE \\
Mapping (dense) & Replica~\cite{straub2019replica} & RTAB-Map~\cite{Labbe2019RTABMap} & IoU / completeness \\
Reasoning \& planning & Solomon VRP~\cite{Solomon1987VRP} & OR-Tools~\cite{Perron2019ORTools} & total cost / distance \\
Reasoning \& planning & OMPL~\cite{Sucan2012OMPL} & Sampling-based planners & planning time / success \\
Manipulation & Amazon Picking~\cite{Correll2016AmazonChallenge} & Integrated perception--grasping team pipelines & pick success / time \\
Manipulation & YCB Object Set~\cite{Calli2015YCB} & Dex-Net~\cite{Mahler2017DexNet} & grasp success rate \\
\bottomrule
\end{tabular}
\end{table}

\subsection{Established Tools and Systems}
\label{subsec:p4-sec9-5-established-tools-and-systems}
The pivotal tasks above are realized in production through a layered ecosystem of simulators, middleware, optimization services, and robot hardware that together constitute the real-world delta over the methods of Part~\ref{part:3}.

\textbf{Simulation and middleware} form the integration substrate. Digital-twin and physics platforms such as NVIDIA Omniverse Isaac Sim \cite{NVIDIA2021Isaac} and Gazebo provide realistic environments for developing and validating perception, localization, and planning pipelines before deployment, while perception SDKs such as NVIDIA Isaac, the Intel RealSense SDK, and Zebra RFID/barcode platforms  supply end-to-end sensing support in real facilities. Interoperability across heterogeneous hardware is provided by ROS~\cite{Quigley2009ROSIndustrial}, ROS\,2, and ROS-Industrial, with commercial localization stacks such as SLAMcore\footnote{\url{https://www.slamcore.com/}} and BlueBotics ANT\footnote{\url{https://bluebotics.com/}} delivering optimized multi-sensor fusion, map management, and fleet coordination for AGVs and AMRs.

\textbf{Warehouse management and fleet orchestration} sit at the enterprise layer. Large-scale optimization services such as SAP Extended Warehouse Management\footnote{\url{https://www.sap.com/products/scm/extended-warehouse-management.html}}, Blue Yonder WMS\footnote{\url{https://blueyonder.com/}}, and Manhattan Associates\footnote{\url{https://www.manh.com/}} integrate task allocation and scheduling with enterprise resource planning, while open toolkits such as Google OR-Tools \cite{Perron2019ORTools} and PyMARL support the underlying combinatorial and learning-based planners.

\textbf{Robot hardware and manipulation} close the loop physically. High-throughput pick-and-place arms from ABB\footnote{\url{https://new.abb.com/products/robotics}}, FANUC\footnote{\url{https://www.fanuc.com/}}, and KUKA\footnote{\url{https://www.kuka.com/}} are widely deployed and integrated through ROS-based middleware, mobile manipulators such as those from Fetch Robotics\footnote{\url{https://fetchrobotics.com/}} unify navigation and grasping, and motion-planning frameworks such as MoveIt \cite{Sucan2010MoveIt} and NVIDIA cuMotion provide inverse kinematics and trajectory optimization. For shared workspaces, collaborative robots from Universal Robots\footnote{\url{https://www.universal-robots.com/}} and AMRs from Locus Robotics\footnote{\url{https://locusrobotics.com/}} provide the safety-rated platforms on which the HRC layer is built.






\subsection{Challenges}
\label{subsec:p4-sec9-5-challenges}
Deployed warehouse systems work because they fix the SKU sets, layouts, and operating conditions they handle, and each stage of the pipeline has its own obstacle to relaxing those limits: small-object and cluttered perception, perceptual aliasing in repetitive aisles, multi-agent scalability under uncertainty, manipulation of variable objects, and the safety--throughput trade-off in human--robot collaboration.

\paragraph{Cluttered and Small-Object Perception.}
Barcode labels and compact SKU packages often occupy only a few pixels in high-resolution images, leading to frequent misdetections, while occlusion and clutter compound the difficulty as items overlap within bins and shelves. 
A further persistent issue is domain shift: models trained on synthetic or generic data degrade under real-world variation in lighting, reflective surfaces, and packaging. 
Research responses include domain adaptation and self-supervised learning for generalization \cite{Tzeng2017DomainAdversarial}, active perception that deliberately adjusts viewpoints or sensing to resolve occlusion \cite{Bajcsy2018ActivePerception}, and continual learning that lets perception adapt over time without exhaustive retraining, increasingly combined with multi-sensor fusion and neuro-symbolic representations for higher-level causal and semantic reasoning.

\paragraph{Perceptual Aliasing in Repetitive Environments.}
Localization and mapping are degraded by dynamic obstacles (humans, forklifts) that generate transient or misleading landmarks corrupting loop-closure hypotheses \cite{Cadena2016SLAMSurvey}, and above all by perceptual aliasing: repetitive shelving and low-texture surfaces produce ambiguous place recognition and false data associations in feature-based SLAM \cite{Campos2021ORBSLAM3,Labbe2019RTABMap}. 
High-speed AGV and AMR operation further demands low-latency, sub-centimeter pose updates that sensor delays, occlusion, and monocular scale drift can readily violate. Robust multi-modal fusion, learning-based estimation, and tighter integration of SLAM with semantic perception and digital twins are the leading mitigations.

\paragraph{Scalability and Uncertainty in Multi-agent Planning.}
The combinatorial complexity of multi-robot coordination grows exponentially with the number of agents and tasks, rendering exact optimization intractable for large fleets, while real-world uncertainty (new order arrivals, blocked aisles, unpredictable human movement) necessitates adaptive re-planning and decentralized decision making under partial observability \cite{Veloso2015CMUCoBot,Sartoretti2019PRIMAL}. 
Hierarchical and hybrid planners, graph-neural and differentiable coordination, and simulation-informed multi-objective optimization are the active directions toward scalable, resilient fleet behavior.

\paragraph{Manipulation of Variable and Deformable Objects.}
Many SKUs exhibit deformable, reflective, or transparent surfaces that interfere with perception and grasp stability \cite{Calli2015YCB,Mahler2017DexNet}, and deformable handling remains especially difficult because grasp dynamics and contact points resist rigid-body modeling \cite{Levine2016DeepGrasp}. 
Tactile sensing and force control improve slip detection and grasp robustness but face cost and calibration barriers in industrial deployment, motivating hybrid control that combines analytical models with data-driven adaptation for generalization across geometries and grasp contexts.

\paragraph{Safety--Throughput Trade-off in Collaboration.}
Robust human detection under occlusion, variable lighting, and cluttered aisles continues to hinder reliable intent estimation \cite{dragan2013legibility,Mainprice2016Intent}, and the trade-off between throughput and safety remains unresolved: conservative margins reduce efficiency while aggressive planning risks unsafe interaction \cite{Lasota2015HRC}. 
Emerging work emphasizes predictive and explainable safety models that justify decisions in human-interpretable terms, and adaptive safety frameworks that fuse learned intent models with certified control policies for dynamic risk assessment and shared autonomy on the warehouse floor.

\newpage
\chapter{Spatial Intelligence for Virtual Realities}

\FramedInitial{V}{irtual} realities are the media in which spatial intelligence becomes constitutive rather than optional: an extended-reality (XR) device cannot draw a single pixel until it knows where it is, what surrounds it, and where the user is looking and reaching. This makes XR the most demanding real-time deployment target for the perception, reasoning, and generation capabilities of the previous part, specialized to run egocentrically, on-device, and within a hard latency budget.

Beneath everything in XR sits the platform itself, the runtime that localizes the device, maps and registers the surrounding world, senses the user's body, and composites the result in real time; layered on top is the application that has driven the medium hardest, entertainment, spanning live AR and VR experiences and virtual production as well as the offline, high-fidelity reconstruction and performance capture that supply them with authored content.



\section{AR/VR/XR}
\label{sec:p4_s10_1}

\subsection{Introduction}
\label{subsec:p4-sec10-1-introduction}

Augmented, virtual, and mixed reality (collectively, extended reality or XR) form a continuum of media in which a head-worn display couples the user's egocentric viewpoint to a computational world that is registered to, or substituted for, the physical surroundings. 
What distinguishes XR from screen-based media is that spatial intelligence is not an optional feature but a constitutive requirement: before any pixel can be drawn, the system must know where the headset is, what geometry surrounds it, and where the user is looking and reaching. 
XR is therefore the most demanding real-time deployment target for the spatial capabilities surveyed in Part~\ref{part:3}. 
The perception machinery (detection, depth, pose, reconstruction) of Chapter~\ref{chapter:3:sec:6}, the spatial reasoning and planning of Chapter~\ref{sec:reasoning}, and the neural reconstruction and rendering of Chapter~\ref{sec:generation} all reappear here, but specialized to run egocentrically, on-device, and within a hard latency budget.

This section treats the XR \emph{medium itself} as a spatial-computing platform: the runtime engine that localizes the device, maps and registers the world, senses the user's body, and presents a composited result. 
The neighboring section on entertainment (§\ref{sec:entertainment}) builds \emph{on top of} this engine, treating entertainment as an application domain: its real-time AR/VR experiences and virtual production, together with the offline, high-fidelity reconstruction and performance capture that supply them with authored content; here we focus on the platform that makes such experiences run at interactive rates.
Concretely, a headset must continuously answer three coupled questions, \ie, where am I (6-DoF self-localization), what is around me (geometry and semantics), and how do I fuse virtual content into that reality so that it appears physically present, while simultaneously sensing the user and refreshing the display. 
Early systems answered the localization question with outside-in tracking, using external base stations or cameras (SteamVR Lighthouse, the original Oculus constellation), but the dominant modern paradigm is inside-out tracking, in which visual--inertial SLAM runs entirely on the device, as in Meta Quest, Apple Vision Pro, and Microsoft HoloLens. 
The three pivotal tasks below correspond to the stages of this loop.


\begin{figure}[!t]
    \centering
    \includegraphics[width=0.98\linewidth]{figures/part4/section10/xr_overview.pdf}
    \caption{Spatial intelligence for XR as a closed real-time loop: world-facing sensing, tracking and mapping, and registration (T1), user-facing input (T2), and rendering and display (T3). The latency values are illustrative end-to-end targets, not per-stage budgets.}
    \label{fig:p4_s10_1_overview}
\end{figure}


\subsection{System Overview}
\label{subsec:p4-sec10-1-system-overview}

As sketched in Figure~\ref{fig:p4_s10_1_overview}, an XR device closes a tight real-time loop. Egocentric sensors first \textbf{sense} the device and the user (outward cameras and an inertial measurement unit for head pose; inward-facing cameras for eyes and outward- or downward-facing cameras for hands and body). 
The system then \textbf{tracks and maps} the environment, recovering a 6-DoF pose and a persistent geometric model, and \textbf{registers} virtual content into that model with correct anchoring, occlusion, and lighting. 
Finally, it \textbf{renders and displays} the composite back through the headset optics, after which the user moves and the loop repeats. 
The pipeline shares an end-to-end motion-to-photon budget; roughly $20$\,ms at a display refresh of $90$\,Hz or higher are the commonly cited targets, though the right values depend on device and application, and exceeding them degrades presence and comfort. 
The three pivotal tasks map onto this loop directly: Task-1 is the world-facing engine (track, map, register), Task-2 is the user-facing input (egocentric sensing and interaction), and Task-3 is the user-facing output (display and rendering).



\subsection{Pivotal Task-1: Spatial Tracking, Mapping, and Registration}
\label{subsec:p4-sec10-1-pivotal-task-1-spatial-tracking-mapping-and-registration}

The first task is the medium's spatial backbone: it answers where the device is, what surrounds it, and how virtual content is pinned to the world so convincingly that the illusion holds. 
It proceeds in three steps: localize the device, model the environment, and register synthetic content into that model. 
Figure~\ref{fig:p4_s10_1_tracking} sketches this progression from inside-out tracking to a persistent scene map and finally to anchored, occluded, and relit virtual content.

\begin{figure}
    \centering
    \includegraphics[width=0.98\linewidth]{figures/part4/section10/xr_tracking.pdf}
    \caption{The XR world-facing engine: (a) inside-out visual--inertial tracking of 6-DoF pose, (b) mapping of geometry, semantic labels, and persistent anchors, and (c) registration of virtual content with depth-tested occlusion and matched lighting. The room and shading are illustrative, not measured reconstructions.}
    \label{fig:p4_s10_1_tracking}
\end{figure}

\textbf{Inside-out 6-DoF tracking.} The foundational layer estimates the headset's six-degree-of-freedom pose at high rate by fusing visual features with inertial measurements. 
Feature-based visual SLAM with loop closure and relocalization remains the academic standard, exemplified by the ORB-SLAM lineage and its visual-inertial extension ORB-SLAM3~\cite{Campos2021ORBSLAM3}, while tightly-coupled visual--inertial odometry such as VINS-Mono~\cite{qin2018vins} underpins many mobile and headset stacks; direct (photometric) methods such as DSO~\cite{engel2017direct} and the earlier LSD-SLAM~\cite{engel2014lsd} offer an alternative that does not depend on sparse keypoints. 
More recent learned dense systems, including DROID-SLAM~\cite{teed2021droid} and transformer-based geometric matchers~\cite{mast3r}, trade computation for robustness in low-texture and fast-motion regimes. 
These methods are evaluated chiefly on EuRoC~\cite{Burri2016EuRoC} and TUM-VI~\cite{Schubert2018TUMVI}, with absolute trajectory error (ATE) as the headline metric, but the binding constraints in XR are drift accumulation, recovery from tracking loss (the kidnapped-device problem), and the power and thermal envelope of on-device inference.

\textbf{Spatial mapping and scene understanding.} The tracked stream is consolidated into a usable model of the environment: dense depth and mesh, planar surfaces (floor, wall, table), and semantic labels. 
Real-time volumetric fusion in the KinectFusion tradition~\cite{newcombe2011kinectfusion} provides the geometric substrate, on top of which plane and mesh extraction expose the surfaces that gameplay and user interfaces attach to. 
(The offline, high-fidelity reconstruction of such scenes as navigable content is treated in §\ref{sec:entertainment} and Chapter~\ref{sec:generation}; here the map is a means to localization and registration, not an end in itself.) 
Semantic and instance understanding is benchmarked on indoor datasets such as ScanNet~\cite{dai2017scannet} and Matterport3D~\cite{chang2017matterport3d}, with mean intersection-over-union (mIoU) for segmentation, and increasingly on mobile-captured data with laser ground truth such as ARKitScenes~\cite{baruch2021arkitscenes} and high-quality reconstructions such as Replica~\cite{straub2019replica} for depth and geometry evaluation.

\textbf{Persistent and shared anchors.} For content to survive across sessions and to be experienced jointly by co-located users, the map must be persistent and shareable. 
Cloud-hosted and peer-to-peer world maps let a device relocalize against a previously built model, and at city scale a Visual Positioning System anchors content to real landmarks in well-mapped areas. 
The open problems are relocalization under appearance change (lighting, season, rearranged furniture), the memory footprint of on-device maps, and the privacy implications of uploading scans of private spaces, a thread we return to in the challenges.

\textbf{Registration and realistic compositing.} Finally, virtual content must be registered both geometrically and photometrically so that it appears to occupy the scene. 
Geometric registration pins objects to anchors and planes in the SLAM frame; depth-based occlusion lets real geometry correctly hide virtual pixels, as in the per-pixel depth produced for mobile interaction by DepthLab and the ARCore Depth API~\cite{du2020depthlab}; and lighting estimation predicts the dominant light direction and an ambient environment map from a single image so that virtual objects are shaded and shadowed consistently~\cite{gardner2017learning}. 
Small errors here are perceptually unforgiving: a few centimeters of depth error at an object boundary, or a mismatched shadow, immediately reveals the composite as synthetic, which is why robust occlusion and real-time inverse lighting remain active research problems.

\subsection{Pivotal Task-2: Egocentric Human Sensing and Spatial Interaction}
\label{subsec:p4-sec10-1-pivotal-task-2-egocentric-human-sensing-and-spatial-interaction}

The second task closes the loop from the user's body back into the system: it senses the hands, eyes, and posture from inside-out cameras and turns those signals into spatial intent. Low-level egocentric estimates of hands, gaze, and body are fused into selection, manipulation, and locomotion. As Figure~\ref{fig:p4_s10_1_sensing} indicates, the same headset sensors that localize the device are repurposed to track the user, and the resulting estimates drive a gaze-and-gesture interaction model.

\begin{figure}[!ht]
    \centering
    \includegraphics[width=0.98\linewidth]{figures/part4/section10/xr_sensing.pdf}
    \caption{Egocentric human sensing for XR interaction: (a) articulated hand tracking, (b) gaze from inward-facing eye cameras, (c) full-body pose inferred from three tracked points, and (d) gaze-plus-pinch selection, in which gaze supplies the target and a pinch confirms it.}
    \label{fig:p4_s10_1_sensing}
\end{figure}

\textbf{Hand tracking.} Articulated 3D hand pose (commonly $21$ joints, often a full mesh) is regressed from headset cameras under heavy self-occlusion and truncation. 
The standard hand model is MANO~\cite{romero2017mano}, and on-device pipelines such as MediaPipe Hands~\cite{zhang2020mediapipe} combine a palm detector with a landmark regressor to run in real time. 
Methods are benchmarked on FreiHAND~\cite{zimmermann2019freihand} and the two-hand InterHand2.6M~\cite{moon2020interhand2}, with 3D mean per-joint position error (MPJPE) in millimeters as the headline metric; robust two-hand interaction and egocentric truncation remain the hard cases.

\textbf{Eye and gaze tracking.} Estimating 3D gaze from infrared eye images serves three roles: calibration, gaze as a pointing modality, and as the driver for foveated rendering (Task-3). 
Appearance-based estimation in the iTracker and GazeCapture tradition~\cite{krafka2016gazecapture} established that gaze can be regressed from commodity cameras, and headset-specific datasets such as OpenEDS~\cite{garbin2019openeds} provide near-eye infrared imagery with dense eye-region segmentation, with the later OpenEDS2020 track adding gaze prediction (angular error in degrees). 
The decisive interaction consequence is that gaze targets quickly but lacks a natural trigger, a long-standing difficulty known as the Midas-touch problem, which motivates the gaze-plus-pinch pattern discussed below.

\textbf{Full-body pose from sparse signals.} Driving a full-body avatar from a headset and two controllers is an underconstrained inverse problem, since the lower body is unobserved and must be synthesized. 
Recent methods infer plausible full-body motion from these three tracked points, including the transformer-based AvatarPoser~\cite{jiang2022avatarposer} and the diffusion-based AGRoL~\cite{du2023agrol}, while a parallel line recovers pose from sparse body-worn inertial sensors, as in DIP~\cite{huang2018dip}, TransPose~\cite{yi2021transpose}, and the physics-aware PIP~\cite{yi2022pip}. 
These models are trained on large motion corpora such as AMASS~\cite{mahmood2019amass} and evaluated with MPJPE and temporal jitter. 
This sparse, live setting is distinct both from the wearable-IMU body sensing of §\ref{sec:p4_s9_3} (which instruments the body more densely) and from the studio performance capture that drives authored avatars in §\ref{sec:entertainment}; the XR-specific challenge is plausibility from radically incomplete observation.

\textbf{Multimodal interaction and locomotion.} The recovered signals are fused into spatial actions. The foundational idea of combining a deictic gesture with a complementary modality dates to Bolt's ``Put-That-There''~\cite{bolt1980putthatthere}, and the modern instantiation decouples targeting from confirmation: the eyes select while a hand pinch or a voice command commits, sidestepping the Midas-touch problem. 
Locomotion through a virtual space larger than the physical room is handled by redirected walking, which applies imperceptible rotational and translational gains to steer the user within a confined area~\cite{razzaque2001redirected}, with teleportation as the common non-natural fallback. 
Spatial user interfaces, in turn, must choose a reference frame (world-locked, body-locked, or head-locked) and increasingly anchor panels to scene semantics so that content lands on appropriate surfaces.
Learning these interaction behaviors at scale is increasingly grounded in large egocentric corpora such as Ego4D~\cite{grauman2022ego4d}, which supply the first-person video of hands, objects, and activities that egocentric-understanding models are trained and benchmarked on.

\subsection{Pivotal Task-3: Spatial Display and Real-Time Rendering}
\label{subsec:p4-sec10-1-pivotal-task-3-spatial-display-and-real-time-rendering}

The third task presents the composited world back to the user within the perceptual budget that the previous two stages must respect. 
Three questions decide it: where to spend rendering effort, how to hide system latency, and how to display the result comfortably. 
Figure~\ref{fig:p4_s10_1_display} previews this pipeline, from gaze-contingent shading through latency compensation to the optics that place the image before the eye.

\begin{figure}[!ht]
    \centering
    \includegraphics[width=0.98\linewidth]{figures/part4/section10/xr_display.pdf}
    \caption{The XR display and rendering pipeline: (a) gaze-contingent foveated rendering, (b) late reprojection of a rendered frame to the latest pose, and (c) alternative video-passthrough and optical see-through display paths. The motion-to-photon target is illustrative and applies end to end.}
    \label{fig:p4_s10_1_display}
\end{figure}

\textbf{Foveated rendering.} Because acuity falls off sharply outside the fovea, rendering effort can be concentrated at the gaze point. Foveated 3D graphics established multi-resolution shading driven by eye tracking~\cite{guenter2012foveated}, and subsequent work added temporal stability and contrast preservation to make the quality reduction imperceptible~\cite{patney2016foveated}. Eye-tracked foveated rendering now ships on consumer headsets, materially reducing shading cost on thermally constrained hardware.

\textbf{Latency compensation.} Even an efficient renderer cannot meet the motion-to-photon budget on every frame, so the last stage warps the rendered image to the most recent head pose. 
Asynchronous time warp applies a rotational reprojection at display refresh, asynchronous space warp adds depth-aware frame extrapolation, and head-mounted optical-see-through systems use per-pixel late-stage reprojection~\cite{vanwaveren2016timewarp}. 
These techniques keep the world stable under fast head motion and are essential to comfort.

\textbf{Passthrough versus optical see-through.} Two display regimes trade off differently. Video passthrough mixed reality (as on Meta Quest~3 and Apple Vision Pro) renders the real world from cameras, granting full control over occlusion and relighting at the cost of camera-to-display latency, parallax, and limited dynamic range; optical see-through (as on HoloLens~2 and Magic Leap~2) presents the real world with zero latency at its true focal depth (though the virtual overlay still incurs render latency and a fixed focal plane) but cannot occlude real light or match real-world brightness. 
The choice propagates back into the registration problem of Task-1.

\textbf{Depth-correct rendering and visual comfort.} Conventional stereo displays fix the focal plane while vergence varies with content, producing the vergence-accommodation conflict and the eye strain that limits session length. 
Varifocal, light-field, and holographic near-eye displays aim to restore correct focus cues~\cite{maimone2017holographic}, but remain largely at the prototype stage.

\textbf{Neural rendering for telepresence.} Finally, photoreal social presence requires driving a faithful avatar of a remote user from headset-mounted cameras that see only oblique crops of the face. 
The codec-avatar line, beginning with deep appearance models~\cite{lombardi2018deep}, learns a person-specific decoder that is driven on-device, a direct specialization of the neural rendering of Chapter~\ref{sec:generation} to the constrained, real-time XR setting.

\begin{table}[!ht]
\centering
\fontsize{8}{10}\selectfont
\setlength{\tabcolsep}{1.0mm}
\caption{Representative benchmarks for spatial intelligence in XR, spanning the three pivotal tasks.}
\label{tab:p4_s10_1_benchmarks}
\begin{tabular}{L{0.16\textwidth} L{0.25\textwidth} L{0.30\textwidth} L{0.20\textwidth}}
\toprule
\textbf{Pivotal task} & \textbf{Benchmark} & \textbf{Representative method} & \textbf{Headline metric} \\
\midrule
Tracking & EuRoC~\cite{Burri2016EuRoC} & ORB-SLAM3~\cite{Campos2021ORBSLAM3} & ATE \\
Tracking & TUM-VI~\cite{Schubert2018TUMVI} & VINS-Mono~\cite{qin2018vins} & ATE \\
Mapping & ScanNet~\cite{dai2017scannet} & dense fusion / semantics & mIoU \\
Mapping & ARKitScenes~\cite{baruch2021arkitscenes} & mobile RGB-D mapping & depth abs-rel \\
Hand tracking & InterHand2.6M~\cite{moon2020interhand2} & InterNet~\cite{moon2020interhand2} & 3D MPJPE \\
Body (sparse) & AMASS~\cite{mahmood2019amass} & AvatarPoser~\cite{jiang2022avatarposer} & MPJPE / jitter \\
Gaze & OpenEDS~\cite{garbin2019openeds} & appearance-based gaze regressor & angular error \\
Interaction & Ego4D~\cite{grauman2022ego4d} & egocentric understanding & task-specific \\
\bottomrule
\end{tabular}
\end{table}

\subsection{Established Tools and Systems}
\label{subsec:p4-sec10-1-established-tools-and-systems}

A mature ecosystem implements these tasks across mobile AR, standalone headsets, and the web, and it is organized into developer frameworks, cross-platform standards, hardware, and the companies that drive them.

\textbf{SDKs and standards.} At the device level, \textbf{Apple's ARKit}\footnote{\url{https://developer.apple.com/augmented-reality/}} and \textbf{Google's ARCore}\footnote{\url{https://developers.google.com/ar}} expose plane detection, anchors, depth, light estimation, and people occlusion, with ARCore's Geospatial API localizing against world imagery for outdoor anchoring; \textbf{Vuforia}\footnote{\url{https://developer.vuforia.com}} specializes in robust image and object targets for industrial AR. 
Cross-engine abstraction layers such as \textbf{Unity AR Foundation}\footnote{\url{https://docs.unity3d.com/Packages/com.unity.xr.arfoundation@6.0/}} and the \textbf{Meta Presence Platform and Spatial SDK}\footnote{\url{https://developers.meta.com/horizon/develop/spatial-sdk}} present the common subset (planes, anchors, meshes, occlusion, hand and body tracking) over multiple devices, while \textbf{8th Wall and WebXR}\footnote{\url{https://immersiveweb.dev}} bring markerless AR into the browser. These abstractions matter because XR content is expensive to author and must migrate across devices and sessions; this is the role of cross-platform standards, principally \textbf{OpenXR}\footnote{\url{https://www.khronos.org/openxr/}} as a unifying runtime API and glTF together with OpenUSD and \textbf{USDZ}\footnote{\url{https://openusd.org}} as asset and scene-description formats. Without a shared runtime and shared asset formats, the ecosystem fragments into incompatible walled gardens.

\textbf{Hardware.} Current headsets divide along the passthrough versus optical-see-through axis. 
\textbf{Apple Vision Pro}\footnote{\url{https://www.apple.com/apple-vision-pro/}} (2024) and \textbf{Meta Quest~3 and 3S}\footnote{\url{https://www.meta.com/quest/}} use high-resolution video passthrough, with the Quest line the consumer volume leader; \textbf{Microsoft HoloLens~2}\footnote{\url{https://www.microsoft.com/en-us/hololens/}} (optical see-through, with production ended in 2024) and \textbf{Magic Leap~2}\footnote{\url{https://www.magicleap.com}} target enterprise, and \textbf{Varjo}\footnote{\url{https://varjo.com}} and \textbf{Pico}\footnote{\url{https://www.picoxr.com}} serve high-fidelity simulation and the Asian consumer market respectively. 
Lightweight true-AR glasses, such as \textbf{Snap Spectacles}\footnote{\url{https://www.spectacles.com}} and \textbf{Meta's Orion and Project Aria} research prototypes, signal the all-day-glasses horizon toward which the field is converging.

\textbf{Companies and platforms.} Meta leads consumer volume and the most aggressive glasses research bet; 
\textbf{Apple} defines the premium spatial-computing tier and the USD content pipeline; \textbf{Google} supplies world-scale localization through ARCore and Android XR; \textbf{Microsoft} has pivoted from first-party HoloLens hardware toward OpenXR software and enterprise cloud services; \textbf{Niantic}\footnote{\url{https://www.nianticspatial.com}} operates a city-scale Visual Positioning System for persistent, multi-user AR; and Snap, Varjo, and ByteDance's Pico round out the social-AR, high-fidelity, and consumer-standalone niches.

\subsection{Challenges}
\label{subsec:p4-sec10-1-challenges}

XR runs every spatial capability egocentrically and on-device under a perceptual real-time budget, and its open problems follow from that constraint: holding the motion-to-photon budget, fitting perception into the compute and thermal envelope of glasses, keeping tracking and registration robust in the wild, sharing persistent maps without compromising privacy, and making extended wear visually comfortable.

\paragraph{Latency and the Motion-to-Photon Budget.} Meeting the motion-to-photon targets of the System Overview is increasingly difficult as pipelines add wireless or split rendering, cloud offload, and learned perception in the loop, each of which contributes its own delay that the reprojection stage of Task-3 can only partially mask.

\paragraph{On-Device Compute, Thermal, and Power.} All-day glasses face a hard energy and heat envelope: a lightweight form factor leaves little room for the silicon that SLAM, hand and eye tracking, and scene understanding demand. The frontier is partitioning perception across on-device accelerators, companion devices such as wristbands or pucks, and edge offload, without violating the latency or thermal budget.

\paragraph{Robust Tracking and Registration in the Wild.} Visual--inertial SLAM and relocalization degrade in dynamic scenes, textureless rooms, repetitive structures, low light, and outdoors, exactly the conditions that curated benchmarks such as EuRoC and ScanNet underrepresent~\cite{Burri2016EuRoC,dai2017scannet}. 
Persistent Visual Positioning Systems push toward city-scale anchoring, but globally drift-free, cross-device localization, together with convincing occlusion and lighting registration, remains unsolved.

\paragraph{Persistent, Shared Maps and Privacy.} Co-located multi-user AR requires a shared, persistent spatial map, an ``AR cloud'', which collides directly with privacy: always-on egocentric cameras continuously scan private spaces and bystanders. 
Designing maps that support shared anchoring while minimizing the retention of raw imagery, and that admit meaningful consent, is both a research and a policy frontier.

\paragraph{Visual Comfort.} Beyond latency, fixed-focus stereo optics decouple vergence from accommodation, causing eye strain that limits session length, and poorly designed locomotion induces cybersickness. 
Varifocal and holographic displays, gaze-contingent depth-of-field, and accurate interpupillary-distance matching are promising but immature, leaving comfort a gating factor for mainstream adoption.

Because XR subjects the perception, reasoning, and generation capabilities of Part~\ref{part:3} to egocentric, on-device, real-time constraints, it exposes failure modes that offline settings rarely reveal.

\newpage
\section{Entertainment}
\label{sec:entertainment}

\subsection{Introduction}
\label{subsec:p4-sec10-2-introduction}

In entertainment, spatial intelligence is exercised chiefly as the ability to perceive and reconstruct 3D environments, objects, and human bodies, and to embed virtual content into the physical world so that it is geometrically and physically consistent. Spatial entertainment (AR games, VR experiences, location-based attractions, and virtual production for film) differs from screen-based media in that it must compute camera pose, scene geometry, and user state continuously and within strict latency limits, without sacrificing aesthetic quality or narrative coherence; the perception and reconstruction machinery involved is treated in Chapter~\ref{chapter:3:sec:6} and Chapter~\ref{sec:generation}.

The field has followed the same arc as the capabilities it borrows: commodity depth sensing made real-time dense mapping practical for AR games, neural radiance and Gaussian representations made captured sets photoreal and navigable, parametric body models made full-body avatars routine, and world-scale visual positioning let location-based games such as Pok\'emon GO anchor characters to real landmarks across sessions and players. Each of these threads returns in the pivotal tasks below.

\begin{figure}[!ht]
    \centering
    \includegraphics[width=0.98\linewidth]{figures/part4/section10/entertainment_overview.pdf}
    \caption{Spatial intelligence for entertainment as a real-time closed loop: sensor inputs feed world reconstruction, performance capture, and content authoring, whose composited render is presented back to the user.}
    \label{fig:p4_s10_2_overview}
\end{figure}

\subsection{System Overview}
\label{subsec:p4-sec10-2-system-overview}

Entertainment sits on top of the XR platform of §\ref{sec:p4_s10_1}: it does not re-implement that runtime but builds on it, and its distinctive demand is production-quality spatial content, whether rendered live or authored offline. The domain spans two regimes. Real-time interactive experiences (AR games, VR titles, location-based attractions, and live virtual production) close the tight sensing-to-display loop sketched in Figure~\ref{fig:p4_s10_2_overview}, whereas offline high-fidelity creation (scanning sets, capturing performers, building volumetric footage) supplies those experiences with assets.
Across both, three spatial tasks recur as the frequent bottlenecks.
\textbf{World reconstruction} builds the geometric and radiance model of the environment, supplying localization and occlusion for live experiences and navigable captured content for authored ones;
\textbf{embodied performance capture} digitizes the user or performer into a controllable avatar; and
\textbf{spatial content authoring} decides where virtual elements appear and how cameras move so that the result is both consistent and expressive.
The interactive regime runs all three under hard latency budgets, while offline authoring trades that latency for fidelity.



\subsection{Pivotal Task-1: World Reconstruction and Spatial Mapping}
\label{subsec:p4-sec10-2-pivotal-task-1-world-reconstruction-and-spatial-mapping}
The first task builds and maintains an accurate, consistent model of the environment that supports localization, occlusion, and physics for live experiences and navigable content for authored ones, specializing the dense-reconstruction and scene-understanding techniques of Chapter~\ref{chapter:3:sec:6} and Chapter~\ref{sec:generation} to interactive rates. Where §\ref{sec:p4_s10_1} treated mapping as a means to on-device localization, here the reconstruction is itself the deliverable: the navigable set or stage that players and film-makers inhabit.
Figure~\ref{fig:p4_s10_2_reconstruction} contrasts geometric fusion, neural and point-based scene representations, and world-scale localization as complementary capabilities rather than obligatory sequential stages.

\begin{figure}
    \centering
    \includegraphics[width=0.98\linewidth]{figures/part4/section10/entertainment_reconstruction.pdf}
    \caption{Complementary capabilities for entertainment reconstruction, not sequential stages: (a) TSDF fusion of posed RGB-D frames into a mesh, (b) radiance-field or 3D Gaussian appearance models for view synthesis, and (c) visual positioning for persistent place-based content. Samples, surfaces, and place markers are schematic.}
    \label{fig:p4_s10_2_reconstruction}
\end{figure}

\textbf{Geometric reconstruction} forms the base: geometric SLAM and dense methods fuse depth or RGB data into volumetric maps, as in KinectFusion~\cite{newcombe2011kinectfusion} and the reconstructions underlying datasets such as ScanNet~\cite{dai2017scannet}, while DynamicFusion extends this to moving users by tracking non-rigid deformations \cite{newcombe2015dynamicfusion}. 
\textbf{Neural and point-based representations} such as NeRF, instant-NGP, and 3D Gaussian Splatting then represent scene radiance to produce compact, high-fidelity reconstructions that render at interactive rates \cite{mildenhall2021nerf,muller2022instant,kerbl20233d}; dynamic and 4D extensions that deform or time-index the primitives carry this to moving scenes, so a performance or an entire set can be captured as free-viewpoint volumetric video for later replay \cite{yang2024deformable,wu20244d}. At \textbf{global scale}, visual positioning systems (VPS) combine local reconstructions into a worldwide map that lets a mobile device localize from a single image, as in Niantic's Lightship platform. These methods are evaluated on indoor datasets such as ScanNet and Matterport3D, which provide large-scale reconstructions with semantic labels for scene understanding \cite{dai2017scannet,chang2017matterport3d}, and on NeRF-style view-synthesis benchmarks (\eg, synthetic Blender scenes) for photometric fidelity; the latter, however, rarely capture entertainment-specific constraints such as rapid camera motion or cluttered living rooms, so practitioners often fall back on custom internal datasets (Table~\ref{tab:p4_s10_2_benchmarks}). 
The standard metrics are absolute trajectory error (ATE) for localization, surface reconstruction error for geometry, and PSNR/SSIM for neural rendering \cite{mildenhall2021nerf,kerbl20233d}, but in entertainment these are complemented by interaction-centric measures: the accuracy of occlusion masks, the physical plausibility of collisions, and user reports on realism (for instance, whether a virtual character convincingly hides behind a sofa).

\subsection{Pivotal Task-2: Embodied Performance Capture and Avatar Control}
\label{subsec:p4-sec10-2-pivotal-task-2-embodied-performance-capture-and-avatar-control}
The second task estimates full-body pose, shape, and facial expression to drive avatars in real time, specializing the human pose and motion machinery of Chapter~\ref{chapter:3:sec:6} (and the motion-generation methods of Chapter~\ref{sec:generation}) to live interaction. Unlike the sparse three-point body tracking of §\ref{sec:p4_s10_1}, driven from a headset and two controllers, entertainment capture targets studio-grade, fully-observed performance.
As Figure~\ref{fig:p4_s10_2_capture} indicates, approaches range from skeleton-based joint prediction to surface-based model fitting and full volumetric capture.

\begin{figure}
    \centering
    \includegraphics[width=0.98\linewidth]{figures/part4/section10/entertainment_capture.pdf}
    \caption{Three routes to performance capture, each able to drive avatar playback or control: (a) skeleton joints retargeted to a rig, (b) a parametric surface (SMPL) driving a body mesh, and (c) volumetric capture for reconstructed or animatable appearance. Performer illustrations and motion traces are qualitative, not measured scores.}
    \label{fig:p4_s10_2_capture}
\end{figure}

\textbf{Skeleton-based} methods predict joint positions from sensors, while \textbf{surface-based} approaches fit parametric models such as SMPL directly to video \cite{loper2015smpl}, with datasets like Human3.6M and AMASS supplying the supervised motion priors that make this robust \cite{IonescuPOS14,mahmood2019amass}. 
\textbf{Volumetric capture} systems such as Holoportation go further, streaming textured 3D avatars into VR/AR scenes by combining reconstruction with tight time synchronization \cite{orts2016holoportation}. More recent Gaussian-splatting avatars learn drivable, photorealistic humans from multi-view capture, with rigged 3D Gaussians giving expressive control over head avatars \cite{qian2024gaussianavatars} and pose-dependent Gaussian maps modeling full-body avatars with dynamic clothing \cite{li2024animatablegaussians}, narrowing the gap between a captured performance and a real-time renderable asset. These capabilities are benchmarked chiefly on Human3.6M, which offers millions of 3D poses with calibrated cameras \cite{IonescuPOS14}, and on AMASS, which provides over forty hours of motion data for learning priors \cite{mahmood2019amass}; commercial scenarios remain underrepresented, however, since common cases such as players using controllers in small rooms or performers in costumes lack sufficient public benchmark data (Table~\ref{tab:p4_s10_2_benchmarks}).
Quantitative evaluation uses metrics like mean per-joint position error (MPJPE), but for interactive entertainment latency and temporal stability are decisive, as even minor jitter breaks the sense of embodiment, and user studies show that enjoyment depends on the alignment between perceived and actual motion \cite{dongas2023vrrhythm}; production systems therefore pair algorithmic metrics with feedback from expert animators.

\subsection{Pivotal Task-3: Spatially-Aware Content Authoring and Virtual Production}
\label{subsec:p4-sec10-2-pivotal-task-3-spatially-aware-content-authoring-and-virtual-production}
The third task helps creators decide where virtual elements appear and how cameras move, turning the reconstructed world and captured performers into a coherent, expressive scene, and connecting to the spatial-reasoning and planning paradigms of Chapter~\ref{sec:reasoning}. 
Figure~\ref{fig:p4_s10_2_authoring} previews the spectrum from SDK-level authoring primitives to camera-tracked LED-wall virtual production and AI-assisted cinematography.

\begin{figure}
    \centering
    \includegraphics[width=0.98\linewidth]{figures/part4/section10/entertainment_authoring.pdf}
    \caption{Three related spatial authoring capabilities: (a) engine primitives (detected planes and persistent anchors) for asset placement, (b) virtual production, where tracked camera pose updates the LED-wall perspective, and (c) camera planning from narrative and scene constraints to a shot plan. The scene, camera path, and shot settings are illustrative.}
    \label{fig:p4_s10_2_authoring}
\end{figure}

\textbf{Engine-level authoring} provides designers with high-level primitives such as plane detection and anchors through AR SDKs in real-time engines like Unity and Unreal.
\textbf{Virtual production} links physical cameras to 3D scenes on LED walls so that lighting and parallax respond correctly to camera movement, and \textbf{intelligent cinematography} research is exploring AI-assisted camera planning that respects film grammar while adapting to dynamic user actions \cite{azzarelli2025intelligentcine}. Standardized benchmarks for spatial authoring are rare: some studies evaluate camera planning on synthetic storyboards, but most assessment happens internally at studios, and virtual-production workflows are typically tested on complex reference sequences rather than public datasets, which makes objective comparison of authoring tools difficult (Table~\ref{tab:p4_s10_2_benchmarks}). 
Evaluation accordingly mixes system metrics such as rendering latency with expert judgment on cinematic quality, with intelligent-cinematography studies often convening expert panels to analyze framing and continuity \cite{azzarelli2025intelligentcine}; for AR/VR tools usability is paramount, and HCI evaluations measure how quickly a user completes a level-design task or how rich the resulting spatial layouts are under different tool abstractions.
Generative world models are beginning to bypass authoring by synthesizing playable worlds directly: GameNGen~\cite{valevski2024diffusion} simulates the game DOOM interactively in real time with a diffusion model, and Genie~3~\cite{genie3} generates navigable environments that respond to user actions, so that creating content starts to mean conditioning a world model (§\ref{sec:p3_s8_6}).

\begin{table}[!th]
\centering
\fontsize{8}{10}\selectfont
\setlength{\tabcolsep}{1.0mm}
\caption{Representative benchmarks for spatial intelligence in entertainment, spanning the three pivotal tasks.}
\label{tab:p4_s10_2_benchmarks}
\begin{tabular}{L{0.19\textwidth} L{0.22\textwidth} L{0.30\textwidth} L{0.20\textwidth}}
\toprule
\textbf{Pivotal task} & \textbf{Benchmark} & \textbf{Representative method} & \textbf{Headline metric} \\
\midrule
World reconstruction & ScanNet~\cite{dai2017scannet} & Dense TSDF fusion~\cite{newcombe2011kinectfusion} & recon.\ error \\
World reconstruction & Matterport3D~\cite{chang2017matterport3d} & Dense SLAM & recon.\ error / mIoU \\
View synthesis & NeRF-synthetic~\cite{mildenhall2021nerf} & 3D Gaussian Splatting~\cite{kerbl20233d} & PSNR / SSIM \\
Performance capture & Human3.6M~\cite{IonescuPOS14} & SMPL fitting~\cite{loper2015smpl} & MPJPE \\
Performance capture & AMASS~\cite{mahmood2019amass} & Motion-prior learning & MPJPE / jitter \\
Content authoring & (no public benchmark) & Intelligent cinematography~\cite{azzarelli2025intelligentcine} & expert panel / latency \\
\bottomrule
\end{tabular}
\end{table}

\subsection{Established Tools and Systems}
\label{subsec:p4-sec10-2-established-tools-and-systems}

A robust set of tools implements these tasks for games, film, and live events. At the engine level, Unity's AR Foundation\footnote{\url{https://docs.unity3d.com/Packages/com.unity.xr.arfoundation}} offers a unified API over ARKit, ARCore, and other SDKs, exposing essential features such as point clouds and depth meshing, while Google's ARCore Geospatial API\footnote{\url{https://developers.google.com/ar/develop/unity-arf/geospatial/developer-guide}} adds global localization, allowing applications to attach anchors to outdoor landmarks with meter-level accuracy.
These abstractions let designers build AR gameplay around real-world surfaces without managing low-level sensor fusion.

Niantic's Lightship platform packages spatial mapping and multiplayer synchronization into turnkey services: Reality Blending\footnote{\url{https://nianticlabs.com/news/realityblending-announcement}} in Pok\'emon GO uses learned depth to let virtual creatures move behind real-world obstacles, and Pok\'emon Playgrounds\footnote{\url{https://nianticlabs.com/news/pokemon-playgrounds}} extends this to persistent AR, where players leave virtual items at specific locations for others to find, with Niantic's Visual Positioning System (VPS) keeping that content anchored consistently across users and times.

In VR gaming, consumer headsets with inside-out tracking enable room-scale interaction; games like Beat Saber use precise 6-DoF tracking to create rhythm gameplay that doubles as exercise, and studies show a strong link between spatial design choices and player exertion, underscoring the connection between tracking quality and user experience \cite{dongas2023vrrhythm}. 
Virtual production stages have likewise transformed film and sports broadcasting: productions such as \emph{The Mandalorian} (built on ILM's StageCraft) and sports studios such as ESPN's Catalyst Stage\footnote{\url{https://espnpressroom.com/us/press-releases/2023/06/catalyst-stage-disneys-new-immersive-espn-studio-opens-storytelling-and-production-frontiers/}} use LED volumes driven by tracked cameras, where spatial intelligence involves calibrating physical cameras to digital scenes and maintaining low-latency tracking. 
Finally, neural capture tools inspired by NeRF and 3D Gaussian Splatting are entering the pipeline, letting artists scan a location with a phone and import navigable 3D assets into game engines \cite{muller2022instant,kerbl20233d}; while this bridges scouting and production, it introduces new challenges around lighting control and integration.

\subsection{Challenges}
\label{subsec:p4-sec10-2-challenges}

Entertainment asks for production-quality spatial content at interactive rates, and its open problems stem from that demand. The tension between real-time performance and visual fidelity is the most immediate; close behind come robustness once the system leaves controlled spaces, the cognitive load that authoring places on creators, and the human-factors and ethics questions that surface the moment these systems scan real users and real rooms.

\paragraph{Real-time Performance versus Fidelity.} Neural representations have improved visual quality, but many methods remain too computationally heavy for mobile or standalone VR. While techniques such as 3D Gaussian Splatting achieve real-time rendering on high-end hardware \cite{kerbl20233d}, porting them to thermally constrained devices is difficult, and entertainment applications must additionally handle variable network conditions while guaranteeing stable frame rates, since frame drops that are acceptable in research benchmarks can cause physical discomfort in VR.

\paragraph{Robustness in the Wild.} Existing benchmarks such as ScanNet focus on standard indoor spaces, yet entertainment experiences often occur in unpredictable environments such as festivals, dim clubs, or crowded stadiums. World-scale AR platforms address this by training on millions of crowdsourced scans, which raises questions about coverage bias (particularly toward urban areas) and the difficulty of maintaining maps as the world changes; designing models that remain robust under these shifts is an open research area.

\paragraph{Authoring Complexity.} Although modern SDKs are powerful, creating high-quality spatial experiences imposes a high cognitive load, as creators must simultaneously manage geometry, interaction, cinematography, and technical constraints. Intelligent-cinematography tools attempt to assist by suggesting camera paths \cite{azzarelli2025intelligentcine}, but these are not yet standard in creative workflows; there is a pressing need for tools that let designers specify high-level intentions (for example, ``the character should hide behind the sofa'') which the system then realizes via underlying planning modules.

\paragraph{Human Factors and Ethics.} Spatial systems frequently scan users and private spaces. While datasets like AMASS demonstrate the utility of large-scale motion capture, collecting similar data in the wild raises privacy concerns, and from an HCI perspective issues such as cybersickness and fatigue limit the duration of VR and AR sessions, while physically active titles introduce comfort and safety considerations that spatial layout must respect. Addressing these problems requires a collaborative approach spanning algorithm design, user-feedback systems, and clear data governance.

\newpage
\chapter{Spatial Intelligence for Geo-spatial Systems }


\FramedInitial{G}{eo-spatial} systems operate on spatial structure at the scale of assets, cities, and the planet, where the task is to build a faithful model of a physical system and keep it aligned with live observation. This chapter applies the spatial intelligence of the previous part to that task, across three settings that share the same closed loop between model and world but differ in scale and stakes.

The digital twin binds a metric, physically grounded replica to a specific asset through a live bidirectional loop, so that the physical system and its model track each other. GeoAI widens the frame to the planet, where streams of Earth-observation data drive perception and decision-making over terrain, infrastructure, and environmental change. Aerospace systems carry the same loop off the planet, where a spacecraft or rover must perceive, reason about, and act on its surroundings with little or no ground support.


\section{Digital Twins}
\label{sec:p4_s11_1}
\label{sec:Digital Twins}

\subsection{Introduction}
\label{subsec:p4-sec11-1-introduction}

A \textbf{digital twin} is a 3D representation faithful enough to stand in for its physical counterpart in analysis and operation.
Building one requires more than accurate geometric reconstruction and photorealistic appearance under varying illumination; it requires realistic modeling of physical properties and interaction, so that the replica behaves as the original would.
What separates a digital twin from any other 3D asset, however, is not fidelity alone but \emph{coupling}.
Following the now-standard taxonomy of Kritzinger \etal~\cite{kritzinger2018digital}, a \emph{digital model} carries no automated data exchange with the physical object, a \emph{digital shadow} carries an automated one-way flow from the physical to the virtual, and only a \emph{digital twin} closes the loop with an automated, bidirectional flow: the virtual model is continuously synchronized to the physical asset, simulated forward to predict, and used to act back upon it.

A digital twin therefore reuses, but specializes, much of the machinery surveyed earlier in this handbook. The perception stack of Chapter~\ref{chapter:3:sec:6} supplies the captures from which a twin is built, and the neural reconstruction and generation methods of Chapter~\ref{sec:generation} supply its geometry and appearance; what the twin adds is metric accuracy, physical grounding, and a persistent correspondence to a specific real object over its lifecycle.
Entertainment (§\ref{sec:entertainment}) builds the same high-fidelity captures as creative \emph{content}; here they become an \emph{operational engineering model} wired to a live counterpart, and the city- and planet-scale twins that recur below connect to the GeoAI methods of §\ref{sec:p4_s11_2}.
The concept itself traces from Grieves' product-lifecycle lectures, formalized with Vickers~\cite{grieves2017digital}, through its aerospace codification for NASA and U.S. Air Force vehicles~\cite{glaessgen2012digital}, to its industrial maturation~\cite{tao2019digital,wang2023survey}, where the five-dimension model of Tao \etal makes the data connections between physical entity, virtual entity, and services a first-class element of the architecture.

\subsection{System Overview}
\label{subsec:p4-sec11-1-system-overview}

Figure~\ref{fig:p4_s11_1_overview} draws the twin as a closed loop around its physical asset. Sensing (images, LiDAR, RGB-D scans, onboard IoT telemetry) feeds the construction of a replica that carries geometry, appearance, materials, and physical properties; the replica is simulated forward, kept synchronized by live data assimilation, and queried to predict, and its decisions operate the asset, which is then sensed anew.
The three pivotal tasks map onto this loop directly. Task-1 builds the twin (sense and construct), Task-2 makes it behave (model and simulate), and Task-3 keeps it synchronized and turns it to use (assimilate, predict, and act). Table~\ref{tab:p4_s11_1_benchmarks} collects the benchmarks on which the first two are evaluated; synchronization has no shared benchmark.

\begin{figure}
    \centering
    \includegraphics[width=0.98\linewidth]{figures/part4/section11/digital_twins_overview.pdf}
    \caption{The digital twin as a closed loop around its physical asset: construction (T1) and modeling (T2) establish the replica, synchronization (T3) corrects its state from live observations, and forecasts drive decisions through a control path separate from the observation path.}
    \label{fig:p4_s11_1_overview}
\end{figure}

\subsection{Pivotal Task-1: High-Fidelity Twin Construction}
\label{subsec:p4-sec11-1-pivotal-task-1-high-fidelity-twin-construction}

The first task builds the twin from sensor captures, and it must deliver four properties that a generic one-shot reconstruction does not guarantee. The twin must be \emph{metric} (correct at real-world scale), \emph{relightable} (its appearance separable from capture-time illumination), \emph{physically grounded} (carrying mass, material, and friction), and \emph{actionable} (its movable parts and joints encoded so a simulator can drive it).
Figure~\ref{fig:p4_s11_1_construction} sketches the progression from raw captures through geometry and material recovery to a simulatable model.

\begin{figure}
    \centering
    \includegraphics[width=0.98\linewidth]{figures/part4/section11/digital_twins_reconstruction.pdf}
    \caption{Three components of high-fidelity twin construction: (a) metric geometry from multi-view, LiDAR, or RGB-D capture; (b) relightable material channels separated from illumination; and (c) physical articulation, here a prismatic drawer joint with mass and friction. Material swatches and mesh overlays are schematic.}
    \label{fig:p4_s11_1_construction}
\end{figure}

\textbf{Geometry and appearance acquisition.} The earliest replicas were built by photogrammetry: multiple overlapping photographs of an object, terrain, or structure are captured, and a Structure-from-Motion algorithm recovers camera poses and a sparse metric point cloud, as systematized in COLMAP~\cite{schonberger2016structure}, after which dense multi-view stereo and meshing produce a surface.
Real-time RGB-D fusion in the KinectFusion tradition~\cite{newcombe2011kinectfusion} yields metric, watertight geometry over the observed region, and neural implicit-surface methods such as NeuS~\cite{wang2021neus} and the detail-preserving Neuralangelo~\cite{li2023neuralangelo} recover watertight meshes as the zero-level-set of a learned signed-distance field.
A second family optimizes for view synthesis rather than surface quality: NeRF~\cite{mildenhall2021nerf}, its hash-grid acceleration Instant-NGP~\cite{muller2022instant}, and the explicit, real-time 3D Gaussian Splatting~\cite{kerbl20233d} render novel views with high fidelity but historically produced no clean surface and bake appearance into the representation; surface-aligned variants such as SuGaR~\cite{guedon2024sugar} and 2D Gaussian Splatting~\cite{huang20242d} narrow this gap by extracting watertight meshes directly from a Gaussian field.
Because a twin requires metric scale and an extractable surface, the photogrammetric and implicit-surface lines are load-bearing, with meshing operators such as Poisson surface reconstruction~\cite{kazhdan2006poisson} and Marching Cubes~\cite{lorensen1998marching} converting a field or point set into geometry; the radiance-field family enters as an appearance and speed substrate, treated in full in Chapter~\ref{sec:generation}.

\textbf{Inverse rendering and relightable materials.} A captured asset whose appearance is baked at scan time only reproduces the object under the one illumination it was scanned in, because each pixel entangles albedo, geometry, and lighting.
A twin, by contrast, lives in a simulator where viewpoint and illumination change every frame, so its construction must recover separated physically based rendering (PBR) channels (albedo, roughness, normals) together with an estimate of the environment lighting.
Inverse-rendering methods supply exactly this factorization: NeRFactor decomposes a captured radiance field into shape, reflectance, and a single environment light~\cite{zhang2021nerfactor}, while the differentiable, engine-ready NVDiffRec jointly extracts a triangle mesh, an SVBRDF, and a lighting estimate directly optimizable for downstream rendering~\cite{munkberg2022extracting}.
The same factorization has been ported onto explicit Gaussians, as in Relightable 3D Gaussians, which attaches normals and a BRDF to each primitive and ray-traces shadows for real-time relighting~\cite{gao2024relightable}. Relightability is what lets a twin be composited under novel lighting and rendered consistently inside a physics or sensor simulation.

\textbf{Physical-property and articulation estimation.} A static mesh is viewable but not operable: it has no notion of which surfaces move, with what effort, or under what dynamics. Closing this gap requires a movable-part inventory and the kinematic joints (revolute or prismatic, with axes and limits) that a physics engine can integrate, together with physical parameters such as mass and friction. Ditto recovers part geometry and an articulation model jointly from before-and-after observations of an interaction~\cite{jiang2022ditto}, and URDFormer predicts a full URDF, the standard interactive description, from a single real image, turning a photograph into a simulatable articulated twin~\cite{chen2024urdformer}. Dynamical parameters can be identified directly in a differentiable field representation, as in PAC-NeRF, which estimates geometry and physical parameters together through a differentiable material-point simulator~\cite{li2023pacnerf}, and PhysGaussian, which integrates Newtonian dynamics into a Gaussian representation so that what is reconstructed can also be simulated~\cite{xie2024physgaussian}. Such methods lean heavily on articulated-object datasets, principally the PartNet-Mobility models exposed through the SAPIEN simulator~\cite{xiang2020sapien}, the fine-grained part hierarchy of PartNet~\cite{mo2019partnet}, and the cross-category, actionable parts of GAPartNet~\cite{geng2023gapartnet}. The resulting articulated, parameterized model is the input to the behavioral simulation of Task-2.

\textbf{Datasets and metrics.} Construction quality is measured against scanned-object and scene corpora (Table~\ref{tab:p4_s11_1_benchmarks}).
At the object level, the recent Digital Twin Catalog provides a large set of digital-twin-quality scans with millimeter geometry, photorealistic materials, and an explicit relighting evaluation~\cite{dong2025dtc}, complementing the real-scanned OmniObject3D~\cite{wu2023omniobject3d}, the simulation-ready Google Scanned Objects~\cite{downs2022googleGSO}, the artist-authored, PBR-annotated ABO~\cite{collins2022abo}, the large Objaverse collection~\cite{deitke2023objaverse}, and the in-the-wild Common Objects in 3D~\cite{reizenstein2021common}.
At the scene level, twins are evaluated on ScanNet~\cite{dai2017scannet}, Matterport3D~\cite{chang2017matterport3d}, the photorealistic Replica reconstructions~\cite{straub2019replica}, and the commodity-LiDAR ARKitScenes~\cite{baruch2021arkitscenes}, with metric-geometry benchmarks such as Tanks and Temples supplying laser-scan ground truth~\cite{knapitsch2017tanks}.
Geometry is reported with Chamfer distance, F-score at a stated threshold, and normal consistency; appearance with PSNR, SSIM, and LPIPS; and, for twins specifically, an added relighting error that penalizes appearance that fails to transfer across illumination.

\subsection{Pivotal Task-2: Physics-Based Modeling and Behavioral Simulation}
\label{subsec:p4-sec11-1-pivotal-task-2-physics-based-modeling-and-behavioral-simulation}

The second task makes the constructed replica behave.
Geometry is not behavior: a twin must reproduce its asset's dynamics, not merely its shape, so the static model of Task-1 is coupled to a physics, behavioral, or surrogate engine and run forward as an \emph{executable} model whose parameters are fit to the real asset (Figure~\ref{fig:p4_s11_1_simulation}).

\begin{figure}
    \centering
    \includegraphics[width=0.98\linewidth]{figures/part4/section11/digital_twins_simulation.pdf}
    \caption{From a constructed replica to a behaving twin: (a) constructed model, (b) physics engine, and (c) simulated behavior, with system identification fitting the physical parameters $\theta$ to real-asset observations and an optional fast surrogate $f_\phi$. Response traces are illustrative, not measured results.}
    \label{fig:p4_s11_1_simulation}
\end{figure}

\textbf{Physics engines.} The behavioral substrate of a twin is a simulator that integrates its dynamics.
MuJoCo provides accurate rigid-body contact dynamics and is the de facto standard for model-based control~\cite{todorov2012mujoco}; the open-source PyBullet offers a widely used free alternative spanning rigid and soft bodies~\cite{coumans2021}; and SAPIEN specializes in part-level interaction with large articulated objects~\cite{xiang2020sapien}, the natural runtime for the actionable twins of Task-1.
GPU-accelerated platforms couple photorealistic rendering with physics at scale, exemplified by NVIDIA Isaac Sim\footnote{\url{https://developer.nvidia.com/isaac/sim}} on the Omniverse stack and its successor Isaac Lab\footnote{\url{https://developer.nvidia.com/isaac/lab}}, with newer GPU-first simulators such as Genesis\footnote{\url{https://genesis-embodied-ai.github.io/}} and the accelerator-native MuJoCo~MJX pushing throughput further; deformable phenomena such as cloth and fluids are handled by unified particle and material-point (MPM) solvers.

\textbf{Differentiable simulation and system identification.} A twin is only behaviorally faithful if its physical parameters match the real asset, which makes parameter estimation central.
Differentiable simulators expose gradients of simulated outcomes with respect to physical parameters, turning system identification into gradient-based optimization: DiffTaichi compiles differentiable physical simulators from high-level code~\cite{hu2020difftaichi}, Brax provides massively parallel differentiable rigid-body dynamics on accelerators~\cite{freeman2021brax}, and gradSim couples differentiable simulation with differentiable rendering to backpropagate from video pixels all the way to mass, friction, and elasticity~\cite{jatavallabhula2021gradsim}. A twin can thus be calibrated to its counterpart from observation rather than hand-tuned.

\textbf{Surrogate and physics-informed models.} High-fidelity simulation is often too slow to run in the synchronization loop of Task-3, so twins frequently substitute a reduced or learned surrogate.
Physics-informed neural networks embed governing partial differential equations directly in the training loss, solving forward and inverse problems with scarce data~\cite{raissi2019physics}, a strategy situated within the broader programme of physics-informed machine learning~\cite{karniadakis2021physics}.
Reduced-order modeling, identified early as a key enabling technology for digital twins, compresses a high-fidelity simulation to real-time evaluation~\cite{hartmann2018model}, and the hybrid-twin paradigm couples such a physics-based reduced model with a data-driven correction term to absorb unmodeled effects~\cite{chinesta2020virtual}.

\textbf{Real-to-sim for policy learning.} A faithful behavioral twin is also the substrate on which robot policies are trained and validated before deployment, which connects this task to the manipulation methods of §\ref{sec:p4_s9_1} and the reasoning of Chapter~\ref{sec:reasoning}.
The reality gap between simulation and the physical world is bridged by randomizing visual appearance~\cite{tobin2017domain} and dynamics~\cite{peng2018sim} during training, and recent real-to-sim-to-real pipelines build a twin of the workspace from a quick scan and use it to robustify learned policies, as in RialTo~\cite{torne2024reconciling}.
A complementary line argues that an exact twin is not always necessary: ACDC instead generates geometrically and semantically similar ``digital cousins,'' trading exact correspondence for robustness and zero-shot transfer~\cite{dai2024automated}, sharpening by contrast what a true digital twin provides.

\subsection{Pivotal Task-3: Real-Time Synchronization and Predictive Operation}
\label{subsec:p4-sec11-1-pivotal-task-3-real-time-synchronization-and-predictive-operation}

The third task is the one that earns the name \emph{twin}: it maintains the automated bidirectional coupling between model and asset that defines a twin~\cite{kritzinger2018digital}, and turns that coupling to use.
Figure~\ref{fig:p4_s11_1_synchronization} previews the pipeline, from a live sensor stream through data assimilation to prediction and a control signal fed back to the physical asset.

\begin{figure}
    \centering
    \includegraphics[width=0.98\linewidth]{figures/part4/section11/digital_twins_synchronization.pdf}
    \caption{Real-time synchronization and predictive operation: (a) live measurements from the physical asset are (b) assimilated with the model prior into a posterior state of (c) the synchronized twin, which feeds forecasts, what-if, anomaly, and useful-life outputs. The upper feedback path propagates the updated state into the next model prior; the lower path returns decisions or controls to the physical asset. Dashed forecast branches are possible futures, not observed trajectories; all analytical outputs are schematic.}
    \label{fig:p4_s11_1_synchronization}
\end{figure}

\textbf{Data assimilation and state estimation.} Keeping the twin synchronized is a state-estimation problem: the model's latent state must be continually corrected toward the incoming measurements.
The classical recursive estimator is the Kalman filter~\cite{kalman1960new}, and for the nonlinear, high-dimensional systems typical of physical twins the ensemble Kalman filter provides a Monte-Carlo data-assimilation scheme that propagates an ensemble of states to forecast error statistics~\cite{evensen1994sequential}.
At the level of an entire engineered system, a probabilistic graphical model casts data-driven model adaptation as Bayesian state estimation with end-to-end uncertainty quantification~\cite{kapteyn2021probabilistic}, so that predictive twins can be assembled and updated at scale and their predictions carry calibrated confidence rather than point estimates.

\textbf{Predictive operation across domains.} The synchronized twin is queried to anticipate rather than merely to monitor, and the applications follow the asset.
In manufacturing, the twin couples machine geometry to live telemetry for prognostics and health management, estimating remaining useful life and scheduling maintenance before failure~\cite{tao2018digital}. At urban scale, a city digital twin fuses a semantically rich 3D model with sensor and simulation layers to support planning and what-if analysis, as in the Herrenberg study~\cite{dembski2020urban} and national programmes built on CityGML foundations~\cite{soon2017citygml}, a layer that connects directly to the GeoAI methods of §\ref{sec:p4_s11_2}.
In healthcare, patient-specific anatomical twins driven by physiological models support precision diagnosis and treatment planning, with cardiology an early exemplar~\cite{corralacero2020digital} and a growing body of work surveyed across the field~\cite{katsoulakis2024digital}; the medical twin is developed in §\ref{sec:p4_s12_1}.
In aerospace, where the paradigm was first codified, a structural twin mirrors an as-flown airframe to predict fatigue~\cite{glaessgen2012digital}; and in agriculture, field and crop twins fuse geospatial layers with growth models for precision management~\cite{verdouw2021digital}.

\begin{table}[!th]
\centering
\fontsize{8}{10}\selectfont
\setlength{\tabcolsep}{1.0mm}
\caption{Representative benchmarks for twin construction and simulation. Synchronization (Task-3) has no shared benchmark.}
\label{tab:p4_s11_1_benchmarks}
\begin{tabular}{L{0.20\textwidth} L{0.22\textwidth} L{0.30\textwidth} L{0.20\textwidth}}
\toprule
\textbf{Pivotal task} & \textbf{Benchmark / dataset} & \textbf{Representative method} & \textbf{Headline metric} \\
\midrule
Construction (object) & DTC~\cite{dong2025dtc} & neural reconstruction + PBR & geometry + relighting error \\
Construction (object) & OmniObject3D~\cite{wu2023omniobject3d} & multi-view reconstruction & Chamfer / F-score \\
Construction (materials) & ABO~\cite{collins2022abo} & NVDiffRec~\cite{munkberg2022extracting} & material / PSNR \\
Construction (scene) & ScanNet~\cite{dai2017scannet} & KinectFusion-style RGB-D fusion~\cite{newcombe2011kinectfusion} & Chamfer / F-score \\
Construction (metric) & Tanks and Temples~\cite{knapitsch2017tanks} & MVS / neural surfaces & F-score \\
Construction (articulation) & PartNet-Mobility~\cite{xiang2020sapien} & Ditto~\cite{jiang2022ditto} / URDFormer~\cite{chen2024urdformer} & joint / part error \\
Construction (city-scale) & GauU-Scene / MatrixCity & CityGaussianV2~\cite{liu2025citygaussian} & PSNR / geometry F-score \\
Simulation (real-to-sim) & real-world manipulation tasks & RialTo~\cite{torne2024reconciling} / ACDC~\cite{dai2024automated} & policy success rate \\
\bottomrule
\end{tabular}
\end{table}

\subsection{Established Tools and Systems}
\label{subsec:p4-sec11-1-established-tools-and-systems}

\textbf{Capture, authoring, and render runtimes.} The construction stage is served by spatial-capture products such as \textbf{Matterport}\footnote{\url{https://matterport.com/}} for interiors and sites, the mobile scanner \textbf{Polycam}\footnote{\url{https://poly.cam/}}, the photogrammetry pipelines \textbf{RealityScan}\footnote{\url{https://www.realityscan.com/}} (formerly RealityCapture) and Agisoft Metashape\footnote{\url{https://www.agisoft.com/}}, and, as an integrating substrate, NVIDIA Omniverse with \textbf{OpenUSD}\footnote{\url{https://www.nvidia.com/en-us/omniverse/}}, which provides the scene description and real-time rendering on which twins are assembled. Game and render engines such as \textbf{Unreal Engine}\footnote{\url{https://www.unrealengine.com/}} and Unity, the latter packaged for industry as \textbf{Unity Industry}\footnote{\url{https://unity.com/products/unity-industry}}, serve as the interactive runtimes for the resulting models.

\textbf{Industrial twin platforms.} The operational layer is dominated by enterprise platforms that wire CAD geometry to live telemetry. \textbf{Siemens Teamcenter}\footnote{\url{https://plm.sw.siemens.com/en-US/teamcenter/}} drives product twins from a PLM backbone; \textbf{Ansys Twin Builder}\footnote{\url{https://www.ansys.com/products/digital-twin/ansys-twin-builder}} couples reduced-order models to 3D physics; Dassault's 3DEXPERIENCE markets sensor-driven ``virtual twin'' experiences\footnote{\url{https://www.3ds.com/virtual-twin}}; and cloud services including \textbf{PTC ThingWorx}\footnote{\url{https://www.ptc.com/en/products/thingworx}}, \textbf{AWS IoT TwinMaker}\footnote{\url{https://aws.amazon.com/iot-twinmaker/}}, \textbf{Azure Digital Twins}\footnote{\url{https://azure.microsoft.com/en-us/products/digital-twins}}, and \textbf{Bentley iTwin}\footnote{\url{https://www.bentley.com/software/itwin-platform/}} maintain live twin graphs of factories, buildings, and infrastructure, with \textbf{GE Vernova}\footnote{\url{https://www.gevernova.com/software/}} specializing in energy-asset twins.

\textbf{City-scale systems and interoperability standards.} At geo-scale, procedural and streaming platforms such as \textbf{Esri ArcGIS CityEngine}\footnote{\url{https://www.esri.com/en-us/arcgis/products/arcgis-cityengine/overview}}, \textbf{Cesium with 3D Tiles}\footnote{\url{https://cesium.com/}}, \textbf{Google's photorealistic 3D Tiles}\footnote{\url{https://developers.google.com/maps/documentation/tile/3d-tiles}}, and \textbf{NVIDIA Earth-2}\footnote{\url{https://www.nvidia.com/en-us/high-performance-computing/earth-2/}} build and serve city- to planet-scale twins, with national programmes such as Virtual Singapore as flagship deployments. What keeps such twins portable across this fragmented tooling are open standards: OpenUSD for scene description, IFC (ISO 16739) for building information, CityGML for semantic city models, and OGC 3D Tiles for streaming massive geospatial geometry, alongside the framework and terminology standards ISO 23247 and ISO/IEC 30173. The data layer feeding these geo-scale twins is the subject of §\ref{sec:p4_s11_2}.

\subsection{Challenges}
\label{subsec:p4-sec11-1-challenges}

A twin must be high-fidelity, live, and scalable at once, and these three demands pull against one another. The open problems follow from that tension: scaling construction from an object to a city, keeping the twin current, validating that it is accurate rather than merely photorealistic, coupling several physical domains in real time, and making twins interoperable without exposing the spaces they capture.

\paragraph{Scalability to Large and City Scale.} Building a twin of a single object is now routine; building one of a city is not.
City twins must convert heterogeneous source data into semantically structured models (CityGML levels of detail) and stream them efficiently (3D Tiles), and the reconstruction itself must scale, motivating divide-and-conquer pipelines for large-scene Gaussian splatting~\cite{liu2025citygaussian}.
Data integration and cross-format interoperability remain the named bottlenecks at this scale.

\paragraph{Synchronization Latency and Data Assimilation.} A twin is only as useful as it is current, and maintaining real-time correspondence across an edge-fog-cloud hierarchy is hard: assimilation must keep pace with the asset's dynamics, and the ``liveness gap'' between the asset's true state and the twin's estimate degrades every downstream prediction. Scaling data assimilation to high-dimensional physical systems without violating the update budget is an open problem.

\paragraph{Model Fidelity, Validation, and the Photorealism-Accuracy Tradeoff.} Verification, validation, and uncertainty quantification are the maturity frontier for twins used in consequential decisions.
A specific tension recurs: methods such as vanilla Gaussian splatting optimize visual fidelity, not metric accuracy, so a photorealistic reconstruction can still fail a metric tolerance check; surface-aligned variants~\cite{guedon2024sugar,huang20242d} narrow but do not close this gap.

\paragraph{Multi-Physics and Multi-Scale Coupling.} Real assets couple mechanics, thermics, fluids, and electromagnetics across scales, and no universal model spans them; coupled systems risk solver instability, and the reduced-order and physics-informed surrogates~\cite{hartmann2018model,raissi2019physics} that make twins real-time must do so without discarding the cross-domain interactions that matter.

\paragraph{Standardization, Interoperability, Privacy, and Security.} The proliferation of scene and asset formats (USD, IFC, CityGML, 3D Tiles) keeps twins from composing across vendors, and the always-on capture that a live twin requires raises acute concerns: scanning private interiors and urban space implicates bystanders, and a compromised twin server can leak an asset's dynamics, configuration, and credentials~\cite{wang2023survey}. Designing twins that interoperate while minimizing retained raw data is both a research and a policy frontier.

\newpage
\section{GeoAI}
\label{sec:p4_s11_2}


\subsection{Introduction}
\label{subsec:p4-sec11-2-introduction}
Geospatial artificial intelligence (GeoAI) applies machine learning to geographic data (satellite and UAV imagery, LiDAR, GNSS traces, and IoT sensor streams) to model and act on processes that unfold over terrain, cities, and the Earth system.
It makes the geospatial domain a large-scale deployment target for the perception and scene-understanding capabilities surveyed in Part~\ref{part:3}: the detection, segmentation, and remote-sensing techniques of Chapter~\ref{chapter:3:sec:6} (and the aerial and remote-sensing setting of Section~\ref{sec:p4_s9_4}) are here specialized to Earth-observation data and decision workflows at planetary scale.

\subsection{System Overview}
\label{subsec:p4-sec11-2-system-overview}
A GeoAI workflow couples heterogeneous geospatial data to models that turn raw observations into spatially explicit decisions (Figure~\ref{fig:p4_s11_2_overview}).
Following the application taxonomy of the review by Boutayeb \etal~\cite{boutayeb2024geoai}, we treat six domains as pivotal tasks: precision agriculture, the urban built environment, environmental monitoring, water resources, disaster management, and public health.
The six share their inputs, largely the free Sentinel, Landsat, and MODIS archives, and until recently they shared a modeling pattern as well: a task-specific supervised network or tree ensemble trained for one region and one target, which the foundation models at the end of the section are beginning to replace.
Table~\ref{tab:p4_s11_2_benchmarks} samples representative data sources, methods, and reported results across the six.

\begin{figure}[!ht]
    \centering
    \includegraphics[width=0.98\linewidth]{figures/part4/section11/geoai-overview.pdf}
    \caption{GeoAI as a data-to-decision loop: satellite, UAV, LiDAR, GNSS, and IoT observations feed alternative or combined CNN, transformer, and ensemble models, whose outputs serve six application domains and act back, through planning and policy, on the sensed world.}
    \label{fig:p4_s11_2_overview}
\end{figure}

\subsection{Pivotal Task-1: Precision Agriculture}
\label{subsec:p4-sec11-2-pivotal-task-1-precision-agriculture}
Precision agriculture tailors irrigation, fertilization, and pest control to the conditions of each field, or each plant, instead of managing a field uniformly.
GeoAI supplies the maps and predictions this requires through the complementary mapping, yield-estimation, and irrigation tasks of Figure~\ref{fig:p4_s11_2_agriculture}.

\begin{figure}[!ht]
    \centering
    \includegraphics[width=0.98\linewidth]{figures/part4/section11/geoai-agriculture.pdf}
    \caption{Three GeoAI tasks for precision agriculture: (a) crop mapping, (b) yield and leaf-area-index (LAI) estimation, and (c) variable-rate precision irrigation. Maps and fitted points are schematic.}
    \label{fig:p4_s11_2_agriculture}
\end{figure}

\textbf{Crops Mapping.} Crop mapping turns a time series of satellite or UAV imagery into a per-pixel crop-type label, and the accuracy gains in recent work track the richness of the temporal signal.
Han \etal~\cite{han2023spatio} combine a ResNet feature extractor with cross-attention and spatio-temporal self-attention over Sentinel-1 SAR series, reaching 96\% overall accuracy on the Brandenburg benchmark.
3D convolutional networks over Sentinel-2~\cite{gallo2023season} and Landsat~\cite{MOHAMMADI2023272} series treat time as a third convolution axis, and fusing Sentinel-1 SAR with Sentinel-2 optical data maps rice phenological stages within the season~\cite{tiwari2024automated}.
Multi-season and SAR--optical inputs consistently beat single-date classification, in part because SAR sees through the cloud that interrupts optical series.
At the other end of the resolution range, UAV imagery supports row-level segmentation, where U-Net outperforms LinkNet and PSPNet on sugarcane crop lines~\cite{ribeiro2024automated}.

\textbf{Yields Prediction.} Yield prediction regresses a continuous productivity estimate from the same imagery, and the usual design has two stages: extract a vegetation index or phenotype, then regress.
Random forests over Sentinel-2, RTK GPS, and LiDAR variables predict cotton yield~\cite{carneiro2023soil}, and an ensemble over Landsat-8 NDVI and agro-meteorological data reports R$^2$ 0.99 for tea yield on its test set~\cite{ramzan2023multimodal}.
UAV imagery earns its extra cost where the target is a fine-grained phenotype: Mask R-CNN extracts wheat-spike phenotypes that feed a yield regressor~\cite{PengWZYLRLSX23}, and a linear mixed-effect model on fixed-wing UAV multispectral imagery estimates rice leaf-area index and biomass~\cite{rs11111371}.

\textbf{Precision Irrigation.} Irrigation closes the loop, turning maps and estimates into water-delivery decisions~\cite{AbioyeAMBIROOR20}.
One line of work detects irrigation infrastructure at scale, such as center-pivot systems in Sentinel-2 series with Mask R-CNN~\cite{AlbuquerqueCSLB21} or irrigation types in NAIP aerial images with a U-Net and transfer learning~\cite{RaeiANSPA22}.
The other controls water locally: a drone system fuses RGB, thermal-infrared, and GPS data to find the driest zones of a field and sets valve flow with a random forest~\cite{JalajamonyNMF23}.

\subsection{Pivotal Task-2: Urban Planning, Logistics, and Transport}
\label{subsec:p4-sec11-2-pivotal-task-2-urban-planning-logistics-and-transport}
The second task turns GeoAI on the built environment. Figure~\ref{fig:p4_s11_2_urban} presents five complementary tasks on a shared urban context, from built structure and temporal change to services, transport, and green infrastructure.

\begin{figure}[!ht]
    \centering
    \includegraphics[width=0.98\linewidth]{figures/part4/section11/geoai-urban.pdf}
    \caption{GeoAI for the built environment: (a) building and land-use extraction from overhead and street-view imagery; (b) temporal urban change; (c) sensor-supported digital twins and operational services; (d) mobility, routing, and infrastructure analysis; and (e) canopy, parks, green roofs, and biosecurity monitoring. Tasks are complementary rather than a mandatory serial pipeline. Shared map geometry, temporal trends, and the city illustration are schematic; arrows denote time, data flow, or travel direction within the corresponding panels.}
    \label{fig:p4_s11_2_urban}
\end{figure}

\textbf{Urban Data Extraction.} Land use and built form are classified from overhead and street-level imagery, and fusing the two views resolves ambiguities that neither settles alone.
A fine-tuned EfficientNet classifies the 21 land-use classes of the UC Merced set~\cite{AlemK22}, while a dual-branch network that pairs dilated convolutions on Google Earth imagery with an attention-based multi-view network on Street View detects informal urban villages~\cite{ChenFNYGYGL22}.

\textbf{Urban Dynamics.} Temporal models capture how urban form changes over time.
An LSTM with attention simulates urban densification from land-use rasters of 2000, 2010, and 2020~\cite{el2023long}, and a pix2pix GAN predicts land-use and land-cover change by conditioning Landsat imagery on planning indicators~\cite{SunMFWH21}.
A static quantity such as urban solar exposure needs no temporal model; a feed-forward network on building geometry, occupancy, and orientation suffices~\cite{tehrani2024predicting}.

\textbf{Smart Cities.} The service layer turns these representations into live operation: video georeferencing and small-object detection inside an urban digital twin~\cite{HuLYHCYLY23}, LSTM parking-occupancy prediction served to a mobile application~\cite{CanliT21}, and detection of commercial signage in Street View imagery to map urban commerce~\cite{YeWKXC19}.

\textbf{Logistics and Transport.} Movement through the city splits into two problems with different tools.
Flow forecasting is temporal: bus arrival times are predicted from a traffic-density matrix with recurrent and convolutional models~\cite{PanovskiZ20}, a hybrid GRU--LSTM forecasts traffic from OpenStreetMap and GIS features~\cite{ZafarHSCM22}, IODRNN decomposes its outputs incrementally to remove the lag in GNSS-based forecasts~\cite{LuMPXC24}, and itinerary recommendation adds a denoising autoencoder with collaborative learning~\cite{ChenZCWC20}.
Infrastructure mapping is a segmentation problem: open LiDAR point clouds inventory rockfall-prone road slopes~\cite{RuaNAM23}, SegNet on RapidEye imagery updates rural road networks~\cite{kearney2020maintaining}, and U-Net with PointNet++ performs panoptic segmentation of railway environments from multimodal LiDAR~\cite{grandio2023multimodal}.

\textbf{Green Infrastructure and Bio Security.} Vegetation monitoring moves from coverage to function to threat, and each step demands finer spatial and spectral resolution than the last.
A semi-supervised DeepLabv3+ maps urban tree cover nationwide from high-resolution satellite imagery~\cite{guo2023nationwide}; a Faster R-CNN tree inventory identifies seven canopy species and, through geographically weighted regression, their effect on land-surface temperature~\cite{das2022geoai}; and biosecurity surveillance of invasive forest pests uses multispectral aerial imagery, where an instance-segmentation model reached an F1-score of 79\% against 58\% for a tile classifier~\cite{carnegie2023airborne}.

\subsection{Pivotal Task-3: Environmental Monitoring}
\label{subsec:p4-sec11-2-pivotal-task-3-environmental-monitoring}
The third task observes the Earth system itself, across the three fronts of Figure~\ref{fig:p4_s11_2_environment}.

\begin{figure}[!ht]
    \centering
    \includegraphics[width=0.98\linewidth]{figures/part4/section11/geoai-environment.pdf}
    \caption{Three parallel GeoAI tasks for environmental monitoring: (a) climate and vegetation (NDVI/LAI trends), (b) atmospheric chemistry (NO$_2$, SO$_2$, aerosols), and (c) environmental impact of land-cover change. Maps and trends are illustrative, not measured changes.}
    \label{fig:p4_s11_2_environment}
\end{figure}

\textbf{Climate Change.} Long-term land-surface change is read mostly through vegetation records spanning decades of Landsat, Sentinel, and MODIS data, and the methods trade spatial detail against record length.
Lou \etal~\cite{lou2023deep} map land use and land cover on Google Earth Engine with a deep network and find rapid greening of the Qinghai--Tibet Plateau from 1988 to 2018.
Deep-learning-enhanced spatiotemporal fusion reconstructs long-term NDVI in the Danjiang River Basin~\cite{wang2023applying}, random forests retrieve vegetation phenology from MODIS more accurately than rule-based methods~\cite{XinLLLZ20}, and a bidirectional LSTM produces the GLASS 250\,m leaf-area-index product (version 6) from MODIS reflectance~\cite{ma2022development}.

\textbf{Atmospheric Chemistry.} Atmospheric composition is retrieved by learning the map from satellite radiances and meteorology to ground-measured concentrations.
Nitrogen dioxide is the most studied target, estimated hourly from Sentinel-5P, meteorology, and sensor data with XGBoost~\cite{kim2021importance} and daily across China with a random forest and K-means clustering over OMI data~\cite{dou2021estimates}.
A convolutional autoencoder coupled to WRF-Chem models volcanic SO$_2$ dispersion~\cite{opio2023modeling}, AQNet fuses Sentinel-2 and Sentinel-5P to predict NO$_2$, O$_3$, and PM$_{10}$~\cite{rowley2023predicting}, an LSTM forecasts five pollutants from ground-station series~\cite{seng2021spatiotemporal}, and a deep network derives aerosol absorption from the OMI ultraviolet aerosol index, validated against AERONET~\cite{SunVVL21}.

\textbf{Environmental Impacts.} Impact assessment links land-use change to ecological outcomes, and it is judged not only on accuracy but on whether it can name the responsible driver, which is why explainability recurs here.
Neural networks on multi-temporal Landsat land-cover maps assess urban ecological quality~\cite{alqadhi2024applying}; multiscale geographically weighted regression, random forests, and SHAP relate agricultural land to PM$_{2.5}$ concentrations~\cite{gianquintieri2024implementation}; and interpretable ensembles measure the effect of climate change on agricultural land suitability~\cite{ShevchenkoLTBGKNSKM24}.

\subsection{Pivotal Task-4: Water Resources Management and Precipitation Forecasting}
\label{subsec:p4-sec11-2-pivotal-task-4-water-resources-management-and-precipitation-forecasting}
The fourth task addresses four complementary aspects of water resources in Figure~\ref{fig:p4_s11_2_water}.

\begin{figure}[!ht]
    \centering
    \includegraphics[width=0.98\linewidth]{figures/part4/section11/geoai-water.pdf}
    \caption{Four GeoAI tasks for water resources: (a) surface hydrology, (b) groundwater from GRACE storage anomalies and auxiliary data, (c) water quality calibrated against in-situ samples, and (d) radar-based precipitation nowcasting. The panels are parallel tasks, not successive stages of the water cycle; all values are schematic.}
    \label{fig:p4_s11_2_water}
\end{figure}

\textbf{Hydrological Modeling.} Convolutional models dominate local surface geometry.
Space-time image velocimetry on UAV video estimates river velocity and discharge, validated with LiDAR-derived terrain~\cite{islam2023mutual}; U-Net extracts drainage ditches from LiDAR in low-relief terrain~\cite{du2024drainage}; quadtree decomposition with multi-scale CNNs extracts coastlines at lower cost than FCN and U-Net baselines~\cite{sun2024quadtree}; and a reflectance-transformation CNN maps bathymetry from Landsat and Sentinel-2 imagery~\cite{ChanZX0WA24}.
At basin scale, where streamflow depends on several precipitation products (MODIS, CHIRPS, MSWEP), super-ensembles that weight multiple regressors take over~\cite{wegayehu2023super}.

\textbf{Groundwater.} No optical sensor sees the water table, so groundwater models work from proxy variables and gravimetry.
Static delineation of recharge or shallow-groundwater zones uses tree ensembles and deep networks over hydrogeological and remote-sensing variables~\cite{al2024hybrid,sahour2022identification,nhu2024new}.
GRACE/GRACE-FO measures total-water-storage changes; isolating groundwater requires auxiliary estimates of other water-storage components. Storage forecasting uses these gravity-derived anomalies, which a CNN--LSTM forecasts with an RMSE of 44.92\,mm/month~\cite{seo2021predicting} and which machine-learning regressors downscale to catchment resolution~\cite{sabzehee2023enhancing}.

\textbf{Water Quality.} Water quality is regressed from spectral reflectance against in-situ samples, and accuracy follows the optical signature of the constituent.
Pigments with a clear signature are tractable: fusion of data and classifiers through higher-order SVD estimates chlorophyll-a in Lake Nicaragua~\cite{sun2021ensemble}, and on UAV multispectral and hyperspectral data of a karst wetland, transformers and gradient boosting predict phycocyanin, dissolved oxygen, and chlorophyll-a, interpreted with SHAP~\cite{yang2023monitoring}.
Dissolved nutrients remain harder because their optical fingerprint is faint, whether estimated from MODIS with a spatiotemporal deep belief network~\cite{wu2022spatiotemporal} or from Landsat-8 with an ANN~\cite{vakili2020determination}.
Optical water types themselves can be classified from Sentinel-2 with an SVM~\cite{dasilva2021machine}.

\textbf{Precipitation Forecasting.} Precipitation nowcasting predicts the next radar frames from the current sequence.
HST-AFP fuses radar echoes with weather forecasts through attention~\cite{niu2023heterogeneous}, STIN adds humidity, wind, and temperature~\cite{jin2023spatiotemporal}, and PFST-LSTM adds pseudo-flow estimation to follow echo motion~\cite{luo2020pfst}; at longer range, multi-model ensembles post-process numerical weather prediction~\cite{sengoz2023machine}.

\subsection{Pivotal Task-5: Disaster Management}
\label{subsec:p4-sec11-2-pivotal-task-5-disaster-management}
The fifth task covers the disaster lifecycle (hazard or susceptibility mapping, event observation, and damage assessment) for earthquakes, floods, and forest fires, arranged as a grid in Figure~\ref{fig:p4_s11_2_disaster}.

\begin{figure}[!ht]
    \centering
    \includegraphics[width=0.98\linewidth]{figures/part4/section11/geoai-disaster.pdf}
    \caption{GeoAI for disaster management as a lifecycle-by-hazard grid: earthquakes, floods, and forest fires (columns) across hazard or susceptibility mapping, event observation, and post-event assessment (rows). All maps and severity scales are schematic.}
    \label{fig:p4_s11_2_disaster}
\end{figure}

\textbf{Earthquakes.} For earthquakes the lifecycle includes hazard mapping, event-related observation, and damage assessment; these tasks should not be equated with reliable prediction of an exact earthquake.
A CNN over elevation, vector, and attribute data maps seismic risk in north-eastern India~\cite{jena2021earthquake}, and encoder--decoder LSTMs~\cite{xiong2022gnss} and neural networks~\cite{hafeez2021possible} detect ionospheric anomalies in GNSS total-electron-content data around events such as the 2019 Pakistan earthquake.
After the event, LiDAR, orthophotos, and multispectral imagery estimate building vulnerability~\cite{torres2019integration}, and change detection on multi-temporal photogrammetric point clouds classifies structural damage into grades~\cite{zahs2023classification}.

\textbf{Floods.} Floods lean on radar.
Forecasting fuses meteorological, hydrological, and crowdsourced data, either comparing classical learners under cross-validation~\cite{PuttinaovaratH20} or with a ConvLSTM that learns spatial rainfall patterns and their temporal dynamics jointly~\cite{MoishinDPRA21}.
Inundation mapping falls to SAR segmentation, which sees through cloud, as in an attention-augmented U-Net on dual-polarized Sentinel-1 imagery~\cite{LiWCWJG22}.

\textbf{Forest Fires.} Fire has the fullest lifecycle.
Susceptibility is predicted from drone imagery with temperature and humidity sensing~\cite{rathod2023multipurpose}, or from elevation, vegetation indices, and hydrometeorology with a neuro-fuzzy system tuned by an artificial bee colony~\cite{PHAM2024102392}.
Detection is trained on drone RGB and infrared imagery such as the FLAME 2 dataset~\cite{ChenHWOARFCRW22} and can be distributed across sites with federated learning~\cite{SiddiqueADBAA24}.
Burn scars are delineated afterward by a U-Net on PlanetScope RGB and near-infrared imagery, with F1 up to 0.94~\cite{ChoPKKLS23}.

\subsection{Pivotal Task-6: Healthcare}
\label{subsec:p4-sec11-2-pivotal-task-6-healthcare}
The sixth task brings GeoAI to public health, where spatial signals couple disease and environment (Figure~\ref{fig:p4_s11_2_health}).

\begin{figure}[!ht]
    \centering
    \includegraphics[width=0.98\linewidth]{figures/part4/section11/geoai-health.pdf}
    \caption{Three parallel GeoAI tasks for public health: (a) disease surveillance and incidence forecasting, (b) spatial spread modeled by graph message passing over regions, and (c) pollution exposure assessment. The panels are not causally linked; maps and intervals are illustrative, not surveillance results.}
    \label{fig:p4_s11_2_health}
\end{figure}

\textbf{Infectious Diseases and Toxins.} Disease surveillance gains most from leading indicators, since search queries and geotagged posts surface an outbreak before clinical reporting does.
Influenza forecasting combines CDC surveillance data with Google Flu Trends and climate in a multi-stage LSTM~\cite{DBLP:journals/access/VennaTGRMN19}, with geotagged Twitter posts in a neural network tuned by an artificial-tree algorithm~\cite{hu2018prediction}, and with search frequencies and health records in a network-based regression that exploits synchrony across U.S. states~\cite{lu2019improved}; dengue in China is forecast from climate and weekly case counts, with support-vector regression performing best~\cite{guo2017developing}.
FINDER shows the payoff most clearly: it geolocates web-search signals of foodborne illness to restaurants, and the restaurants it flagged were three times more likely to pose food-safety risks than those found by complaint-based inspection~\cite{SadilekCDMSBJRG18}.

\textbf{Pandemics.} A pandemic shifts the question to transmission across geography.
Geospatial similarity learning predicts respiratory-virus spread in new regions~\cite{ZhaoHJCWWYR23}, and spatio-temporal regression tracks COVID-19 in Brazil~\cite{da2021covid}.
Detection networks catch what an epidemic curve misses: gatherings of people and vehicles in satellite imagery~\cite{atek2022geospatial}, proximity between re-identified individuals in video~\cite{ZhangLCZJ23}, and the drop in PM$_{2.5}$ during lockdown, measured from MODIS and Sentinel-5P~\cite{TanLGJ22}.

\textbf{Public Health.} Chronic exposure links where people are to what they breathe.
On the mobility side, mobile health uses wireless devices for clinical and public-health purposes~\cite{kahn2010mobile}, and the MD2K center turns mobile-sensor streams into markers of smoking, craving, and stress~\cite{KumarAAaBCCCEEG15}.
On the pollutant side, a geographically weighted gradient-boosting machine~\cite{zhan2017spatiotemporal} and a neural network over MODIS aerosol optical depth and land-use features~\cite{di2016assessing} map daily PM$_{2.5}$ across China and the United States.
Exposure assessment is only as sharp as the coarser of the two spatial scales.

\begin{table}[!th]
\centering
\fontsize{8}{10}\selectfont
\setlength{\tabcolsep}{1.0mm}
\caption{Representative data sources and methods for GeoAI, one or two per pivotal task. Numbers are each paper's own test set and are not comparable across rows; cells without a number name the metric the paper reports.}
\label{tab:p4_s11_2_benchmarks}
\begin{tabular}{L{0.20\textwidth} L{0.24\textwidth} L{0.28\textwidth} L{0.20\textwidth}}
\toprule
\textbf{Pivotal task} & \textbf{Representative data source} & \textbf{Representative method} & \textbf{Headline metric} \\
\midrule
Precision agriculture & Sentinel-1 SAR & STMA~\cite{han2023spatio} & 96\% OA (Brandenburg) \\
Precision agriculture & Landsat-8 NDVI + meteorology & Ensemble~\cite{ramzan2023multimodal} & R$^2$ 0.99 \\
Urban planning & UC Merced land use & EfficientNet (fine-tuned)~\cite{AlemK22} & accuracy \\
Environmental monitoring & Landsat, 1988--2018 & Deep LUCC mapping on GEE~\cite{lou2023deep} & LUCC accuracy \\
Water resources & GRACE / precipitation & CNN--LSTM~\cite{seo2021predicting} & RMSE 44.92\,mm/month \\
Disaster (flood) & Sentinel-1 SAR & Attention U-Net~\cite{LiWCWJG22} & precision / recall / F1 \\
Disaster (forest fire) & PlanetScope (RGB+NIR) & U-Net~\cite{ChoPKKLS23} & F1 0.94 \\
Healthcare & CDC / Google Flu / climate & Multi-stage LSTM~\cite{DBLP:journals/access/VennaTGRMN19} & RMSE / MAPE \\
\bottomrule
\end{tabular}
\end{table}

\subsection{Established Tools and Systems}
\label{subsec:p4-sec11-2-established-tools-and-systems}

\textbf{Cloud geospatial platforms} provide planetary-scale storage and analysis next to the data. Google Earth Engine\footnote{\url{https://earthengine.google.com/}} offers a petabyte-scale archive of satellite imagery with server-side analysis, Microsoft's Planetary Computer\footnote{\url{https://planetarycomputer.microsoft.com/}} pairs open Earth-observation data with a Python/STAC API, and commercial providers such as Esri ArcGIS\footnote{\url{https://www.arcgis.com/}} and Descartes Labs\footnote{\url{https://www.descarteslabs.com/}} deliver hosted GeoAI pipelines for enterprise and government users.

\textbf{Open-source GIS and ML toolkits} carry the modeling layer. QGIS\footnote{\url{https://qgis.org/}} is the de facto open desktop GIS, while deep-learning libraries tailored to raster geospatial data, such as TorchGeo\footnote{\url{https://github.com/microsoft/torchgeo}} and Raster Vision\footnote{\url{https://rastervision.io/}}, bridge PyTorch to satellite imagery, and Sentinel Hub\footnote{\url{https://www.sentinel-hub.com/}} streamlines on-demand access and processing of Copernicus archives.

\textbf{Open Earth-observation programmes} supply the data foundation: the European Copernicus programme's Sentinel constellation and NASA/USGS Landsat provide free, calibrated, multi-decadal imagery, complemented by gravimetric (GRACE/GRACE-FO), atmospheric (OMI, Sentinel-5P), and high-resolution commercial (WorldView, PlanetScope) sources that recur throughout the methods above.

\subsection{Toward Geospatial Foundation Models and GeoLLMs}
\label{subsec:p4-sec11-2-toward-geospatial-foundation-models-and-geollms}
Almost every method surveyed above is a task-specific supervised model: a CNN, random forest, LSTM, or GAN trained from scratch for one crop, one hazard, or one pollutant. That design is being displaced by the same pretraining-and-transfer paradigm that reshaped the rest of the field (Chapter~\ref{sec:reasoning}), and the geospatial community has begun building its own foundation models. Masked-autoencoding backbones lead the shift: SatMAE~\cite{cong2022satmae} adapts the masked image model to the temporal and multi-spectral structure of satellite imagery, Scale-MAE~\cite{reed2023scalemae} makes the pretext task explicitly scale-aware so a single backbone serves the huge ground-sample-distance range Earth observation spans, and Prithvi~\cite{jakubik2023foundation} packages such a backbone as a generalist model fine-tunable across flood mapping, burn-scar delineation, and crop segmentation. A single pretrained encoder thus stands in for the dozens of bespoke feature extractors that populate the tasks above.

The language side is moving just as fast. RemoteCLIP~\cite{liu2024remoteclip} aligns remote-sensing imagery with text, giving open-vocabulary retrieval and zero-shot classification without per-class labels; GeoChat~\cite{kuckreja2024geochat} grounds a large vision-language model on overhead imagery so an analyst can query a scene, ask for region-level descriptions, and obtain referring localizations in dialogue; and SegEarth-R1~\cite{li2025segearth} pushes further into pixel-level geospatial reasoning, generating segmentation masks from implicit, reasoning-heavy queries rather than fixed class names. These systems inherit the multimodal reasoning and world-model machinery of Part~\ref{part:3} and specialize it to Earth observation, pointing toward GeoAI workflows driven by natural-language instructions rather than hand-built per-task pipelines.

\subsection{Challenges}
\label{subsec:p4-sec11-2-challenges}
Most of the methods above were trained and tested in one region on one sensor, and that pattern shapes the open problems: whether a model transfers across regions and sensors, how heterogeneous and uneven data are reconciled, whether pipelines run at the scale and latency that operations need, and what new failure modes the geospatial foundation models bring.

\paragraph{Generalization across Regions and Sensors.} Few of the studies surveyed test a model outside the region, season, or sensor it was trained on, so their reported accuracies say little about transfer. Crop calendars, building styles, and spectral responses vary from place to place, and a crop map or flood segmenter trained on one country or one satellite can degrade sharply on another. Pretrained backbones such as Prithvi~\cite{jakubik2023foundation} are motivated by exactly this gap, but systematic cross-region benchmarks remain rare.

\paragraph{Data Heterogeneity and Quality.} Geospatial workflows combine imagery, sensor streams, vector layers, and crowdsourced inputs that differ in resolution, revisit time, and reliability, and the volume of these archives makes curation itself demanding~\cite{boutayeb2024geoai}. Open-access data needs quality control before it can serve as training labels. Geospatial data is also strategically sensitive, which raises questions of data sovereignty, and adversarial manipulation of inputs is a threat that none of the methods above addresses.

\paragraph{Scale and Latency.} Operational services such as flood mapping during an event, precipitation nowcasting, and traffic forecasting must ingest new acquisitions continuously and update their models as satellites deliver data, which strains pipelines built around one-off training runs. When several models are chained in one workflow, validation and deployment cost grow with each link, pushing systems toward distributed, modular architectures.

\paragraph{Open Problems for Geospatial Foundation Models.} The foundation models of the previous subsection bring problems of their own: pretraining that fuses radar, optical, and hyperspectral sensors at very different resolutions; evaluation benchmarks that test geospatial reasoning rather than image-level recognition; and grounding language output in physically consistent, georeferenced coordinates, which dialogue systems such as GeoChat~\cite{kuckreja2024geochat} do not yet guarantee.

\newpage
\section{Aerospace Engineering}
\label{sec:p4_s11_3}

\subsection{Introduction}
\label{subsec:p4-sec11-3-introduction}

Aerospace engineering designs and operates systems that work within and beyond the Earth's atmosphere, and its missions are spatial problems at their core.
Rendezvous and docking (RVD), on-orbit servicing (OOS), debris inspection, and planetary landing and roving all depend on estimating three-dimensional structure and relative motion, often of a non-cooperative target that carries no markers, while communication delays rule out continuous ground control.
Both spacecraft autonomy (\eg, guidance, navigation, and control) and human-in-the-loop operations (\eg, situational awareness for ground controllers) therefore rest on spatial perception and reasoning over geometry, relative motion, semantics, and physical constraints.

Aerospace engineering is thus a demanding deployment target for the perception and reasoning capabilities surveyed in Part~\ref{part:3}: the detection, segmentation, pose-estimation, and 3D-reconstruction techniques of Chapter~\ref{chapter:3:sec:6} and the planning and decision-making methods of Chapter~\ref{sec:reasoning} are here specialized to non-cooperative orbital targets, extreme illumination, and long-horizon autonomy under radiation-hardened compute and communication latency.

\subsection{System Overview}
\label{subsec:p4-sec11-3-system-overview}
The value of spatial intelligence in aerospace lies less in any single perception module than in the closed autonomy loop it serves (Figure~\ref{fig:p4_s11_3_overview}): onboard and ground-based sensing (monocular and multi-view cameras, LiDAR, telemetry) feeds a perception-reasoning-decision core, whose commands change the scene that is sensed next.
The four pivotal tasks treated below are complementary stages of this loop.
\textbf{Satellite pose estimation} supplies the geometric anchor, recovering where a target is and how it is oriented;
\textbf{spacecraft detection, segmentation, and parts recognition} enrich that anchor with structural semantics, identifying what the target is and which functional components it exposes;
\textbf{operator agents} lift this perceptual state into high-level reasoning, planning, and human-facing decision support; and
\textbf{planetary robotics} closes the loop on the surface, coupling perception, localization, and planning into physical exploration and sampling.
Table~\ref{tab:p4_s11_3_benchmarks}, after Task-4, consolidates representative benchmarks and metrics across all four.

\begin{figure}[!th]
    \centering
    \includegraphics[width=0.98\linewidth]{figures/part4/section11/aerospace-overview.pdf}
    \caption{Aerospace spatial intelligence as a closed loop from onboard and ground sensing, through a perception--reasoning--decision core, to complementary pose, structural-semantic, planning, and surface-action outputs. Decisions may request additional observations; mission execution produces subsequent measurements. Pose and semantic estimates are information products rather than physical actions. Satellite views, point clouds, and rover imagery are schematic.}
    \label{fig:p4_s11_3_overview}
\end{figure}


\subsection{Pivotal Task-1: Satellite Pose Estimation}
\label{subsec:p4-sec11-3-pivotal-task-1-satellite-pose-estimation}

Vision-based satellite pose estimation predicts the 6-DoF relative position and orientation of a target spacecraft from monocular or multi-view imagery; it is the geometric anchor of the autonomy loop and a prerequisite for rendezvous, docking, and on-orbit servicing.
It specializes the 6-DoF estimation problem of Chapter~\ref{chapter:3:sec:6} to a hostile setting: harsh illumination, specular reflections, little texture, and a severe gap between the synthetic images models are trained on and the real images they meet in orbit.
Figure~\ref{fig:p4_s11_3_pose} traces the pipeline from method to evaluation.

\begin{figure}[!t]
    \centering
    \includegraphics[width=0.95\linewidth]{figures/part4/section11/aerospace-pose.pdf}
    \caption{Satellite pose estimation: a monocular image of a non-cooperative target passes through alternative feature-based, regression, or hybrid pipelines to a target pose relative to the camera. Translation and relative-rotation errors compare estimated and reference poses; 2D reprojection compares detected image keypoints with projected model points. The satellite render and keypoint residuals are schematic; no numerical benchmark scores are implied.}
    \label{fig:p4_s11_3_pose}
\end{figure}

Traditional methods relied on hand-engineered features (\eg, Harris corners, SIFT, SURF, and ORB) or model-based tracking~\cite{d2014poseestimationuncooperativespacecraft}, which are sensitive to noise and domain shift, and deep learning has since become dominant~\cite{proencca2020deepURSO}.
The first deep models regress the pose directly with CNNs~\cite{garcia2021lspnet}, as in the Spacecraft Pose Network (SPN)~\cite{sharma2020neuralposeestimation}, a GoogLeNet-based network~\cite{phisannupawong2020vision}, and a VGG backbone~\cite{sonawani2020assistive}, trained on synthetic data and on laboratory imagery that simulates an on-orbit assembly operation.
SPNv2~\cite{park2024robustspacecraftposeestimation} made multi-task learning the prevailing design: its shared encoder jointly predicts keypoints, a direct pose regression, and foreground segmentation, and it narrows the synthetic-to-real gap by online domain refinement, adapting the network's normalization layers on unlabeled target imagery at deployment.
Hybrid pipelines estimate a coarse pose with a deep model and refine it with geometry-based optimization, keeping the final estimate physically consistent; early work classified discretized viewpoints with a CNN~\cite{sharma2018pose}, and the same coarse-to-fine pattern now extends to 6-DoF tracking with stereo event cameras~\cite{liu2025stereo}.

SPEED~\cite{Sharma2019SpacecraftPE} and SPEED+~\cite{park2022speed+} are the standard testbeds (Table~\ref{tab:ssa_datasets}): both pair a large synthetic training set with laboratory hardware-in-the-loop test images, so that the synthetic-to-real gap is measured rather than assumed, and the Satellite Pose Estimation Challenge built on them fixed common training regimes, input formats, and metrics~\cite{kisantal2020satellitechallenge}.
CubeSat-CDT~\cite{musallam2022cubesatcdt} adds trajectories of a symmetric target, so that temporal methods and pose ambiguity can be evaluated.
Evaluation reports translation and rotation error as defined for SPEED~\cite{Sharma2019SpacecraftPE}, often with 2D reprojection or keypoint error~\cite{garcia2021lspnet}; results are given separately on synthetic and real test sets, and robustness to illumination, occlusion, and background clutter is probed explicitly~\cite{park2024robustspacecraftposeestimation}.

\begin{table}[!ht]
\centering
\caption{Overview of existing pose estimation datasets. \textit{sim} refers to simulated images and \textit{lab} to laboratory data. SPARK's 11 classes are 10 spacecraft and one debris class.}
\label{tab:ssa_datasets}
\fontsize{9}{11}\selectfont
\setlength{\tabcolsep}{1.0mm}
\begin{tabular}{L{0.17\textwidth} C{0.08\textwidth} C{0.11\textwidth} C{0.08\textwidth} C{0.13\textwidth} C{0.10\textwidth} C{0.11\textwidth} C{0.08\textwidth}}
\toprule
\multirow{2}{*}{\textbf{Dataset}} & \multicolumn{2}{c}{\textbf{Image}} & \multirow{2}{*}{\textbf{Classes}} & \multirow{2}{*}{\textbf{Resolution}}  & \multirow{2}{*}{\textbf{Rendering}} & \multirow{2}{*}{\textbf{Trajectories}} & \multirow{2}{*}{\textbf{Public}} \\
\cmidrule(r){2-3}
& \textbf{Synthetic} & \textbf{Non-synthetic} & & & & & \\
\midrule
SPARK~\cite{musallam2021spacecraft} & 150k & - & 11 & $1024 \times 1024$ & sim & \crossmark & \checkmark  \\
SPEED~\cite{Sharma2019SpacecraftPE} & 15k & 305 & 1 & $1920 \times 1200$  & sim + lab & \crossmark & \checkmark \\
SPEED+~\cite{park2022speed+} & 60k & 10k & 1 & $1920 \times 1200$ & sim + lab & \crossmark & \checkmark \\
URSO~\cite{proencca2020deepURSO} & 15k & - & 2 & $1080 \times 960$ & sim & \crossmark & \checkmark \\
SwissCube~\cite{hu2021widedepthrange} & 50k & - & 1 & $1024 \times 1024$ & sim & \checkmark & \checkmark \\
Cygnus~\cite{black2021realCygnus} & 20k & 540 & 1 & $1024 \times 1024$ & sim + real  & \crossmark & \crossmark \\
Prisma12K~\cite{park2019towardsPrisma12K} & 12k & - & 1 & $752 \times 580$ & sim  & \crossmark & \crossmark \\
Prisma25~\cite{park2019towardsPrisma12K} & - & 25 & 1 & -NA- & real  & \crossmark & \crossmark \\
CubeSat-CDT~\cite{musallam2022cubesatcdt} & 14k & 8k & 1 & $1440 \times 1080$ & sim + lab & \checkmark & \checkmark \\
\bottomrule
\end{tabular}
\end{table}


\subsection{Pivotal Task-2: Spacecraft Detection, Segmentation, and Parts Recognition}
\label{subsec:p4-sec11-3-pivotal-task-2-spacecraft-detection-segmentation-and-parts-recognition}

Detection, segmentation, and parts recognition enrich the geometric anchor of Task-1 with structural semantics, telling what the target is and which functional components it exposes; proximity operations, on-orbit servicing, debris inspection, and collision avoidance all depend on this.
The detection and segmentation machinery of Chapter~\ref{chapter:3:sec:6} meets the obstacles of Task-1 here (strong illumination variation, specular metallic surfaces, little texture) and, in addition, a scarcity of annotated real imagery.
The stack runs from coarse presence to fine structure: detection localizes the spacecraft, segmentation delineates its silhouette, and parts recognition resolves its components, as Figure~\ref{fig:p4_s11_3_detection} illustrates.

\begin{figure}
    \centering
    \includegraphics[width=0.95\linewidth]{figures/part4/section11/aerospace-detection.pdf}
    \caption{The ``detect--segment--recognize'' perception stack for non-cooperative spacecraft: (1) object localization, (2) encoder--decoder silhouette segmentation, and (3) recognition of solar panels, the main body, antenna, and thruster. The same synthetic target is used throughout. The silhouette is an illustrative foreground mask, not a measured prediction; arrows indicate increasing semantic detail rather than requiring every implementation to run these stages serially.}
    \label{fig:p4_s11_3_detection}
\end{figure}

Convolutional and transformer architectures now carry all three stages~\cite{shao2023satellite}.
Detection is formulated as bounding-box or keypoint localization, with lightweight saliency-driven detectors for onboard compute, while encoder--decoder networks, pyramid attention modules, and hybrid 2D/3D architectures delineate silhouettes and component boundaries~\cite{zhao20223dnoncooperativespacecraft}.
Parts recognition (solar panels, antennas, thrusters, reflectors) often adds semantic priors or 3D geometric reasoning to stay robust under occlusion.
Two recent lines reduce the dependence on manual labels: prompt-guided Segment Anything (SAM) automates the annotation of space targets~\cite{zhang2024automaticspacetarget}, and diffusion priors support few-shot segmentation and adaptation to unseen spacecraft~\cite{li2025exploitingdiffusion}.

The datasets in Table~\ref{tab:satellite_datasets_datasetwise} differ along a few axes.
Labels range from whole-object masks (SwissCube~\cite{hu2021widedepthrange}, a pose dataset that also provides CubeSat masks, and SPARK~\cite{musallam2021spacecraft}) and object-level boxes (STAR-24K~\cite{zhang2022star24K}) to component labels, from a single component (the solar panels of the Tsinghua dataset~\cite{li2020noncooperativeobjectdetection}) up to nine (Spacecraft-DS~\cite{cao2024spacecraftds}).
Most are rendered (Unity3D for SPARK, Unreal Engine~4 for UESD~\cite{zhao2023intelligentUESD}), whereas Spacecraft-DS uses hardware-in-the-loop capture to narrow the synthetic-to-real gap.
The Adelaide dataset~\cite{dung2021spacecraftdataset} fixed the common three-component scheme (solar panel, main body, antenna), which Shao \etal~\cite{shao2023satellite} extended to video sequences rendered under three illumination conditions and several relative-motion modes; Zhao \etal~\cite{zhao20223dnoncooperativespacecraft} carried the same labels to 3D point clouds, and NCSTP~\cite{liu2025largencstp} raised the scale to 200,000 images of satellites, debris, and rocks.
Detection is scored by mean average precision and segmentation by IoU, as in Chapter~\ref{chapter:3:sec:6}.

\begin{table}[!ht]
\centering
\caption{Summary of satellite recognition datasets. `-' indicates that the exact quantity is unknown; the Zhao \etal dataset consists of point clouds rather than images.}
\label{tab:satellite_datasets_datasetwise}
\fontsize{9}{11}\selectfont
\setlength{\tabcolsep}{1.5mm}
\begin{tabular}{L{0.32\textwidth} C{0.12\textwidth} C{0.16\textwidth} C{0.16\textwidth} C{0.13\textwidth}}
\toprule
\textbf{Dataset}
& \textbf{\#Images}
& \textbf{\#Satellite Models}
& \textbf{Segmentation Level}
& \textbf{\#Components}   \\
\midrule
SwissCube~\cite{hu2021widedepthrange} & 50{,}000 & 1 & Global & 0  \\
UESD~\cite{zhao2023intelligentUESD} & 10{,}000 & 33 & Component & 5 \\
SPARK~\cite{musallam2021spacecraft} & 150{,}000 & 10 & Global & 0  \\
STAR-24K~\cite{zhang2022star24K} & 24{,}451 & - & Global & - \\
NCSTP~\cite{liu2025largencstp} & 200{,}000 & 16 & Component & 4 \\
Tsinghua Satellite Dataset~\cite{li2020noncooperativeobjectdetection} & 18{,}381 & - & Component & 1 \\
Adelaide Satellite Dataset~\cite{dung2021spacecraftdataset} & 3{,}117 & - & Component & 3   \\
Satellite Component Video Dataset~\cite{shao2023satellite} & 32{,}402 & 13 & Component & 3  \\
Spacecraft-DS~\cite{cao2024spacecraftds} & 1{,}602 & 2 & Component & 9 \\
3D component point clouds (Zhao \etal)~\cite{zhao20223dnoncooperativespacecraft} & - & 33 & Component & 3 \\
\bottomrule
\end{tabular}
\end{table}


\subsection{Pivotal Task-3: Operator Agents}
\label{subsec:p4-sec11-3-pivotal-task-3-operator-agents}

Operator agents lift the perceptual state recovered by Tasks 1 and 2 into reasoning, planning, and decision support, monitoring, interpreting, and controlling spacecraft systems to reduce operator workload and allow more autonomous mission execution.
They are the aerospace instance of the planning and decision-making capabilities of Chapter~\ref{sec:reasoning}.
Rooted in rule-based and model-driven architectures, their role has expanded from health monitoring and scheduling to high-level control in dialogue with human experts, a history that runs from symbolic planning through multi-agent coordination to foundation-model reasoning (Figure~\ref{fig:p4_s11_3_operator}).

\begin{figure}
    \centering
    \includegraphics[width=0.95\linewidth]{figures/part4/section11/aerospace-operator.pdf}
    \caption{Three paradigms in the development of spacecraft operator agents: (1) symbolic planning and expert systems, (2) multi-agent coordination across a constellation, and (3) LLM- and RAG-based agents that combine telemetry with retrieved domain knowledge to propose natural-language plans. Inter-panel arrows indicate historical development, not required dataflow. In the rule example, observations and action outcomes support the goal; blue peer links and dashed ground links denote communication. Satellite imagery and rules are illustrative.}
    \label{fig:p4_s11_3_operator}
\end{figure}

Early operator-agent architectures were built on symbolic reasoning, constraint-based planning, and expert-system principles.
The Remote Agent Experiment~\cite{bernard1998designremoteagent} and early autonomous spacecraft-agent prototypes~\cite{pell1997autonomousspacecraftagent} demonstrated onboard planning, diagnosis, and execution in deep-space missions.
As autonomy requirements grew, multi-agent planning frameworks~\cite{colby2016autonomousmultiagent,picard2021autonomousAgents} were introduced to coordinate heterogeneous spacecraft, manage distributed sensing, and schedule jointly across satellite constellations.
Large language models changed the design: recent work~\cite{rodriguez2024languagemodelsspacecraft,carrasco2025visualoperatoragents} uses LLMs as operator agents that interpret telemetry, generate action plans, respond to anomalies, and follow natural-language mission specifications, and fine-tuned LLMs~\cite{zucchelli2025finetunellmSpace,carrasco2024finetuneAutonomousSpacecraft,foutter2024spacellava,carrasco2025llmAutonomousSpacecraft} have been tested in simulators such as Kerbal Space Program (KSP) for guidance, navigation, and control. To meet the interpretability and safety requirements of mission-critical decisions, other studies integrate retrieval-augmented generation (RAG)~\cite{belo2025languageragSpaceSpace}, domain knowledge graphs~\cite{powell2025questionknowledgegraphssatellite}, and structured symbolic reasoning, while hybrid architectures such as MCP-based multi-component agents~\cite{wu2025spacemind} combine small specialized models with foundation models to balance computational cost and reasoning depth, particularly for on-orbit servicing.

Evaluation environments followed the same path.
The Remote Agent Experiment~\cite{bernard1998designremoteagent} supplied real mission scenarios for onboard planning and diagnosis, and multi-agent planning studies~\cite{colby2016autonomousmultiagent} supply constellation-level scheduling, resource-allocation, and communication constraints.
LLM agents are tested in simulation, chiefly in KSP~\cite{carrasco2025llmAutonomousSpacecraft} and in the KSPDG suite of SpaceGym~\cite{allen2023spacegym}, a set of multi-agent pursuit--evasion differential games built in the KSP engine.
Classical agents are scored on planning optimality, constraint satisfaction, fault-diagnosis accuracy, and task completion; LLM agents are additionally judged on adherence to safety constraints, robustness to shifts in telemetry or mission context, and the interpretability of their decisions, which is why uncertainty estimation and human-in-the-loop verification recur throughout this work.


\subsection{Pivotal Task-4: Planetary Robotics}
\label{subsec:p4-sec11-3-pivotal-task-4-planetary-robotics}

Planetary robotics closes the autonomy loop on the surface, coupling the perception, localization, and planning capabilities of Chapters~\ref{chapter:3:sec:6} and~\ref{sec:reasoning} into exploration, navigation, and sampling on the Moon, Mars, and small bodies.
Unlike orbital operations, surface missions face uneven and unstructured terrain, delayed and limited communication, restricted sensing, and long stretches of autonomy without continuous human supervision.
Figure~\ref{fig:p4_s11_3_planetary} layers the autonomy stack from perception and mapping through hierarchical decision-making to vision-guided manipulation.

\begin{figure}
    \centering
    \includegraphics[width=0.99\linewidth]{figures/part4/section11/aerospace-planetary.pdf}
    \caption{The planetary-robotics autonomy stack: visual-odometry feature tracks and 3D terrain mapping support hierarchical decisions and vision-guided manipulation. Mission reasoning handles goals, task decomposition, resources, and replanning; path planning uses global/local maps and constraints; reactive locomotion handles obstacles, stability, and terrain adaptation. Solid downward links convey plans, and dashed upward links convey status. Manipulation yields samples, material observations, environment updates, and task progress. Scenes and map geometry are illustrative.}
    \label{fig:p4_s11_3_planetary}
\end{figure}

Surface robots are organized around a mobility stack that couples perception, localization, mapping, planning, and control, with manipulation and science operations integrated at the task level.
Feature tracking, visual odometry, and stereo reconstruction estimate rover motion and build local 3D models of rugged terrain, a capability flown since the Mars Exploration Rovers~\cite{maimone2007two}.
The Enhanced AutoNav (ENav) system on NASA's Perseverance rover~\cite{toupet2025enav} builds onboard 3D terrain maps, evaluates hazards, and plans safe local traverses under severe computational and communication constraints, and it evaluated terrain over a large majority of the rover's traversed distance.
This autonomy remains embedded in a human-supervised workflow: ground operators specify science goals and route segments, and onboard autonomy resolves local navigation and hazard avoidance.

Research and analogue-testbed systems extend these capabilities to multiple robots, where cooperative visual navigation and shared mapping improve coverage and robustness~\cite{bahraini2021autonomousplanetaryrobots}.
Hierarchical frameworks decompose behavior across temporal and functional layers, balancing reactive locomotion and obstacle avoidance with deliberative path planning, task scheduling, and mission-level reasoning under safety, power, communication, and resource constraints.

Vision-guided manipulation supports sample collection and instrument deployment by coupling image-based pose estimation with position-controlled or compliant manipulation suited to contact with natural terrain~\cite{nickels2010visionplanetaryrobotics}.
Learning-based methods are increasingly explored for terrain classification, slip prediction, and traversability estimation, but model-based perception, localization, and planning remain central: safety requirements are stringent, mission-representative labeled data is scarce, onboard computation is constrained, and purely data-driven methods are hard to validate for mission-critical use.
Evaluation therefore combines physics-based simulators, mission-derived datasets, and laboratory or analogue-field testbeds that model illumination, dust, wheel--soil interaction, and communication delay, and it reports localization accuracy, map quality, hazard-aware traversal, task completion, and energy use.

\begin{table}[!ht]
\centering
\caption{Representative benchmarks and headline metrics for the four pivotal aerospace tasks. The headline column names the standard evaluation axis, not a single best score.}
\label{tab:p4_s11_3_benchmarks}
\fontsize{8}{10}\selectfont
\setlength{\tabcolsep}{1.0mm}
\begin{tabular}{L{0.18\textwidth} L{0.24\textwidth} L{0.30\textwidth} L{0.20\textwidth}}
\toprule
\textbf{Pivotal task} & \textbf{Benchmark} & \textbf{Representative method / system} & \textbf{Headline metric} \\
\midrule
Pose estimation & SPEED / SPEED+~\cite{Sharma2019SpacecraftPE,park2022speed+} & Multi-task / domain-adaptive CNN~\cite{park2024robustspacecraftposeestimation} & Translation + rotation error \\
Pose estimation & CubeSat-CDT~\cite{musallam2022cubesatcdt} & Cross-domain regression~\cite{garcia2021lspnet} & 2D reprojection error \\
Detection \& parts & Adelaide / SPARK~\cite{dung2021spacecraftdataset,musallam2021spacecraft} & Encoder-decoder segmentation~\cite{zhao20223dnoncooperativespacecraft} & mAP \\
Detection \& parts & NCSTP~\cite{liu2025largencstp} & SAM-based prompt segmentation~\cite{zhang2024automaticspacetarget} & mAP / IoU \\
Operator agents & KSPDG / SpaceGym~\cite{allen2023spacegym} & Fine-tuned LLM agent~\cite{carrasco2025llmAutonomousSpacecraft} & Maneuver success rate \\
Operator agents & Deep Space 1 flight experiment~\cite{bernard1998designremoteagent} & Symbolic planner / diagnosis~\cite{pell1997autonomousspacecraftagent} & Task completion \\
Planetary robotics & Mars rover traverses~\cite{toupet2025enav} & Visual odometry~\cite{maimone2007two} & Localization drift \\
\bottomrule
\end{tabular}
\end{table}


\subsection{Established Tools and Systems}
\label{subsec:p4-sec11-3-established-tools-and-systems}

\textbf{Simulators and rendering engines.}
Because real orbital imagery is scarce and flight opportunities are expensive, most perception research is trained and validated in simulation.
General-purpose game and physics engines have been repurposed as space testbeds: Unity3D and Blender drive the synthetic rendering pipelines behind the SPARK~\cite{musallam2021spacecraft} and satellite-component~\cite{shao2023satellite} datasets, while the Kerbal Space Program\footnote{\url{https://www.kerbalspaceprogram.com/}} engine underpins LLM-agent guidance-navigation-control studies~\cite{carrasco2025llmAutonomousSpacecraft}.
For decision-making, SpaceGym\footnote{\url{https://github.com/mit-ll/spacegym-kspdg}}~\cite{allen2023spacegym} packages its OrbitDefender2D gridworld games and the KSPDG suite as standardized non-cooperative game environments, and robotics simulators such as NVIDIA Isaac Sim\footnote{\url{https://developer.nvidia.com/isaac/sim}} and Gazebo\footnote{\url{https://gazebosim.org/}} provide the wheel-soil and sensor models used to prototype planetary-rover autonomy.

\textbf{Public challenges.}
The Satellite Pose Estimation Challenge~\cite{kisantal2020satellitechallenge}, run by Stanford's Space Rendezvous Laboratory and the European Space Agency on its Kelvins platform\footnote{\url{https://kelvins.esa.int/}}, set the protocols that pose estimation now follows, and its SPEED+ edition~\cite{park2022speed+} made the synthetic-to-real gap the explicit test.

\textbf{Flight-proven autonomy systems.}
The NASA Remote Agent~\cite{bernard1998designremoteagent}, flown on Deep Space 1, remains the canonical reference for onboard planning, diagnosis, and execution.
On planetary surfaces, rover visual odometry~\cite{maimone2007two} and the Enhanced AutoNav stack flown on Perseverance~\cite{toupet2025enav} set the operational baseline against which hierarchical-autonomy and learning-augmented methods are measured, and current work pairs such flight-qualified stacks with foundation-model agents~\cite{wu2025spacemind}.


\subsection{Challenges}
\label{subsec:p4-sec11-3-challenges}
Orbital and planetary environments differ from terrestrial ones in ways that terrestrial data does not cover, and the open problems follow the pipeline: the data gap that limits perception, the compute budget that limits deployment, and the trustworthiness requirements that limit certification.

\paragraph{The Unresolved Sim-to-Real Domain Gap.}
Autonomous space systems, such as those performing on-orbit servicing or debris removal, are trained largely on simulated data, and simulators do not reproduce the high-dynamic-range lighting of orbit (\eg, solar glare and deep shadow), intricate material reflections, or sensor degradation such as radiation noise.
Models consequently lose accuracy on real imagery, as the hardware-in-the-loop test sets of the pose-estimation challenges show~\cite{kisantal2020satellitechallenge,park2022speed+}, which keeps domain adaptation and domain randomization central to the field.

\paragraph{Computational Constraints and Real-time Operation.}
Spacecraft hardware is bound by size, weight, and power (SWaP) limits and must use radiation-hardened processors that are far less powerful than ground-based GPUs.
Running 6-DoF pose estimation or SLAM at the update rates that proximity operations require on such computers calls for model compression, lightweight architectures, and hardware-aware optimization that preserve accuracy while cutting latency and power consumption.

\paragraph{Assuring Trustworthiness and Robustness in Extreme Conditions.}
Space missions are safety-critical, so perception must hold up under extreme visual ambiguity, such as estimating the pose of highly symmetric targets like CubeSats~\cite{musallam2022cubesatcdt} or navigating where standard reference signals are unavailable.
Deployed systems also explain little about their own failures, so operators cannot easily trace a perception anomaly (\eg, a pose error caused by sensor saturation) to its root cause in time to intervene.
Formal verification and interpretable models that guarantee predictable behavior across operating extremes remain open problems.

\newpage
\chapter{Spatial Intelligence for Biomedicine and Natural Sciences}


\FramedInitial{B}{iomedicine} and the natural sciences are, at their core, spatial disciplines: their objects are structures in space, from organs and tissues down to cells and molecules, and understanding them means reasoning about shape, arrangement, and interaction in three dimensions. This chapter turns the spatial-intelligence capabilities of the previous part onto that setting, where in both of its cases the unit of value is the same: a metrically faithful 3D model of a specific physical system, kept in correspondence with real measurements.

In medicine, that model is patient-specific, and it drives a closed loop running from imaging through planning to image-guided intervention. In the natural sciences the same perception and reconstruction machinery operates several orders of magnitude smaller, at the molecular and cellular scale, where the structures reconstructed are the very objects of scientific explanation.


\section{Medicine and Clinical Healthcare}
\label{sec:p4_s12_1}

\subsection{Introduction}
\label{subsec:p4-sec12-1-introduction}

Anatomy is a three-dimensional arrangement of organs with shape, volume, and topology; disease is \emph{localized}, so that a lesion is a position and an extent in the body frame rather than a label on an image; and intervention is fundamentally an act of navigation, guiding an instrument to a target through tissue that must not be harmed.

The perception machinery of Chapter~\ref{chapter:3:sec:6} (detection, segmentation, depth, pose) and the reconstruction and generation of Chapter~\ref{sec:generation} reappear here, specialized to volumetric data, to registration across modalities and time, and to a regulatory setting in which a wrong answer harms a patient. 
The unit of analysis is the individual patient and the individual procedure: the molecular and cellular scale belongs to the natural sciences (§\ref{sec:Natural Science}), the generic digital twin, of which the patient-specific medical twin below is a specialization, to §\ref{sec:p4_s11_1}, and the population scale of spatial epidemiology to the GeoAI of §\ref{sec:p4_s11_2}.

\subsection{System Overview}
\label{subsec:p4-sec12-1-system-overview}

Figure~\ref{fig:p4_s12_1_overview} draws clinical spatial intelligence as a closed loop around the patient: multimodal imaging (computed tomography, magnetic resonance, ultrasound, endoscopy, pathology) is acquired and the anatomy perceived, a patient-specific 3D model is built and a plan formed, the plan is executed under image guidance, and the outcome is imaged anew. 
Registration connects the image, patient, and instrument frames throughout this pipeline; the recurring difficulty is keeping the virtual model aligned with the physical patient as tissue moves and deforms. The three pivotal tasks map onto the loop directly (Task-1 perceives, Task-2 models and plans, Task-3 acts), and Table~\ref{tab:p4_s12_1_benchmarks} collects their benchmarks.

\begin{figure}
    \centering
    \includegraphics[width=0.98\linewidth]{figures/part4/section12/medical_overview.pdf}
    \caption{Clinical spatial intelligence with two connected information paths: acquired observations support anatomical perception and patient-specific planning, while live imaging and tracking update the alignment used for guided intervention. Physical intervention changes the patient state, prompting subsequent observation. Anatomical scenes are schematic; pathology may require specimen correspondence.}
    \label{fig:p4_s12_1_overview}
\end{figure}

\subsection{Pivotal Task-1: Volumetric Medical Image Perception}
\label{subsec:p4-sec12-1-pivotal-task-1-volumetric-medical-image-perception}

The first task perceives the anatomy. Medical perception is intrinsically volumetric: it converts a voxel grid into delineated 3D anatomical objects with measurable shape and volume, localizes disease in the body frame, and aligns observations across modalities and time. Figure~\ref{fig:p4_s12_1_perception} illustrates three complementary outputs: segmented organs, a localized candidate finding, and a cross-modal alignment; these tasks need not form a serial pipeline.

\begin{figure}
    \centering
    \includegraphics[width=0.98\linewidth]{figures/part4/section12/medical_perception.pdf}
    \caption{Three input-operation-output views of medical perception: (a) image volumes yield voxel labels, organ masks, and 3D anatomy; (b) detected candidate regions are mapped from voxel indices to physical patient coordinates using image geometry; and (c) moving and reference modalities jointly determine a spatial transform for anatomical alignment. Synthetic CT/MRI-style images and the hypothetical ROI illustrate task structure, not measured clinical findings or a validated registration result.}
    \label{fig:p4_s12_1_perception}
\end{figure}

\textbf{Volumetric segmentation.} Segmentation is the spatial core of medical imaging, and its history is a move from two-dimensional slices to full volumes. 
The encoder-decoder U-Net~\cite{ronneberger2015unet} founded the area, and its volumetric successors, the 3D U-Net~\cite{cciccek20163d} and the V-Net~\cite{milletari2016vnet} (which introduced the now-ubiquitous Dice loss to handle voxel class imbalance), moved it into three dimensions. 
The self-configuring nnU-Net~\cite{isensee2021nnunet} adapts preprocessing, architecture, and training to each dataset and remains the baseline that is hard to beat, while transformer encoders such as UNETR~\cite{hatamizadeh2022unetr} and Swin-UNETR~\cite{hatamizadeh2022swinunetr} add global context. 
At the whole-body scale, TotalSegmentator~\cite{wasserthal2023totalsegmentator} delineates more than a hundred anatomical structures in a single CT, functioning as a spatial atlas in routine use. 
These methods are evaluated by the Dice similarity coefficient for volumetric overlap and by boundary-geometry metrics, the 95th-percentile Hausdorff distance and the average symmetric surface distance, that capture shape errors which overlap alone misses~\cite{taha2015metrics}.

\textbf{Medical foundation models.} The recent shift mirrors the rest of vision: from per-task specialists to promptable foundation models that segment arbitrary anatomy from spatial prompts. 
The Segment Anything Model~\cite{kirillov2023sam} catalyzed the move, and MedSAM~\cite{ma2024medsam} adapted it by fine-tuning on more than a million medical image--mask pairs. 
Because anatomy is volumetric, the decisive step is into 3D: SAM-Med3D~\cite{wang2023sammed3d} replaces the slice-wise encoder with a fully three-dimensional one, and SegVol~\cite{du2024segvol} supports both spatial and textual prompts over hundreds of anatomical categories. 
Much of this work is built on the open MONAI framework~\cite{cardoso2022monai}, whose medical-aware 3D transforms and model zoo have become standard infrastructure.

\textbf{Detection and landmark localization.} Detection localizes a finding in the body frame, and landmark localization fixes canonical anatomical coordinates; both are spatial tasks distinct from whole-image classification. 
Lung-nodule detection is anchored by the LIDC-IDRI reference database~\cite{armato2011lidcidri} and the LUNA16 benchmark~\cite{setio2017luna16}, which standardized evaluation on the free-response operating characteristic (sensitivity against false positives per scan). 
Anatomical landmarks can be found by treating localization as a sequential spatial search, as in the reinforcement-learning agent of Ghesu \etal~\cite{ghesu2016landmark} that navigates the volume to a target.

\textbf{Registration as spatial alignment.} Registration brings two images into a common coordinate frame and is foundational for multimodal fusion, atlas-based labeling, longitudinal tracking, and surgical planning. 
Learning-based methods reframe it as predicting a dense deformation field in a single forward pass: VoxelMorph~\cite{balakrishnan2019voxelmorph} established the unsupervised CNN formulation, and TransMorph~\cite{chen2022transmorph} added a transformer backbone. 
These build on classical diffeomorphic and intensity-based methods, principally the symmetric diffeomorphic SyN~\cite{avants2008syn} and the intensity-based toolbox elastix~\cite{klein2010elastix}, and are measured by the target registration error between corresponding landmarks. 
Segmentation and detection are benchmarked on generalist suites such as the Medical Segmentation Decathlon~\cite{antonelli2022msd}, the brain-tumor benchmark BraTS~\cite{menze2015brats}, the abdominal multi-organ datasets AMOS~\cite{ji2022amos} and BTCV~\cite{landman2015btcv}, the kidney-tumor KiTS19~\cite{heller2021kits19}, and the cardiac ACDC~\cite{bernard2018acdc}.

\subsection{Pivotal Task-2: Anatomical Reconstruction, Modeling, and Surgical Planning}
\label{subsec:p4-sec12-1-pivotal-task-2-anatomical-reconstruction-modeling-and-surgical-planning}

The second task turns the perceived anatomy into a patient-specific 3D model that can be measured, simulated, and planned upon, bridging from pixels to clinical decisions. Figure~\ref{fig:p4_s12_1_modeling} previews this progression from a segmentation to an anatomical mesh, a patient-specific simulation, and a surgical plan.

\begin{figure}
    \centering
    \includegraphics[width=0.99\linewidth]{figures/part4/section12/medical_reconstruction.pdf}
    \caption{Anatomical reconstruction, modeling, and planning: a volumetric segmentation becomes a 3D anatomical mesh, is personalized into a medical digital twin (\eg, a cardiac electrophysiology or blood-flow model), and is used to simulate outcomes and design a surgical or radiotherapy plan. New patient observations update model personalization. Cardiac targeting and radiotherapy dose prediction are alternative planning examples; synthetic anatomy and vector overlays illustrate the workflow rather than clinical results.}
    \label{fig:p4_s12_1_modeling}
\end{figure}

\textbf{3D anatomical reconstruction.} A segmentation becomes a usable model once it is turned into geometry. 
The classical route extracts an isosurface mesh by Marching Cubes~\cite{lorensen1998marching}, optionally constrained by a statistical shape model that encodes anatomical priors as a learned distribution over landmarks, an idea founded by Active Shape Models~\cite{cootes1995asm} and surveyed for medical use by Heimann and Meinzer~\cite{heimann2009ssm}. 
Deep methods deform a template mesh directly from the volume, as in Voxel2Mesh~\cite{wickramasinghe2020voxel2mesh} and the whole-heart reconstruction of Kong \etal~\cite{kong2021meshdeformnet}. 
Implicit neural representations supply a continuous alternative, with DeepSDF~\cite{park2019deepsdf} and Occupancy Networks~\cite{mescheder2019occupancy} as the substrate (Chapter~\ref{sec:generation}) and medical adaptations such as the neural attenuation fields of NAF~\cite{zha2022naf} reconstructing volumes from sparse-view cone-beam CT.

\textbf{Patient-specific modeling and the medical digital twin.} A geometric model becomes a \emph{twin} when it is endowed with patient-specific physiology and kept faithful to a particular patient, a medical specialization of the generic digital twin of §\ref{sec:p4_s11_1}. 
What is clinically specific is that the model is built from the patient's own imaging and physiological signals, governed by mechanistic physiology rather than CAD, and bound by regulatory credibility. 
Cardiac twins are the most mature instance: the vision was set out by Corral-Acero \etal~\cite{corralacero2020digital}, the mechanistic backbone is the multiscale cardiac modeling reviewed by Niederer \etal~\cite{niederer2019cardiacmodels}, and a personalized virtual heart built from scar imaging has been used to guide ablation of ventricular tachycardia~\cite{prakosa2018personalized}. 
Hemodynamic twins compute fractional flow reserve from coronary CT by simulating blood flow~\cite{taylor2013ffrct}, and the paradigm extends to the brain, where a virtual epileptic patient estimates the epileptogenic zone to guide surgery~\cite{wang2023vep}. 
Such models also underwrite in-silico trials, in which virtual patient cohorts partially substitute for physical ones~\cite{viceconti2016avicenna}, a route that regulators have begun to accept, as in the VICTRE imaging trial~\cite{badano2018victre}. 
The broader landscape is surveyed in recent reviews of digital twins for health~\cite{katsoulakis2024digital}.

\textbf{Surgical and treatment planning.} The patient model is finally used to plan. 
Open platforms such as 3D Slicer~\cite{fedorov2012slicer} assemble segmentation, registration, and 3D modeling into a planning workflow, and real-time finite-element engines such as SOFA~\cite{faure2012sofa} support interactive surgical simulation and rehearsal on deformable tissue. 
In radiotherapy, planning is increasingly cast as a learning problem of predicting an achievable dose distribution from anatomy, standardized by the OpenKBP challenge~\cite{babier2021openkbp}. 
Across these settings the plan is a spatial object, a trajectory, a dose map, or a resection boundary, defined on the patient's own anatomy.

\subsection{Pivotal Task-3: Image-Guided Intervention and Surgical Navigation}
\label{subsec:p4-sec12-1-pivotal-task-3-image-guided-intervention-and-surgical-navigation}

The third task brings the plan into the operating room and acts within the patient's live coordinate frame. 
The central problem is patient-to-image registration under real-world deformation and latency: the preoperative model must be aligned to the physical patient, and instruments tracked within that frame, in real time. 
Figure~\ref{fig:p4_s12_1_navigation} previews registered preoperative anatomy and augmented-reality guidance alongside the complementary reconstruction of a deforming endoscopic scene.

\begin{figure}[!ht]
    \centering
    \includegraphics[width=0.98\linewidth]{figures/part4/section12/medical_navigation.pdf}
    \caption{Patient-frame guidance combines registered anatomy and a calibrated tracked-tool pose. New intraoperative observations update anatomical registration; display calibration supports AR presentation. A separate endoscopic branch reconstructs dynamic geometry and registers it to the patient before optional fusion. The robot controller consumes the fused state rather than raw endoscopic imagery. Dashed connections denote optional downstream use. Anatomical scenes, overlays, and silicone-phantom reconstruction are schematic; transforms act on homogeneous coordinates.}
    \label{fig:p4_s12_1_navigation}
\end{figure}

\textbf{Surgical navigation.} Navigation aligns the preoperative model to the patient and tracks instruments within that frame, typically by optical tracking, which is sub-millimeter but requires line of sight, or electromagnetic tracking, which tolerates occlusion but is perturbed by metal. 
The defining difficulty is deformation: in neuronavigation, the brain shifts by up to centimeters once the skull is opened, invalidating the rigid preoperative registration, which is why intraoperative imaging and learning-based compensation are active concerns.

\textbf{Augmented-reality guidance.} Rather than divert the surgeon's gaze to a screen, augmented reality overlays the segmented anatomy directly on the surgical field, a use of the XR runtime treated in §\ref{sec:p4_s10_1}. 
The approach is most mature in spine surgery, where the Augmedics xvision system, the first FDA-cleared augmented-reality navigation platform, projects the patient's CT-derived spine onto a transparent head-mounted display for pedicle-screw placement~\cite{molina2021xvision}.

\textbf{Endoscopic spatial perception.} Minimally invasive surgery presents a hard perception problem: a monocular, moving camera views a deforming, specular, poorly lit scene. 
Surgical workflow and instruments are recognized from video, as standardized by the Cholec80 dataset and the EndoNet architecture for joint tool and phase recognition~\cite{twinanda2017endonet}. 
Geometry is recovered by self-supervised depth and visual odometry, benchmarked by the SCARED challenge~\cite{allan2021scared} and the EndoSLAM dataset~\cite{ozyoruk2021endoslam}, and the deforming surgical scene is reconstructed in 3D by dynamic neural rendering, as in EndoNeRF~\cite{wang2022endonerf}, a clinical specialization of the neural reconstruction of Chapter~\ref{sec:generation}.

\textbf{Surgical robotics.} Robotic systems execute the plan within the patient frame; their manipulation and control are the subject of §\ref{sec:p4_s9_1}, and the concern here is the spatial perception and autonomy they enable. 
The da Vinci Research Kit~\cite{kazanzides2014dvrk} is the standard open research platform, and autonomy is commonly graded on a six-level scale from teleoperation to full autonomy~\cite{yang2017autonomy}. 
Most deployed systems remain teleoperated, but the frontier of autonomous spatial action is exemplified by the Smart Tissue Autonomous Robot, which performs supervised-autonomous soft-tissue suturing using 3D tissue tracking~\cite{shademan2016star} and has been extended to in-vivo intestinal anastomosis~\cite{saeidi2022star}.

\begin{table}[!th]
\centering
\fontsize{8}{10}\selectfont
\setlength{\tabcolsep}{1.0mm}
\caption{Representative benchmarks for spatial intelligence in medicine, spanning the three pivotal tasks.}
\label{tab:p4_s12_1_benchmarks}
\begin{tabular}{L{0.16\textwidth} L{0.28\textwidth} L{0.28\textwidth} L{0.20\textwidth}}
\toprule
\textbf{Pivotal task} & \textbf{Benchmark / dataset} & \textbf{Representative method} & \textbf{Headline metric} \\
\midrule
Segmentation & MSD~\cite{antonelli2022msd} & nnU-Net~\cite{isensee2021nnunet} & Dice / HD95 \\
Segmentation & BraTS~\cite{menze2015brats} & Swin-UNETR~\cite{hatamizadeh2022swinunetr} & Dice \\
Multi-organ & AMOS~\cite{ji2022amos} / BTCV~\cite{landman2015btcv} & UNETR~\cite{hatamizadeh2022unetr} & Dice \\
Detection & LUNA16~\cite{setio2017luna16} & 3D CNN detectors (challenge entries)~\cite{setio2017luna16} & FROC sensitivity \\
Registration & multi-site brain MRI (OASIS, ADNI, ABIDE) & VoxelMorph~\cite{balakrishnan2019voxelmorph} & TRE / Dice \\
Reconstruction & MM-WHS~\cite{zhuang2019mmwhs} & whole-heart mesh~\cite{kong2021meshdeformnet} & surface distance \\
Modeling & patient-specific cardiac & virtual heart~\cite{prakosa2018personalized} & ablation-target agreement \\
Endoscopy & SCARED~\cite{allan2021scared} & EndoNeRF~\cite{wang2022endonerf} & depth error \\
Workflow & Cholec80~\cite{twinanda2017endonet} & EndoNet~\cite{twinanda2017endonet} & phase accuracy \\
\bottomrule
\end{tabular}
\end{table}

\subsection{Established Tools and Systems}
\label{subsec:p4-sec12-1-established-tools-and-systems}

\textbf{Open-source frameworks.} The research stack runs from voxels to a planned model: \textbf{3D Slicer}\footnote{\url{https://www.slicer.org/}} and \textbf{MITK}\footnote{\url{https://www.mitk.org/}} provide segmentation, registration, and planning environments over the lower-level \textbf{ITK} and \textbf{SimpleITK}\footnote{\url{https://itk.org/}} engines; \textbf{MONAI}\footnote{\url{https://monai.io/}} supplies 3D-native deep-learning components and \textbf{nnU-Net}\footnote{\url{https://github.com/MIC-DKFZ/nnUNet}} the standard segmentation pipeline; and the web viewer \textbf{OHIF}\footnote{\url{https://ohif.org/}} renders and annotates volumes in the browser.

\textbf{Navigation and surgical-robotics vendors.} The clinical stack carries the model from preoperative registration to constrained in-OR action. Navigation systems from \textbf{Brainlab}\footnote{\url{https://www.brainlab.com/}}, \textbf{Medtronic} (StealthStation)\footnote{\url{https://www.medtronic.com/}}, and \textbf{Stryker}\footnote{\url{https://www.stryker.com/}} register and track instruments, and \textbf{Augmedics}\footnote{\url{https://augmedics.com/}} adds an augmented-reality head-mounted overlay. Surgical robots span teleoperated soft-tissue systems (\textbf{Intuitive} da Vinci and Ion\footnote{\url{https://www.intuitive.com/}}, \textbf{Medtronic} Hugo, \textbf{CMR Surgical} Versius\footnote{\url{https://cmrsurgical.com/}}) and CT-planned, haptically constrained orthopedic systems (\textbf{Stryker} Mako, \textbf{Zimmer Biomet} ROSA\footnote{\url{https://www.zimmerbiomet.com/}}).

\textbf{Imaging-AI and clinical platforms.} A deployed layer turns spatial findings into care. \textbf{Aidoc}\footnote{\url{https://www.aidoc.com/}} triages critical findings on CT, \textbf{Viz.ai}\footnote{\url{https://www.viz.ai/}} alerts stroke teams to large-vessel occlusions, \textbf{HeartFlow}\footnote{\url{https://www.heartflow.com/}} computes non-invasive fractional flow reserve from coronary CT, \textbf{PathAI}\footnote{\url{https://www.pathai.com/}} quantifies biomarkers on whole-slide pathology, and \textbf{Tempus}\footnote{\url{https://www.tempus.com/}} fuses imaging with genomic and clinical data.

\subsection{Challenges}
\label{subsec:p4-sec12-1-challenges}

Every capability in the clinical loop must be accurate, robust to deformation, and trustworthy enough to act on a patient. The open problems follow from those demands: generalizing across scanners and sites, obtaining labels under privacy constraints, registering deforming tissue during surgery, earning regulatory and clinical trust, and meeting real-time budgets while fusing imaging with non-imaging evidence.

\paragraph{Generalization across Scanners and Sites.} Models trained at one institution degrade at another, because scanner vendor, reconstruction kernel, and acquisition protocol induce a covariate shift that curated benchmarks underrepresent. 
Domain generalization and adaptation are therefore prerequisites for deployment, and leaderboard results routinely overstate performance at a new site.

\paragraph{Data Scarcity, Annotation Cost, and Privacy.} Volumetric expert annotation is slow and expensive, and patient data is locked in privacy-restricted silos. 
Self-supervised pretraining reduces the labeling burden~\cite{huang2023ssl}, and federated learning trains shared models without moving data across institutions~\cite{rieke2020federated}, but label efficiency and privacy-preserving collaboration remain limiting.

\paragraph{Non-Rigid Intraoperative Registration.} The hardest open spatial problem is keeping the model aligned with the patient as tissue moves. Brain shift after craniotomy and soft-tissue deformation during surgery invalidate rigid preoperative registration, and while learning-based deformable registration~\cite{balakrishnan2019voxelmorph} addresses the static case, real-time non-rigid intraoperative alignment under sparse, partial observation is not solved.

\paragraph{Validation, Regulation, and Trust.} A model that changes a clinical decision needs prospective validation and regulatory clearance before use, and at the bedside it needs calibrated uncertainty and explanations that support appropriate rather than blind reliance. Retrospective benchmark accuracy establishes none of these.

\paragraph{Real-Time Constraints and Multimodal Fusion.} Intraoperative guidance must run within a tight latency budget while the scene deforms, a constraint that few surgical-video algorithms yet meet without specialized hardware. 
At the same time, the richest clinical picture comes from fusing imaging with electronic health records and genomics~\cite{acosta2022multimodal}, so spatial perception must ultimately integrate with non-spatial evidence.


\newpage
\section{Natural Science}
\label{sec:Natural Science}
\label{sec:p4_s12_2}


\subsection{Introduction}
\label{subsec:p4-sec12-2-introduction}

A physical system is a field over space; a molecule or a material is a three-dimensional arrangement of atoms together with an electronic field over that arrangement; a protein, a cell, or a tissue is a three-dimensional structure whose function follows from its form.
In every case the scientific quantity of interest, whether an energy, a reaction rate, a biological function, or a material property, is a function of spatial structure under strict symmetry constraints, so that understanding nature is to a large degree a problem of recovering and reasoning about that structure.

The geometric representations of Chapter~\ref{chapter:single_modality} and the reconstruction and generation of Chapter~\ref{sec:generation} reappear here specialized to data that obey physical law: to fields constrained by differential equations, to atomic systems constrained by Euclidean and permutation symmetry, and to biomolecules whose geometry determines their function.
A single capability triad recurs inside every domain and gives the section its spine: spatial \textbf{representation and reconstruction} recovers a continuous field or a 3D structure from sparse, indirect, and noisy measurements; \textbf{dynamics, simulation, and law discovery} learns how that structure evolves and what governs it; and \textbf{design, discovery, and control} either generates new structure or inverts observations for hidden causes. The third leg reads as discovery and control in physics and as generative design in chemistry and life science.

The section is organized by scientific domain (physics, chemistry, and life science); the earth and environmental sciences are treated at geographic scale in the GeoAI of §\ref{sec:p4_s11_2}, the clinical, patient-scale counterpart of the life-science material in §\ref{sec:p4_s12_1}, and the generic digital twin in §\ref{sec:p4_s11_1}.
Figure~\ref{fig:p4_s12_2_overview} previews the section as a cross-scale map in which the three domains share the same vertical triad, and Table~\ref{tab:p4_s12_2_benchmarks} collects a representative benchmark for each capability.

\begin{figure}[!ht]
    \centering
    \includegraphics[width=0.98\linewidth]{figures/part4/section12/natural_science_overview_v2.pdf}
    \caption{Spatial intelligence across physics, chemistry, and life science (columns), each repeating the same capability triad (rows). The lower electron-to-tissue strip indicates illustrative scales rather than exclusive domain boundaries. Fields, molecular and periodic structures, and biomolecular and tissue structures illustrate representation; the equations exemplify wave, quantum, and biological dynamics. Bottom-row design examples are schematic; molecule/material and biomolecular/tissue examples are alternatives rather than literal conversion pipelines. Biological imagery is synthetic.}
    \label{fig:p4_s12_2_overview}
\end{figure}


\subsection{Physics}
\label{subsec:p4-sec12-2-physics}

Spatial intelligence in physics constructs spatial structure and discovers the laws that govern its motion, and it earns its place where first-principles solvers are too costly or the governing equations are unknown: fields are reconstructed from limited observations, their dynamics are modeled and their laws discovered, and the systems are optimized and controlled.
Figure~\ref{fig:p4_s12_2_physics} previews this progression from a sparsely observed field to a learned dynamical rollout and a differentiable-control loop.

\begin{figure}
    \centering
    \includegraphics[width=0.98\linewidth]{figures/part4/section12/natural_science_physics_v2.pdf}
    \caption{The physics triad: (1) reconstruct a field from sparse sensors, (2) fit dynamics and roll out the field, with a separate law-extraction step, and (3) optimize controls through a differentiable simulator. $\mathcal{N}[\mathbf{u};\theta]$ is the learned operator and $\mathcal{N}^*$ a candidate extracted law. The state $\mathbf u$ and control $\mathbf a$ are distinct; the control objective $\mathcal J$ supplies gradients to the optimizer, whereas $\mathcal L_{\rm fit}$ fits observed dynamics. Field contours and temporal shifts are illustrative, not experimental results.}
    \label{fig:p4_s12_2_physics}
\end{figure}

\subsubsection{Spatial Representation and Reconstruction}

Many key quantities in physical systems are available only through sparse and noisy observations, and representation and reconstruction recover continuous spatial fields from such data, the input to law discovery and control.
Embedding physical priors such as governing partial differential equations and symmetry constraints keeps these reconstructions physically consistent in low-data regimes.

A classical task is physical field reconstruction, which aims to recover continuous physical fields such as temperature, pressure, and velocity from sparse sensor readings or boundary measurements; the central challenge lies in ensuring the physical plausibility of the solution under the ill-posed conditions caused by insufficient observations.
\textbf{Physics-informed neural networks} are a classical approach for this task~\cite{raissi2019physics}, embedding partial differential equations such as the Navier--Stokes equations as soft constraints in the loss function and computing physical residuals via automatic differentiation, which enables accurate reconstruction of fluid and thermal fields in data-scarce settings; the same networks serve as real-time surrogates in digital twins (§\ref{sec:p4_s11_1}).
\textbf{Neural fields} implicitly represent continuous physical fields by using coordinate-based networks that map spatial locations to physical quantities.
Sitzmann \etal~\cite{sitzmann2020implicit} proposed SIREN, which employs periodic activation functions to capture high-frequency spatial details with high fidelity and has been applied to solving Poisson equations and reconstructing wave fields.
The general machinery of coordinate-based neural representation is developed in Chapter~\ref{sec:generation}.

Inverse scattering aims to infer the internal structure or material parameters of an object from measured wave-scattering data, and it remains a canonical ill-posed problem in electromagnetics and acoustics~\cite{chen2020review}.
Wei and Chen proposed deep-learning schemes~\cite{wei2019deep} for full-wave nonlinear inverse scattering that directly invert dielectric-permittivity profiles from scattered fields, thereby moving beyond linearized Born-type approximations in practical reconstructions.
SwitchNet~\cite{khoo2019switchnet} introduced a unified neural architecture that learns both the forward map from scatterers to scattered fields and the inverse map from measurements to scatterer profiles, using a switching layer that exploits low-rank structure to reduce model complexity and facilitate training.
Physics-guided losses that inject multiple-scattering near-field priors into the objective further improve robustness.

\subsubsection{Dynamics, Simulation, and Law Discovery}

\textbf{Law discovery} aims to learn the generative structures or explicit laws that govern the evolution of physical systems from observational data, emphasizing the automatic extraction of conservation laws, symmetries, and interpretable mechanisms beyond accurate prediction.
Hamiltonian Neural Networks~\cite{greydanus2019hamiltonian} directly learn the Hamiltonian of a system, thereby enforcing energy conservation and symplectic structure at the model level and providing a principled foundation for interpretable dynamical modeling.
Lagrangian Neural Networks~\cite{cranmer2020lagrangian} parameterize arbitrary Lagrangians with neural networks, learning energy conserving dynamics without canonical coordinates, as validated on the double pendulum and a relativistic particle.
Cranmer \etal~\cite{cranmer2020discovering} combine graph networks with symbolic regression on their internal representations, recovering known force laws and Hamiltonians, and applying the approach to cosmology.
The SINDy algorithm~\cite{brunton2016discovering} employs sparse regression to identify governing differential equations directly from time-series data, enabling data-driven discovery of interpretable physical laws.

\textbf{Efficient physical solvers} focus on predicting the spatio-temporal evolution of physical systems under known governing equations or interaction rules, with the goal of achieving fast and scalable computation while preserving physical consistency beyond the limits of traditional numerical methods.
The Fourier Neural Operator~\cite{li2021fourier} learns mappings between function spaces rather than individual discretized instances, enabling trained models to predict the evolution of partial differential equations at arbitrary spatial resolutions without relying on a fixed computational grid.
DeepONet~\cite{lu2021learning}, grounded in the universal approximation theorem for operators, learns nonlinear mappings from input functions to output functions and thereby approximates the solution operators that govern physical systems.
Interaction Networks~\cite{battaglia2016interaction} introduced the relational reasoning architecture that learns physical interactions between objects, laying the foundation for graph-based physical simulation.
The Graph Network Simulator~\cite{sanchez2020learning} employs graph neural networks to model spatial interactions among particles or objects, achieving high-fidelity simulations of fluids, rigid bodies, and deformable materials, and MeshGraphNets~\cite{pfaff2020learning} extend graph-based models to unstructured meshes, allowing efficient simulation of continuum mechanics on irregular geometries.
For the classical challenge of turbulence modeling, Mohan \etal~\cite{mohan2020spatio} proposed spatio-temporal deep-learning models with physics-informed diagnostics that improve the prediction accuracy of three-dimensional turbulent flows, and Lusch \etal~\cite{lusch2018deep} developed a deep-learning framework for universal linear embeddings of nonlinear dynamics, enabling low-dimensional representations and long-term prediction of nonlinear fluid systems.
Neural representations of the many-body wavefunction are treated with electronic structure in the chemistry subsection.

\subsubsection{Control and Optimization}

Optimization and control in physical systems seek optimal trajectories, shapes, or control policies under governing physical laws and environmental constraints.
Differentiable physics formulates physical simulators as differentiable operators, enabling gradient-based optimization over control inputs, design variables, and physical parameters, and providing a unified framework for optimal control, system identification, and design.
Neural Ordinary Differential Equations~\cite{chen2018neural} provide a differentiable formulation of continuous-time dynamical systems by parameterizing the vector field with a neural network and enabling gradient computation through adjoint sensitivity analysis, serving as a flexible building block for learning and control in continuous time.
ChainQueen~\cite{hu2019chainqueen} demonstrates differentiable simulation of soft and deformable bodies based on the moving least squares material point method, enabling gradient-based optimization in high-degree-of-freedom systems with contact and supporting applications in soft-robot control and co-design.
DiffTaichi~\cite{hu2020difftaichi} generalizes this paradigm through a differentiable programming framework for high-performance physical simulation, allowing efficient construction of differentiable solvers and gradient-driven optimization for fluid control, shape optimization, and material-parameter estimation; the same tools calibrate digital twins (§\ref{sec:p4_s11_1}).
In practice, a mature software substrate of differentiable and physics-informed libraries, such as \textbf{DeepXDE}\footnote{\url{https://github.com/lululxvi/deepxde}} and \textbf{NVIDIA Modulus}\footnote{\url{https://developer.nvidia.com/modulus}} for operator and PINN training, and \textbf{JAX-MD}\footnote{\url{https://github.com/jax-md/jax-md}} and \textbf{DeePMD-kit}\footnote{\url{https://github.com/deepmodeling/deepmd-kit}} for differentiable and learned-potential molecular dynamics, turns this triad into runnable physics.




\subsection{Chemistry}
\label{subsec:p4-sec12-2-chemistry}

Chemistry is intrinsically spatial: a molecule or a material is a three-dimensional arrangement of atoms together with an electronic field over it, and its behavior, whether energy, reactivity, or a bulk property, is a function of that geometry under strict symmetry.
Physical observables must be invariant or equivariant to the Euclidean group of rotations, translations, and reflections and to permutations of identical atoms; for crystalline solids they must also respect periodic lattice symmetry, and for electrons they must obey antisymmetry under exchange.
Spatial intelligence in chemistry is therefore largely the art of building these symmetries into models that represent atomic and electronic structure, simulate its dynamics, and generate new structure under chemical feasibility.
Figure~\ref{fig:p4_s12_2_chemistry} previews this triad from an equivariant representation of a 3D structure to a learned-potential simulation and a 3D generative design.

\begin{figure}
    \centering
    \includegraphics[width=0.98\linewidth]{figures/part4/section12/natural_science_chemistry_v2.pdf}
    \caption{The chemistry triad: (1) symmetry-aware molecular and periodic representations, illustrated by E(3)-equivariant features; (2) a learned invariant interatomic energy $U$ whose coordinate gradients provide forces for molecular dynamics, with related representations supporting property prediction; and (3) complementary 3D generative design by reverse diffusion. Molecular graphs, force directions, and the one-dimensional energy profile are schematic, not a calculated trajectory or energy surface. The denoising example preserves the atom count; generative methods may model coordinates and atom types jointly.}
    \label{fig:p4_s12_2_chemistry}
\end{figure}

\subsubsection{Spatial Representation and Reconstruction}

The central representational problem of chemistry is to map a 3D atomic geometry to energies, forces, and properties while respecting symmetry, and successive representations have captured progressively more of that geometry.
\textbf{Geometry-grounded potentials.} The Behler--Parrinello neural network potential founded the area by decomposing the total energy into atomic contributions computed from rotation-, translation-, and permutation-invariant symmetry functions of each atom's local environment~\cite{behler2007generalized}.
SchNet~\cite{schutt2017schnet} replaced hand-designed descriptors with continuous-filter convolutions over interatomic distances, giving a smooth invariant representation for energy and force prediction, while DimeNet~\cite{klicpera2020directional} added directional message passing that encodes bond angles through spherical harmonics, moving beyond pairwise distances to angular structure, and PaiNN~\cite{schutt2021equivariant} extended message passing to rotationally equivariant vector features for tensorial properties and spectra.
The Deep Potential model~\cite{zhang2018deep} preserves all natural symmetries at linear cost, and the ANI potentials~\cite{smith2017ani} demonstrated transferable organic-chemistry accuracy at force-field cost.
\textbf{Equivariant interatomic potentials.} The decisive recent step was to carry geometric tensors, not merely invariant scalars, through the network.
NequIP~\cite{batzner2022e3} introduced E(3)-equivariant convolutions over geometric tensors and achieved striking data efficiency, often requiring orders of magnitude less training data; Allegro~\cite{musaelian2023learning} made the representation strictly local and thereby scalable to 100 million atoms, and MACE~\cite{batatia2022mace} used higher body-order equivariant messages to reach high accuracy in only two message-passing layers.
A further shift is from per-system fits to \textbf{universal potentials} that span the periodic table: M3GNet~\cite{chen2022universal} couples a graph network with explicit three-body interactions trained on materials relaxations, and CHGNet~\cite{deng2023chgnet} augments the 3D-structure graph with atomic magnetic moments to capture electronic degrees of freedom alongside geometry. These equivariant backbones share machinery with the geometric representations of Chapter~\ref{chapter:single_modality}.
\textbf{Neural electronic structure.} At the quantum scale the spatial object is the many-electron wavefunction itself.
Neural quantum states parameterize that wavefunction with a neural network over spin configurations~\cite{carleo2017solving}, and deep variational ans\"atze such as FermiNet~\cite{pfau2020ferminet} and PauliNet~\cite{hermann2020deep} operate directly on 3D electron coordinates while enforcing exact antisymmetry, solving the electronic Schr\"odinger equation to chemical accuracy from nuclear positions alone.

\subsubsection{Dynamics, Simulation, and Property Prediction}

The potentials of the previous thread are also the simulators: once a model maps a 3D configuration to energies and forces, it drives molecular dynamics at a fraction of the cost of first-principles calculation, and the same representations predict quantum properties such as energies, dipoles, and frontier-orbital gaps directly from coordinates.
The gradient-domain machine-learning force field made this explicit by exploiting energy conservation and the geometry of conformations, and it established the widely used MD17 trajectory benchmark of energies and forces for small molecules~\cite{chmiela2017machine}.
Heterogeneous catalysis pushes the same energy-and-force prediction task to a far larger and more diverse chemical space, spanning many combinations of metal surfaces and adsorbates. 
GemNet-OC~\cite{gasteiger2022gemnet} extends directional message passing to this setting, delivering accurate energies and forces across the surfaces and adsorbates of the Open Catalyst datasets.
In parallel, structure-to-property prediction is standardized on QM9~\cite{ramakrishnan2014quantum}, a dataset of roughly 134{,}000 small organic molecules with density-functional geometries and quantum properties, on which the equivariant networks above are routinely evaluated by the mean absolute error of predicted properties.
This ecosystem rests on open toolkits, including the equivariant-network library \textbf{e3nn}\footnote{\url{https://e3nn.org/}}, the atomistic-simulation environment \textbf{ASE}\footnote{\url{https://wiki.fysik.dtu.dk/ase/}}, and the quantum-chemistry package \textbf{PySCF}\footnote{\url{https://pyscf.org/}}, together with universal-potential foundation models such as MACE-MP, and on reference corpora such as the Materials Project~\cite{jain2013commentary}, the Open Catalyst dataset~\cite{chanussot2021open}, and the MoleculeNet benchmark suite~\cite{wu2018moleculenet}.

\subsubsection{Generation and Inverse Design}

The inverse of property prediction is design: generating new 3D structures that satisfy desired properties under chemical feasibility, by sampling on the appropriate geometric manifold rather than enumerating candidates.
G-SchNet~\cite{gebauer2019symmetry} autoregressively places atoms in 3D Euclidean space with each placement respecting the rotational invariance and local symmetry of those already placed, the first strong 3D-native generator.
Diffusion models then became the dominant paradigm: the equivariant diffusion model~\cite{hoogeboom2022equivariant} jointly denoises atom coordinates and types under E(3) equivariance, GeoDiff~\cite{xu2022geodiff} generates conformers from the molecular graph through roto-translation-equivariant Markov kernels that induce an invariant distribution, and torsional diffusion~\cite{jing2022torsional} restricts the process to the hypertorus of torsion angles, the true intrinsic conformational degrees of freedom, achieving state-of-the-art conformer quality on GEOM-Drugs~\cite{axelrod2022geom} at a fraction of the cost.
For crystalline materials, CDVAE~\cite{xie2022crystal} generates stable structures under periodicity and SE(3) invariance and supports inverse design by optimizing target properties in latent space.
The line has since moved to diffusion directly on the periodic structure: DiffCSP~\cite{jiao2023crystal} jointly denoises lattice and fractional coordinates under periodic equivariance for crystal structure prediction, and MatterGen~\cite{zeni2025generative} scales this into a property-conditioned foundation model for inorganic materials design, with generated candidates validated by experimental synthesis.
The general diffusion machinery underlying these models is developed in Chapter~\ref{sec:generation}; what is chemistry-specific is the constraint manifold (valence and bonding feasibility, energy-minimum geometries, lattice periodicity) on which generation must take place.


\subsection{Life Science}
\label{subsec:p4-sec12-2-life-science}

At the molecular, cellular, and tissue scale of basic biology the object of interest is again a three-dimensional structure: a protein fold, a molecular complex, a spatial map of gene expression, or the wiring of a neural circuit, and in each case function follows from spatial form.
Spatial intelligence here recovers that structure from indirect measurement, models how biomolecules interact and move, and designs new structure.
Figure~\ref{fig:p4_s12_2_lifescience} previews the triad from structure recovery through interaction modeling to generative design.

\begin{figure}
    \centering
    \includegraphics[width=0.98\linewidth]{figures/part4/section12/natural_science_lifescience_v2.pdf}
    \caption{The life-science triad: (1) complementary sequence-to-structure prediction and cryo-EM projection-to-density reconstruction, (2) ligand binding-pose geometry with translation, rotation, and bond torsion, and (3) de novo backbone generation followed by sequence design. These are complementary capabilities, not an obligatory serial pipeline. Protein, pocket, and density renderings, noisy projection examples, and sequence fragments are schematic; they do not constitute experimentally validated structures or designs.}
    \label{fig:p4_s12_2_lifescience}
\end{figure}

\subsubsection{Spatial Representation and Reconstruction}

The defining achievement of spatial intelligence in biology is the recovery of 3D structure from indirect signals.
\textbf{Protein structure prediction} maps a one-dimensional sequence to a three-dimensional fold.
The original AlphaFold predicted inter-residue distance distributions as a learned potential to fold the chain~\cite{senior2020improved}, and AlphaFold2~\cite{jumper2021highly} reached near-experimental accuracy through an end-to-end architecture that reasons explicitly over residue frames and pairwise geometry, the exemplar of spatial intelligence in the natural sciences.
AlphaFold3~\cite{abramson2024accurate} generalized this to the joint structure of proteins, nucleic acids, and small molecules through a diffusion decoder; RoseTTAFold~\cite{baek2021accurate} let one-dimensional, two-dimensional, and three-dimensional information flow between tracks, and ESMFold~\cite{lin2023evolutionary} showed that atomic-level structure emerges from a protein language model without explicit alignments, enabling metagenomic-scale folding.
\textbf{Cryo-electron microscopy} solves the inverse problem of reconstructing a 3D density map from many noisy 2D projections.
RELION cast this as Bayesian refinement~\cite{scheres2012relion}, CryoSPARC enabled rapid ab-initio reconstruction without a starting model~\cite{punjani2017cryosparc}, and cryoDRGN used a coordinate-based neural decoder to reconstruct continuous distributions of conformations and thereby resolve molecular heterogeneity and motion~\cite{zhong2021cryodrgn}.
\textbf{Spatial transcriptomics} places gene expression back into tissue space.
The founding barcoded-array assay preserved the positional information of the transcriptome in tissue sections~\cite{stahl2016visualization}; Tangram learns a spatial mapping that aligns dissociated single cells onto a spatial reference~\cite{biancalani2021deep}, and novoSpaRc reconstructs spatial arrangement from expression alone through optimal transport, assuming that expression-similar cells lie physically close~\cite{nitzan2019gene}.
\textbf{Connectomics} reconstructs neural wiring from electron-microscopy volumes, where flood-filling networks iteratively trace the 3D morphology of individual neurites through the volume~\cite{januszewski2018high}.

\subsubsection{Dynamics, Simulation, and Interaction}

Beyond static structure, spatial intelligence models how biomolecules bind and interact.
Molecular docking predicts the 3D pose in which a ligand binds a protein, and recent methods treat it as a geometric problem on the natural manifold of poses.
DiffDock~\cite{corso2023diffdock} formulates docking as diffusion over the product space of ligand translation, rotation, and torsion, sampling binding poses in their intrinsic degrees of freedom, while EquiBind~\cite{stark2022equibind} regresses the binding location and pose in a single shot using an SE(3)-equivariant network whose equivariance encodes the rotational and translational symmetry of the binding geometry.
Protein--protein and complex geometry is predicted by extending folding models to multiple chains~\cite{evans2022protein}, and AlphaFold3 above provides a generative alternative for protein--ligand and protein--nucleic-acid complexes.
The dynamics of biomolecules in turn rest on the learned interatomic potentials of the chemistry subsection, which accelerate the molecular dynamics that sample conformational change.

\subsubsection{Generation and Design}

The generative leg designs new 3D biological structure.
RFdiffusion~\cite{watson2023novo} adapts structure-denoising diffusion to generate novel protein backbones conditioned on functional and geometric constraints such as motif scaffolding, binders, and symmetric assemblies, and ProteinMPNN~\cite{dauparas2022robust} solves the inverse-folding problem of designing an amino-acid sequence that folds into a given 3D backbone, the standard partner to backbone generation.
Chroma~\cite{ingraham2023illuminating} couples equivariant graph networks with a diffusion process to sample novel all-atom structures and sequences steerable toward desired properties, illuminating regions of protein space never explored by evolution.
As in chemistry, the underlying diffusion machinery is that of Chapter~\ref{sec:generation}; what is biology-specific is that the generated geometry must be foldable and functional, which is why design is paired with structure prediction and, ultimately, wet-lab validation.
The field is supported by open infrastructure, including the AlphaFold Protein Structure Database and AlphaFold3 server\footnote{\url{https://alphafoldserver.com/}}, the ESM Metagenomic Atlas\footnote{\url{https://esmatlas.com/}}, and the open-source folding stack OpenFold\footnote{\url{https://github.com/aqlaboratory/openfold}}, together with reference corpora such as the Protein Data Bank, the CASP blind-assessment benchmark~\cite{kryshtafovych2021critical}, and the PDBBind docking benchmark~\cite{wang2004pdbbind}.

\begin{table}[!th]
\centering
\fontsize{8}{10}\selectfont
\setlength{\tabcolsep}{1.0mm}
\caption{Representative benchmarks for spatial intelligence across the three natural-science domains.}
\label{tab:p4_s12_2_benchmarks}
\begin{tabular}{L{0.11\textwidth} L{0.17\textwidth} L{0.24\textwidth} L{0.19\textwidth} L{0.21\textwidth}}
\toprule
\textbf{Domain} & \textbf{Capability} & \textbf{Benchmark / dataset} & \textbf{Representative method} & \textbf{Headline metric} \\
\midrule
Physics & PDE surrogate & Navier--Stokes / Darcy flow & FNO~\cite{li2021fourier} & relative $L_2$ error \\
Physics & Particle dynamics & Learning-to-Simulate suite (Water / Sand / Goop) & GNS~\cite{sanchez2020learning} & rollout MSE \\
Chemistry & Property prediction & QM9~\cite{ramakrishnan2014quantum} & DimeNet~\cite{klicpera2020directional} & property MAE \\
Chemistry & Force field & MD17~\cite{chmiela2017machine} & NequIP~\cite{batzner2022e3} & energy / force MAE \\
Chemistry & Catalysis & Open Catalyst~\cite{chanussot2021open} & GemNet-OC~\cite{gasteiger2022gemnet} & energy / force MAE \\
Chemistry & Conformer generation & GEOM-Drugs~\cite{axelrod2022geom} & torsional diffusion~\cite{jing2022torsional} & coverage / matching RMSD \\
Life science & Protein folding & CASP~\cite{kryshtafovych2021critical} & AlphaFold2~\cite{jumper2021highly} & GDT-TS \\
Life science & Docking & PDBBind~\cite{wang2004pdbbind} & DiffDock~\cite{corso2023diffdock} & top-1 RMSD $<2$\,\AA{} (\%) \\
Life science & Cryo-EM heterogeneity & EMPIAR single-particle datasets & cryoDRGN~\cite{zhong2021cryodrgn} & resolution (FSC) \\
Life science & Protein design & Motif-scaffolding / binder design & RFdiffusion~\cite{watson2023novo} & in-silico design success rate \\
\bottomrule
\end{tabular}
\end{table}


\subsection{Challenges}
\label{subsec:p4-sec12-2-challenges}

A learned model of nature has to respect the symmetries and conservation laws of its domain and remain correct outside the data it was fitted to. The six problems below follow from that requirement: how to build symmetry in, how to bridge scales, how to solve ill-posed inverse problems, how to learn from scarce and expensive labels, how to generalize to genuinely new systems, and how to make predictions physically consistent and trustworthy.

\paragraph{Symmetry and Equivariance.} Physical observables are constrained by symmetry: Euclidean invariance and equivariance, permutation of identical particles, antisymmetry of electrons, and periodicity of crystals. Whether to build these symmetries into the architecture or to learn them from data is the central design tension, trading the cheaper but less expressive invariant features~\cite{schutt2017schnet} against the richer but costlier equivariant tensor features~\cite{batzner2022e3}, and it recurs from interatomic potentials to protein structure.

\paragraph{Extreme Multi-Scale Structure.} Spatial structure in nature spans many orders of magnitude, from electrons and atoms to molecules and materials, and from biomolecules through cells to tissue. Learned potentials now reach a hundred million atoms~\cite{musaelian2023learning}, still far below cellular scale; no single representation natively bridges electronic-structure accuracy and mesoscale reachability, and the handoffs between scales, through coarse-graining or learned-potential dynamics, are lossy and remain only partially solved.

\paragraph{Ill-Posed Inverse Problems.} Much of natural science is an inverse problem in which a high-dimensional spatial field must be recovered from sparse, indirect, and noisy measurements: field reconstruction from a few sensors, inverse scattering~\cite{wei2019deep}, and the recovery of 3D density from 2D cryo-electron-microscopy projections~\cite{zhong2021cryodrgn} all share this under-determined character and demand strong learned priors to yield a unique and physical solution.

\paragraph{Data Scarcity and the Cost of Ground Truth.} High-fidelity labels are expensive: accurate energies and forces require costly quantum-chemistry calculations, which is why force-field benchmarks such as MD17 cover only small molecules~\cite{chmiela2017machine}, and validated 3D biological structures require slow experiments, so the space of trustworthy training data is small relative to the space of possible systems. This drives the field's emphasis on data efficiency, multi-fidelity learning, and physics-guided objectives.

\paragraph{Generalization Beyond the Training Distribution.} Models must extrapolate to unseen chemistry, novel protein folds, off-equilibrium and reactive geometries, and configurations far from those they were trained on. Discovery is by definition out of distribution, so leaderboard accuracy says little about it; this is why generative materials design is now judged by whether its candidates can actually be synthesized~\cite{zeni2025generative}.

\paragraph{Physical Consistency and Trust.} Predictions that violate conservation laws or come without calibrated uncertainty cannot support a scientific conclusion. Building conservation into the model, as Hamiltonian networks do for energy~\cite{greydanus2019hamiltonian}, and quantifying confidence are what let a learned model extend the scientific record rather than interpolate it.

\newpage
\chapter{Spatial Intelligence for Intelligent Assistance}


\FramedInitial{I}{ntelligent} Assistance is the assistant-facing application of spatial intelligence: an agent must perceive a user's lived environment, remember its changing state, ground natural instructions in that state, and act in ways that are useful, legible, and safe.
Unlike a task-specific robot in a factory, vehicle, or clinic, a personal assistant operates in homes, offices, streets, and wearable or XR interfaces where the scene is private, cluttered, and shaped by individual routines.
This chapter therefore treats assistance as a persistent human-centered loop over spatial memory, language grounding, action, and feedback.
The same spatial model supports tidying and search, close collaboration, and accessibility support for vulnerable users, but each setting changes how much the assistant must model the person as well as the place.





\section{Introduction}
\label{subsec:p4-sec13-1-introduction}\label{sec:p4_s13_1}

An intelligent assistant is a generalist embodied agent that shares the user's own private, unstructured world and acts on the user's behalf. 
Unlike the domain agents of the preceding chapters, which are optimized for a bounded operational envelope (a road network, a warehouse aisle, an operating theater), the assistant must function inside the open-ended geometry of a home, an office, or a sidewalk, in the continuous physical presence of the person it serves. 
Its defining spatial challenge is therefore not a new perception, reasoning, generation, or memory primitive: those capabilities are surveyed in Chapter~\ref{chapter:3:sec:6} (perception), Chapter~\ref{sec:reasoning} (reasoning), Chapter~\ref{sec:generation} (generation), and Section~\ref{subsec:p3_s7_3_6} (retrieval- and memory-augmented spatial reasoning).
The challenge here is rather to build and maintain a persistent, personalized, and queryable spatial-semantic model of the user's environment, to ground human intent onto that model, and to act on it legibly and safely while a person occupies the same space.

Three properties separate the assistant from the embodied agents discussed elsewhere in this part. The environment is the user's own, so the assistant cannot assume an engineered layout and must instead generalize across the long tail of real homes. 
It is also personalized: the same instruction (``tidy up,'' ``take me to the exit'') resolves to different geometric goals for different users, so the spatial model must absorb individual preference. 
And because the human is always co-present, otherwise routine motion becomes a social act that must be predictable, comfortable, and safe. These three demands (an unengineered space, an individual user, and a person sharing that space) recur across every regime in which the assistant operates.

\section{System Overview}
\label{subsec:p4-sec13-1-system-overview}

Across these settings, every assistant closes the same loop. It \textbf{perceives} the environment and the user, \textbf{represents} what it perceives in a persistent spatial-semantic memory, \textbf{grounds} a natural-language or gestural instruction onto that memory, \textbf{reasons} about a goal-directed plan, \textbf{acts} on the physical or informational world, and observes the human's \textbf{feedback}, which re-grounds the next interaction. 
The persistent memory is the part that distinguishes an assistant from a single-shot perception or manipulation system: because the assistant lives in one place over long horizons, it can accumulate a map of where objects belong, where the user tends to be, and how the user prefers tasks to be performed. 
Long-horizon mobile-manipulation benchmarks such as EMMOE instantiate one slice of this loop (language-conditioned planning and action in open environments) but still leave persistent personalization, memory update, and human feedback as broader assistant-level requirements~\cite{li2025emmoe}. Figure~\ref{fig:p4_s13_1_overview} schematizes the loop, with the spatial-semantic memory at its center and the human present at every stage.

\begin{figure}[t]
    \centering
    \includegraphics[width=0.98\linewidth]{figures/part4/section13/assistant_overview.pdf}
    \caption{The intelligent-assistant loop. User and scene observations populate persistent room--object memory; instructions are grounded to specific objects and target states before planning. Alternative outputs illustrate household management, shared-space collaboration, and audio/haptic assistance. Memory links denote spatial associations; the dashed ellipse marks shared interaction space. User responses and observed changes close the feedback loop.}
    \label{fig:p4_s13_1_overview}
\end{figure}

The assistant is deployed in three escalating regimes of human-proximity, which organize the remainder of this chapter. 
In the first, the assistant acts \emph{on} the space while the human is a beneficiary at a distance: it tidies, fetches, and searches. 
In the second, the assistant acts \emph{alongside} the human as a partner that must move and manipulate in shared space without intruding. 
In the third, the assistant acts \emph{for} a vulnerable human, standing in for an impaired capability under hard safety, privacy, and accessibility constraints. 
The three regimes share the loop above but place progressively more weight on modeling the person rather than the place; Table~\ref{tab:p4_s13_1_benchmarks} collects benchmarks for all three.

\section{Pivotal Task-1: Managing the Physical Environment}
\label{subsec:p4-sec13-1-pivotal-task-1-managing-the-physical-environment}

The most autonomous regime asks the assistant to maintain the physical state of a human environment with little real-time supervision. A persistent model of the dwelling supports two complementary operations: grounding instructions for physical action, and retrieving remembered evidence to answer queries. New observations update the model, including after an action changes the environment. Figure~\ref{fig:p4_s13_1_managing} separates these action and recall branches around their shared home memory.

\textbf{Representing the home.} Everything downstream depends on a queryable model of the space, which extends the scene-understanding methods of Chapter~\ref{chapter:3:sec:6} into a persistent, semantically indexed memory. The memory architectures themselves are treated in Section~\ref{subsec:p3_s7_3_6}; here they are the infrastructure that lets an assistant know not only \emph{where} surfaces are, but \emph{what} and \emph{whose}. 
The 3D scene graph \cite{armeni20193dsg} introduced a hierarchical structure that unifies objects, rooms, and cameras with the underlying geometry, and Hydra \cite{hughes2022hydra} brought such graphs into real-time, loop-closure-optimized robot perception so that the memory stays current as the assistant moves. 
To make the memory addressable by free-form language, ConceptFusion \cite{jatavallabhula2023conceptfusion} builds an open-set, multimodal 3D map that answers text, image, or click queries, and ConceptGraphs \cite{gu2024conceptgraphs} fuses foundation-model outputs across views into a captioned, language-queryable 3D scene graph suitable for planning. 
This open-vocabulary memory is what lets an assistant act on objects it was never explicitly trained to recognize.


\begin{figure}[!ht]
    \centering
    \includegraphics[width=0.98\linewidth]{figures/part4/section13/assistant_managing.pdf}
    \caption{Managing the physical environment through a shared home memory. Observations maintain room, object, and relation records; grounded instructions drive manipulation, while queries independently retrieve last-seen evidence. Physical outcomes are re-observed to update memory.}
    \label{fig:p4_s13_1_managing}
\end{figure}

\textbf{Acting on the space: tidy-up and rearrangement.} Restoring order is the canonical management task, framed by Batra \etal \cite{batra2020rearrangement} as bringing an environment from an arbitrary configuration into a target spatial state. 
The benchmark substrate includes the Visual Room Rearrangement (RoomR) challenge \cite{weihs2021visual}, which asks an agent to observe a goal arrangement and then reset one to five displaced objects, and Housekeep \cite{kant2022housekeep}, which removes the explicit goal and instead judges placements against human preferences over nearly two thousand objects. 
Method work has moved from fixed rules to learned commonsense priors. 
TIDEE \cite{sarch2022tidee} tidies a never-before-seen room from pixels using a visuo-semantic out-of-place detector and an associative memory of where receptacles tend to be, and TidyBot \cite{wu2023tidybot} makes the placement \emph{preference-aware}, using few-shot summarization with a large language model to generalize a particular user's tidying rules from a handful of examples. 
The shift these systems mark, from a universal notion of tidiness to one conditioned on the resident, is exactly the personalization property that defines the assistant.

\textbf{Acting on the space: open-vocabulary mobile manipulation.} Tidying a real home requires picking and placing arbitrary objects at arbitrary locations, the task formalized as open-vocabulary mobile manipulation by HomeRobot \cite{yenamandra2023homerobot}, which pairs a simulation and real benchmark with a software stack for the Hello Robot Stretch. 
Notable deployment evidence comes from OK-Robot \cite{liu2024okrobot}, which combines a phone-scanned dense vision-language map with off-the-shelf navigation and grasping to achieve zero-shot pick-and-drop across ten real homes, and from Mobile ALOHA \cite{fu2024mobilealoha}, which shows that low-cost whole-body teleoperation plus imitation learning can acquire dexterous bimanual household skills. 
These systems lean directly on the manipulation primitives of Section~\ref{sec:p4_s9_1}, increasingly instantiated as generalist vision-language-action policies (\eg, OpenVLA \cite{kim2024openvla}, $\pi_0$ \cite{black2024pi0}) that fold perception, grounding, and control into one model; the assistant contribution is to drive them in unstructured homes from open-vocabulary goals rather than in an engineered cell. 
Commercial and pre-commercial examples illustrate the same pressure from different directions: mapping vacuum robots such as iRobot's Roomba, which build semantic room maps for targeted cleaning,\footnote{\url{https://www.irobot.com}} represent the mature narrow end, while emerging home-humanoid platforms such as 1X's Neo\footnote{\url{https://www.1x.tech/neo}} and Figure\footnote{\url{https://www.figure.ai}} show how the field is attempting, so far without a general-purpose result, to generalize mobile manipulation into domestic settings.

\textbf{Searching and recalling the model.} The dual of acting on the space is querying it: an assistant is, among other things, a search engine for the physical world that answers ``where did I leave my keys?'' Object-goal navigation supplies the active half, with goal-oriented semantic exploration (SemExp) \cite{chaplot2020object} building an episodic semantic map and a policy that exploits object co-occurrence to find a target efficiently rather than searching blindly. 
Embodied question answering \cite{das2018embodied} couples navigation with answering, and its open-vocabulary successor OpenEQA \cite{majumdar2024openeqa} spans both an episodic-memory setting, in which the assistant must recall what it has already seen, and an active-exploration setting, over more than a hundred real environments. 
For wearable assistants that record the user's day, the same recall problem becomes natural-language search of long egocentric video, the episodic-memory task defined by Ego4D \cite{grauman2022ego4d}. 
This is the application-facing counterpart of the memory-augmented reasoning systems in Section~\ref{subsec:p3_s7_3_6}.



\section{Pivotal Task-2: Co-habiting and Collaborating}
\label{subsec:p4-sec13-1-pivotal-task-2-co-habiting-and-collaborating}

As the distance between assistant and human shrinks, spatial intelligence must become social intelligence. The assistant must model the human as a co-located agent with intent, viewpoint, and comfort zones, not as a moving obstacle or a bare command source, and act predictably and collaboratively. Running through every part of this task is the reconciliation of reference frames: the assistant's own egocentric view, the human's egocentric view, and a shared allocentric map must be kept aligned.

\begin{figure}[!ht]
    \centering
    \includegraphics[width=0.98\linewidth]{figures/part4/section13/assistant_collaborating.pdf}
    \caption{Collaborating in shared space: (a) person- and context-dependent comfort zones and (b) aligned robot/human views inform an interaction-aware plan; (c) navigation and handover illustrate alternative behaviors. Dashed rings indicate comfort regions, and dashed sight lines associate views with shared objects.}
    \label{fig:p4_s13_1_collaborating}
\end{figure}

\textbf{Socially-aware navigation.} Moving through a populated space is not collision avoidance; it is adherence to the unwritten contracts of proxemics. The classical account is the social force model \cite{helbing1995social}, which casts pedestrian motion as attraction to a goal under repulsion from others, operationalizing Hall's anthropological zones of personal space. 
Data-driven successors learn the interaction directly: Social LSTM \cite{alahi2016social} couples per-person recurrent predictors through a social-pooling layer to forecast jointly plausible trajectories, and deep reinforcement-learning planners such as the attention-based crowd-robot model \cite{chen2019crowdrobot} weigh neighbors to navigate dense crowds while avoiding the ``freezing robot'' failure that purely reactive planners suffer \cite{trautman2010unfreezing}. 
Progress is measured on human-demonstration corpora such as SCAND \cite{karnan2022socially}, which records hours of socially compliant teleoperation for imitation and inverse reinforcement learning. 
The trend is toward anticipatory planning, in which the assistant infers the human's intended path and yields proactively, and toward learnable cost functions that adapt passing distances to individual or cultural preference rather than enforcing one universal margin.

\textbf{Grounded instruction following and spatial dialogue.} A collaborating assistant must resolve spatial language against what it sees, including the deictic references (``put it next to the \emph{other} bowl,'' ``the one over \emph{there}'') that pervade human speech. 
Vision-and-language navigation, introduced with the Room-to-Room benchmark \cite{anderson2018vision}, established the problem of following natural-language route instructions in photorealistic buildings; ALFRED \cite{shridhar2020alfred} extended it to long-horizon household tasks with irreversible state changes, and TEACh \cite{Padmakumar2022TEACh} cast the task as a two-agent dialogue in which a commander and follower must build shared spatial reference over multiple turns. 
Resolving a referring expression to a specific object in a cluttered 3D scene is studied directly by ScanRefer \cite{chen2020scanrefer} and ReferIt3D \cite{achlioptas2020referit3d}, whose same-class distractors force the model to rely on spatial relations rather than appearance. 
These benchmarks build on the reasoning capabilities of Chapter~\ref{sec:reasoning}; the assistant requirement is to maintain a spatial working memory that links current observations to earlier instructions so that deixis can be resolved.

\textbf{Shared autonomy and collaboration.} When full autonomy is brittle, the assistant acts as a third hand under shared control. 
The defining formulation is shared autonomy via hindsight optimization \cite{javdani2018shared}, in which the robot infers a distribution over the user's spatial goal and assists toward it without committing to a single confident prediction, formalizing the earlier policy-blending account of arbitration between human and robot intent \cite{dragan2013policy}. 
Equally important is making the assistant's own motion readable: legible motion \cite{dragan2013legibility} shapes a reaching trajectory so a human can quickly infer its intended target, which is what allows close-quarters collaboration to feel safe. Physical handovers, where giver and receiver must coordinate in space and time \cite{strabala2013handover}, are the elementary collaborative act on which assistive manipulation is built.

\textbf{Perspective-taking and egocentric grounding.} Many assistants observe the world from a wearable, first-person viewpoint, so the frame-alignment problem becomes explicit. 
Ego4D \cite{grauman2022ego4d} and its successors anchor egocentric perception at scale, and dedicated capture platforms such as Project Aria\footnote{\url{https://www.projectaria.com}} supply the calibrated multi-sensor first-person stream such assistants run on, while spatial-VQA backbones such as SpatialVLM \cite{chen2024spatialvlm} give a vision-language model the quantitative 3D reasoning (distances, sizes, directions) needed to answer spatial questions about the user's scene. 
The hardware and pose-estimation side of wearable sensing is treated in Section~\ref{sec:p4_s9_3}, and the augmented-reality runtime that renders guidance into the user's view is treated in Section~\ref{sec:p4_s10_1}; here the contribution is the reasoning that converts a first-person observation into a collaborative action in the human's frame. 
Embodied-AR applications make this alignment problem explicit when directions, warnings, or manipulation cues must be registered into the user's view \cite{yang2025embodiedar}. Figure~\ref{fig:p4_s13_1_collaborating} illustrates the elements of this regime.



\section{Pivotal Task-3: Empowering Vulnerable Users}
\label{subsec:p4-sec13-1-pivotal-task-3-empowering-vulnerable-users}

The highest-impact regime uses spatial intelligence to compensate for a sensory, cognitive, or motor impairment, acting as a prosthesis for the user's perception or agency. The assistant translates the environment into a queryable, navigable, and actionable representation and delivers it through a channel the user can still use (audio, haptics, or physical action). Because the loop closes around a vulnerable person in their own home, safety, privacy, and personalization are design constraints from the start.

\begin{figure}[!ht]
    \centering
    \includegraphics[width=0.98\linewidth]{figures/part4/section13/assistant_empowering.pdf}
    \caption{Personalized assistance combines observed spatial state with user goals and access needs. Consent, confidence, and clearance checks gate audio/haptic cues, stepwise guidance, and shared control; uncertainty or missing permission prompts clarification or a pause. Shared control also receives the user's command. Scenes are illustrative.}
    \label{fig:p4_s13_1_empowering}
\end{figure}

\textbf{Assistance for the blind and visually impaired.} For blind and low-vision users, spatial intelligence turns an unknown and hazardous environment into a queryable one. The oldest strand is sensory substitution, exemplified by the vOICe \cite{meijer1992voice}, which sonifies the spatial structure of an image into sound. 
The modern strand is conversational visual question answering: VizWiz \cite{bigham2010vizwiz} introduced a deployed system in which a blind user photographs their surroundings and asks a spoken question, and the VizWiz Grand Challenge \cite{gurari2018vizwiz} turned the resulting real-world distribution, including blurry and unanswerable images, into the reference benchmark for assistive VQA. 
Indoor wayfinding is addressed by smartphone navigation aids such as NavCog \cite{ahmetovic2016navcog}, which deliver turn-by-turn audio guidance from accurate beacon-based localization, and by three-dimensional spatial audio that places points of interest around the user. 
Commercial assistive applications include Microsoft's Seeing AI for scene and text narration\footnote{\url{https://www.seeingai.com}}, Be My Eyes for volunteer- and model-mediated visual assistance\footnote{\url{https://www.bemyeyes.com}}, and wearable systems from OrCam\footnote{\url{https://www.orcam.com}}, Envision\footnote{\url{https://www.letsenvision.com}}, and biped.ai\footnote{\url{https://biped.ai}} for scene description and obstacle warning. 
These tools move from passive obstacle warning to active interrogation, in which the user asks ``is there a free seat?'' or ``read the sign on the door'' and the assistant retrieves the answer.

\textbf{Cognitive and mobility assistance.} For users with cognitive decline or limited mobility, spatial assistants provide external scaffolding for the activities of daily living. 
Ambient assisted living systems such as the CASAS smart home \cite{cook2013casas} instrument a dwelling to track the resident's location and activity, learning the spatio-temporal patterns whose violation (a fall, a missed meal) signals a need to intervene; wearable accelerometry supplies the canonical fall-detection paradigm \cite{bourke2007evaluation}.

Where monitoring gives way to physical action, the assistant becomes a shared-control effector. 
Smart wheelchairs \cite{simpson2005smart} and wheelchair-mounted assistive arms use collaborative control to enforce spatial safety zones, for example forbidding a chair from driving down a staircase, so that the user retains mobility without constant supervision; the goal-arbitration that underlies such control is the same shared-autonomy machinery introduced above, here cast as a decision-theoretic problem \cite{ghorbel2018decision}. 
Assistive feeding \cite{bhattacharjee2019towards} shows how fine spatial reasoning (where to acquire a bite, how to transfer it to the user's mouth) becomes a caregiving act. 
The platform lineage runs from Robots for Humanity \cite{chen2013robots}, which let a man with quadriplegia perform self-care through a teleoperated PR2, to the affordable Stretch mobile manipulator \cite{kemp2022design}, with physics-based simulation of caregiving tasks consolidated in Assistive Gym \cite{erickson2020assistive}. 
Commercial assistive hardware in this space includes the Kinova JACO arm\footnote{\url{https://www.kinovarobotics.com}} and the Obi feeding device\footnote{\url{https://meetobi.com}}.
This regime overlaps the clinical setting of Section~\ref{sec:p4_s12_1}, but the boundary is clear: there the focus is diagnosis and surgical intervention, whereas here it is daily living in the home. Figure~\ref{fig:p4_s13_1_empowering} summarizes the compensation principle.



\begin{table}[t]
\centering
\fontsize{8}{10}\selectfont
\setlength{\tabcolsep}{1.0mm}
\caption{Representative benchmarks across the three assistant regimes. Where a benchmark has no standard method, the method column names the benchmark paper's own baselines.}
\label{tab:p4_s13_1_benchmarks}
\begin{tabular}{L{0.20\textwidth} L{0.24\textwidth} L{0.26\textwidth} L{0.22\textwidth}}
\toprule
\textbf{Regime / Pivot} & \textbf{Benchmark} & \textbf{Representative method} & \textbf{Headline metric} \\
\midrule
Managing (rearrangement) & Visual Room Rearrangement \cite{weihs2021visual} & TIDEE \cite{sarch2022tidee} & objects restored / episode success \\
Managing (preference tidy) & Housekeep \cite{kant2022housekeep} & LLM-based planner (paper baseline) \cite{kant2022housekeep} & correct-placement rate \\
Managing (mobile manip.) & HomeRobot OVMM \cite{yenamandra2023homerobot} & modular and RL baselines \cite{yenamandra2023homerobot} & overall success rate \\
Managing (search / recall) & EQA \cite{das2018embodied} / OpenEQA \cite{majumdar2024openeqa} & VLM with episodic memory (paper baselines) \cite{majumdar2024openeqa} & LLM-Match score \\
Collaborating (navigation) & SCAND \cite{karnan2022socially} & imitation from demonstrations \cite{karnan2022socially} & demonstration corpus (no leaderboard) \\
Collaborating (instruction) & ALFRED \cite{shridhar2020alfred} / TEACh \cite{Padmakumar2022TEACh} & seq2seq baselines \cite{shridhar2020alfred,Padmakumar2022TEACh} & task / goal-condition success \\
Collaborating (grounding) & ScanRefer \cite{chen2020scanrefer} / ReferIt3D \cite{achlioptas2020referit3d} & ScanRefer baseline \cite{chen2020scanrefer} & Acc@0.25 / 0.5 IoU \\
Empowering (BVI VQA) & VizWiz \cite{gurari2018vizwiz} & VQA baselines \cite{gurari2018vizwiz} & VQA accuracy \\
Empowering (assistive ADL) & Assistive Gym \cite{erickson2020assistive} & RL policies (paper baselines) \cite{erickson2020assistive} & task success / reward \\
\bottomrule
\end{tabular}
\end{table}

\section{Established Tools and Systems}
\label{subsec:p4-sec13-1-established-tools-and-systems}
Deployed assistants have so far approached the persistent, personalized spatial model of this chapter from opposite ends: software agents know a great deal about the user but nothing about the user's space, while home robots occupy the space but act with limited autonomy.

\textbf{Personal agents.} The most widely used assistants now act on a user's behalf in the digital world. OpenAI's ChatGPT agent\footnote{\url{https://openai.com/index/introducing-chatgpt-agent/}} (July 2025) runs on its own virtual computer, browsing and filling in forms across multi-step tasks and asking permission before consequential actions. Meta's Muse\footnote{\href{https://www.cnbc.com/2026/09/08/meta-personal-ai-agents-public-reckoning-privacy-safety.html}{CNBC, 8 September 2026}.} (September 2026), billed as a personal agent for everyday users, likewise runs in an isolated virtual machine and handles errands such as shopping and travel booking, and its request for access to email, calendars, and payments made privacy the main point of public debate around its launch. Voice assistants are being rebuilt on the same models: Amazon's Alexa+\footnote{\url{https://www.aboutamazon.com}}, generally available in the United States since February 2026, adds generative dialogue to the smart-home control that Alexa already provided, and Apple's Siri\footnote{\url{https://www.apple.com/siri/}} is being rebuilt on Google's Gemini models, adding personal context and on-screen awareness. For spatial intelligence, these agents matter mainly as the layer that holds personal context and preference; apart from smart-home control, none of them maintains a model of the physical home.

\textbf{Camera-grounded assistants.} A second line brings the physical world into the conversation through a camera. Gemini Live's camera mode\footnote{\url{https://gemini.google/}}, the product form of Google's Project Astra, lets a user point a phone at the surroundings and ask about what it sees, and Meta AI on camera glasses does the same from the wearer's viewpoint (Section~\ref{sec:p4_s9_3}). These systems answer spatial questions about the current view, and the assistive applications of Pivotal Task-3 build on the same capability, but they keep at most a short history of what they have seen, so the long-horizon spatial memory this chapter calls for remains outside them.

\textbf{Home robots.} Embodied assistants remain the least mature tier. Beyond mapping vacuums, the home humanoids introduced under Pivotal Task-1 are moving toward sale: 1X opened preorders for NEO in October 2025, at \$20{,}000 or a monthly subscription, with first deliveries promised for 2026 and with remote operators teleoperating tasks that its onboard autonomy cannot yet complete, and Figure's third-generation humanoid, Figure~03 (2025), is designed for homes and driven by its Helix vision-language-action model. The reliance on teleoperation shows how far these products are from the autonomous, open-vocabulary mobile manipulation of Pivotal Task-1, and it raises the in-home privacy question discussed below.

\section{Challenges}
\label{subsec:p4-sec13-1-challenges}

The three regimes share a common obstacle: the assistant must be reliable in the one setting that resists engineering, namely a real person's private and unstructured life. 
The challenges below are ordered from the perceptual frontier (generalization, memory) through the human frontier (personalization, social safety) to the deployment frontier (privacy, compute, trust), and each cuts across all three pivotal tasks above. 
Governance benchmarks for embodied agents have begun to measure the deployment end directly, testing whether a system stays controllable, recoverable, and auditable under perturbation rather than only whether it completes tasks \cite{qin2026embodiedgovbench}.

\textbf{Generalization to unstructured homes.} An assistant cannot be retrained per household, yet homes vary without bound in layout, lighting, and contents. 
Zero-shot results such as OK-Robot \cite{liu2024okrobot} are encouraging but remain far from the reliability a daily aide requires, and closing the gap depends on the open-vocabulary representations and sim-to-real transfer surveyed in Part~III rather than on more in-domain data.

\textbf{Long-horizon spatial memory.} Because the assistant persists in one place, it must remember spatial state over days and weeks: where objects were last seen, which routines the user follows, what was already done. 
Real-time scene-graph memories \cite{hughes2022hydra} and episodic-memory benchmarks \cite{majumdar2024openeqa} begin to address this, but the general memory mechanisms are treated in Section~\ref{subsec:p3_s7_3_6}; the assistant-specific unsolved problem is maintaining a consistent, updatable, and forgetting-resistant model over long deployments in a private home.

\textbf{Personalization and preference learning.} The assistant's geometric goals are user-specific, so it must learn an individual's placement rules, comfort distances, and routines from sparse interaction. 
Few-shot preference summarization \cite{wu2023tidybot} and learnable social cost functions are early steps, but acquiring rich preferences without burdening the user remains an open problem.

\textbf{Social compliance and safety.} Operating in continuous contact with people turns motion into a social and safety-critical act. 
Anticipatory, legible navigation and motion \cite{trautman2010unfreezing,dragan2013legibility} and shared-control safety zones \cite{ghorbel2018decision} mitigate risk, yet guaranteeing comfort and physical safety across the diversity of real human behavior is far from settled.

\textbf{Privacy of persistent in-home perception.} An assistant that continuously perceives a private home raises acute privacy concerns, the more so for assistive cameras worn by users who cannot themselves audit what is captured. 
The concern becomes sharper when assistance depends on remote operation or cloud annotation, because the raw spatial stream itself may reveal the user's routines and vulnerabilities. Persistent spatial memory must be reconciled with selective perception, on-device processing, and user control over what is retained.

\textbf{On-device compute and trust.} Wearable and home platforms are power- and thermally constrained, forcing the heavy perception and reasoning stacks of Part~III onto efficient, often on-device, inference. 
Compounding this, assistance for vulnerable users is safety-critical, so calibrated reliability and appropriate user trust, neither over- nor under-reliance, are prerequisites for deployment rather than refinements of it.


\newpage
\chapter{Spatial Intelligence for Others}
\label{sec:p4_s14}

\FramedInitial{O}{ther} application domains of spatial intelligence are too numerous to enumerate, because spatial intelligence is a horizontal capability: any field that must perceive, reconstruct, localize within, or reason about the geometry of the physical world inherits the methods of Part~III. 

This chapter takes five further domains where that inheritance is already visible in deployed systems. Each stresses a capability the earlier chapters do not: registering a building site against its design, tracking many near-identical agents through occlusion, reconstructing fragile objects faithfully, working without clean light or satellite positioning underwater, and binding people to products across a network of cameras.

\section{Architecture, Engineering, and Construction}
\label{sec:part-4-section-14-architecture-engineering-and-construction}

The architecture, engineering, and construction (AEC) sector is built around a single spatial-intelligence problem: continuously registering the messy as-built reality of a site against the structured as-designed model that specifies what it should become. 
This is the same Building-Information-Modeling (BIM) substrate that anchors the city-scale digital twins of Section~\ref{sec:p4_s11_1}, but here the model is the working blueprint of an individual building under construction rather than a synchronized replica of a finished one, and the central task is detecting and acting on the discrepancy between plan and reality.
Figure~\ref{fig:p4_s14_aec} sketches this as-built versus as-designed loop. \textbf{Scan-to-BIM} is its perception front end: dense laser or photogrammetric scans of a structure are semantically parsed into walls, floors, columns, and components, building on the indoor point-cloud understanding of Chapter~\ref{chapter:3:sec:6} and its benchmarks S3DIS \cite{armeni2016s3dis} and ScanNet \cite{dai2017scannet}. 
A specialized sub-problem is vectorizing a scan into an editable floorplan, where synthetic supervision such as Structured3D \cite{zheng2020structured3d} and transformer models that predict all room polygons in parallel \cite{yue2023roomformer} now recover clean geometry from cluttered scans. 
\textbf{Progress monitoring} closes the loop in time: by registering an image-based or scanned as-built point cloud against a four-dimensional IFC-BIM model, deviations between planned and actual construction state can be recognized automatically \cite{golparvarfard2015automated}, the spatiotemporal core of the commercial reality-capture platforms now common on job sites. 
\textbf{Construction robotics} acts on the model directly, with mobile fabricators that localize against the design to print, place, or assemble structural elements at full scale \cite{giftthaler2017mobile}. 
Finally, \textbf{structural inspection} turns the same perception stack toward maintenance, localizing cracks and defects on concrete surfaces from imagery \cite{cha2017deep}, increasingly from a UAV-borne 3D reconstruction of a bridge or facade. Reality-capture and BIM tooling from Matterport,\footnote{\url{https://matterport.com/}} Autodesk,\footnote{\url{https://www.autodesk.com/products/recap/overview}} and Leica,\footnote{\url{https://leica-geosystems.com/products/laser-scanners}} progress-monitoring services such as OpenSpace\footnote{\url{https://www.openspace.ai/}} and Buildots,\footnote{\url{https://buildots.com/}} and on-site robots from Dusty Robotics\footnote{\url{https://www.dustyrobotics.com/}} and Built Robotics\footnote{\url{https://www.builtrobotics.com/}} are representative of the deployed layer.

\begin{figure}[t]
    \centering
    \includegraphics[width=0.98\linewidth]{figures/part4/section14/other_aec.pdf}
    \caption{AEC tasks around an as-designed BIM: (1) semantic parsing of site scans; (2) comparison with the scheduled design state, highlighting a planned element not observed in the scan; (3) model-guided robotic placement; and (4) crack localization linked to a model element. Physical work prompts re-capture of the site. Scenes and point clouds are illustrative.}
    \label{fig:p4_s14_aec}
\end{figure}

\section{Sports Analytics and Performance}
\label{sec:part-4-section-14-sports-analytics-and-performance}

Sports analytics lifts raw broadcast or arena video into a calibrated metric model of the playing surface, then reasons over the resulting positional field. What sets it apart is fine-grained, multi-agent spatiotemporal tracking under severe and persistent occlusion, where many near-identical players and a small fast ball must be followed and assigned consistent identities at sub-second resolution.
Figure~\ref{fig:p4_s14_sports} illustrates three linked capabilities: registered trajectories feed tactical analysis, while calibrated multi-view pose supports biomechanics. \textbf{Tracking and registration} form the geometric foundation: multi-object tracking benchmarks such as SoccerNet-Tracking \cite{cioppa2022soccernet} and the multi-sport SportsMOT \cite{cui2023sportsmot} drive player and ball association, while sports-field registration estimates the homography from the broadcast frame to a metric pitch template so that pixel positions become field coordinates \cite{jiang2020optimizing}. 
\textbf{3D pose and biomechanics} recover athlete motion in three dimensions, with markerless capture systems estimating not only kinematics but joint dynamics from ordinary video, as in OpenCap \cite{uhlrich2023opencap}, typically by triangulating multi-view 2D pose into a 3D skeleton \cite{iskakov2019learnable}. 
\textbf{Tactical and team-space analysis} then treats the positional field itself as the object of study, valuing the space each player controls and forecasting where players will move; deep models of expected possession value \cite{fernandez2019decomposing} and multi-agent trajectory generation \cite{yeh2019diverse} build directly on the spatial-reasoning methods of Chapter~\ref{sec:reasoning}. 
The deployed layer is mature, spanning optical tracking and officiating from Hawk-Eye,\footnote{\url{https://www.hawkeyeinnovations.com/}} player-tracking data from Second Spectrum (now part of Genius Sports)\footnote{\url{https://www.geniussports.com/}} and Stats Perform,\footnote{\url{https://www.statsperform.com/}} and wearable spatial sensing from KINEXON\footnote{\url{https://kinexon-sports.com/}} and Catapult.\footnote{\url{https://www.catapult.com/}}

\begin{figure}[t]
    \centering
    \includegraphics[width=0.98\linewidth]{figures/part4/section14/other_sports.pdf}
    \caption{Sports analytics: (1) player/ball tracking and planar registration; (2) calibrated multi-view pose, kinematics, and model-estimated loads; and (3) tactical analysis from registered trajectories. The tactical view enlarges the dashed pitch region, preserving player IDs. Shading illustrates team control, dotted lines candidate passes, and dashed arrows predicted motion. Scenes and predictions are schematic.}
    \label{fig:p4_s14_sports}
\end{figure}


\begin{figure}[!t]
    \centering
    \includegraphics[width=0.98\linewidth]{figures/part4/section14/other_cultural_heritage.pdf}
    \caption{Three cultural-heritage tasks: (1) multi-view capture and surface reconstruction; (2) geometric and appearance-based matching of fragments for virtual reassembly; and (3) classification of airborne-LiDAR ground returns to reveal candidate structures in bare-earth relief. The tasks are independent; scenes, mesh, and terrain are illustrative.}
    \label{fig:p4_s14_heritage}
\end{figure}


\section{Cultural Heritage and Archaeology}
\label{sec:part-4-section-14-cultural-heritage-and-archaeology}

Cultural heritage and archaeology apply spatial intelligence to the past, running from capture through recovery to discovery. What distinguishes it is high-fidelity, metrically accurate, and usually non-contact reconstruction, coupled to a preservation imperative: the artifact may be fragile, fragmentary, or buried, and the digital record must be faithful enough to stand in for an object that cannot be handled or that no longer exists intact.
Figure~\ref{fig:p4_s14_heritage} organizes the field into three movements. 
\textbf{Digitization} captures artifacts and monuments in three dimensions, a lineage that runs from the landmark in-field laser scanning of the Digital Michelangelo Project \cite{levoy2000digital} through structure-from-motion photogrammetry \cite{schonberger2016structure} to the neural radiance fields \cite{mildenhall2021nerf} and 3D Gaussian splatting \cite{kerbl20233d} now used for photorealistic heritage capture, all drawing on the reconstruction methods of Chapter~\ref{sec:generation}. 
\textbf{Virtual restoration} uses geometric AI to reassemble what time has broken, posing fragment reassembly as a 3D puzzle over fractured parts, with large benchmarks such as Breaking Bad \cite{sellan2022breaking} and the real fresco-reassembly dataset RePAIR \cite{tsesmelis2024repair} driving the methods. 
\textbf{Archaeological prospection} discovers sites at landscape scale, most strikingly through airborne LiDAR whose bare-earth models reveal structures hidden under forest canopy, as in the canopy-penetrating surveys of the Maya lowlands at Caracol \cite{chase2011airborne} and across northern Guatemala \cite{canuto2018ancient}. Organizations such as CyArk\footnote{\url{https://www.cyark.org/}} and the Factum Foundation,\footnote{\url{https://www.factumfoundation.org/}} together with photogrammetry tools including Agisoft Metashape\footnote{\url{https://www.agisoft.com/}} and the Open Heritage 3D archive,\footnote{\url{https://openheritage3d.org/}} anchor the practical layer.



\begin{figure}[t]
    \centering
    \includegraphics[width=0.98\linewidth]{figures/part4/section14/other_marine.pdf}
    \caption{Underwater spatial intelligence under light attenuation/scattering and unavailable submerged GPS fixes. Enhanced visual observations combine with acoustic, inertial, and depth inputs for localization and mapping. Terrain and structure maps support coverage and standoff planning for hull/pipeline inspection; new views close the loop. Blue marks the survey track and selected defect. Scenes and enhancement are illustrative.}
    \label{fig:p4_s14_marine}
\end{figure}

\section{Marine and Underwater Systems}
\label{sec:part-4-section-14-marine-and-underwater-systems}

Underwater spatial intelligence confronts a setting in which the two assumptions that anchor terrestrial methods both fail at once: light no longer propagates cleanly, since wavelength-dependent attenuation and scattering degrade vision with range, and satellite positioning does not penetrate water, so absolute position must be inferred from acoustics, inertia, or the structure being surveyed.
Figure~\ref{fig:p4_s14_marine} depicts the resulting perceive-localize-map-inspect loop. 
\textbf{Navigation and SLAM} solve the GPS-denied localization problem, surveyed comprehensively for autonomous underwater vehicles by Paull \etal \cite{paull2014auv}, with foundational deep-sea visual SLAM demonstrated on the wreck of the Titanic \cite{eustice2005visually} and modern visual-inertial-pressure localization benchmarked by datasets such as AQUALOC \cite{ferrera2019aqualoc}. 
\textbf{Perception and enhancement} contend with the optical medium itself: physically grounded models remove the water from an image to recover true scene color and range \cite{akkaynak2019sea}, learned enhancement restores contrast in real time for downstream perception \cite{islam2020fast}, and benchmarks such as UIEB standardize the evaluation \cite{li2020underwater}, while large annotated image databases like FathomNet \cite{katija2022fathomnet} enable ocean-scale detection. 
\textbf{Mapping and inspection} reconstruct the seafloor, habitats, and submerged structures, with structure-relative perception and coverage planning enabling autonomous inspection of complex ship hulls \cite{hover2012advanced}. The operational layer spans low-cost platforms such as the BlueROV2,\footnote{\url{https://bluerobotics.com/store/rov/bluerov2/}} survey vehicles from Kongsberg\footnote{\url{https://www.kongsberg.com/discovery/autonomous-and-uncrewed-solutions/auv/hugin/}} and research institutions including WHOI\footnote{\url{https://www.whoi.edu/}} and MBARI.\footnote{\url{https://www.mbari.org/}}




\begin{figure}[t]
    \centering
    \includegraphics[width=0.98\linewidth]{figures/part4/section14/other_retail.pdf}
    \caption{Retail and smart-space intelligence across calibrated ceiling cameras: (1) shopper--product association updates a virtual basket; (2) detected shelf contents are compared with a planogram; (3) ground-plane tracks yield shopper-flow and visit-density maps. Dotted lines link scene details; the geometry and density are illustrative.}
    \label{fig:p4_s14_retail}
\end{figure}

\section{Retail and Smart Spaces}
\label{sec:part-4-section-14-retail-and-smart-spaces}

Retail and smart-space intelligence reduces to one persistent spatial problem: jointly associating, across a network of calibrated overlapping cameras and over time, which person is interacting with which product at which shelf location. 
Unlike the warehouse logistics of Section~\ref{sec:p4_s9_5}, where the environment is engineered for the robot, here the space is a crowded store full of untracked people, and its distinguishing challenge is dense multi-object tracking with reliable human-object association through frequent mutual occlusion.
Figure~\ref{fig:p4_s14_retail} shows how the major applications share this core. 
\textbf{Multi-camera tracking} provides the spatial backbone, with benchmarks such as WILDTRACK \cite{chavdarova2018wildtrack} and MMPTRACK \cite{han2023mmptrack} and methods that project per-camera features onto a shared ground plane to detect and track people across views \cite{hou2020multiview}; this is the machinery behind cashierless checkout, where a shopper's path and pickups must be reconstructed without a checkout line. 
\textbf{Shelf monitoring} turns perception toward the products, detecting densely packed items for inventory and planogram compliance, supported by the dense-detection benchmark SKU-110K \cite{goldman2019precise} and by methods that formulate compliance as matching the observed shelf layout against the ideal one as a sub-graph isomorphism \cite{tonioni2017product}. 
\textbf{In-store analytics} aggregates the same tracks into shopper-flow heatmaps and layout insights, building on cross-camera re-identification \cite{zheng2015scalable} to stitch one customer's trajectory across non-overlapping views.
The deployed layer includes cashierless systems such as Amazon's Just Walk Out\footnote{\url{https://www.justwalkout.com/}} and platforms from Trigo,\footnote{\url{https://www.trigoretail.com/}} AiFi,\footnote{\url{https://www.aifi.com/}} and Zippin,\footnote{\url{https://www.getzippin.com/}} alongside shelf-analytics vendors such as Trax.\footnote{\url{https://traxretail.com/}}


\begin{table}[t]
\centering
\fontsize{8}{10}\selectfont
\setlength{\tabcolsep}{1.0mm}
\caption{Representative benchmarks for the five domains of this chapter.}
\label{tab:p4_s14_benchmarks}
\begin{tabular}{L{0.16\textwidth} L{0.30\textwidth} L{0.28\textwidth} L{0.20\textwidth}}
\toprule
\textbf{Domain} & \textbf{Benchmark / dataset} & \textbf{Representative method} & \textbf{Headline metric} \\
\midrule
AEC & S3DIS \cite{armeni2016s3dis} & point-cloud semantic segmentation & mIoU \\
AEC & Structured3D \cite{zheng2020structured3d} & RoomFormer \cite{yue2023roomformer} & room / corner / angle F1 \\
Sports & SoccerNet-Tracking \cite{cioppa2022soccernet} / SportsMOT \cite{cui2023sportsmot} & multi-object trackers & HOTA \\
Sports & broadcast field registration & learned-error optimization \cite{jiang2020optimizing} & IoU (whole / part) \\
Heritage & Breaking Bad \cite{sellan2022breaking} / RePAIR \cite{tsesmelis2024repair} & learned fragment reassembly & rotation / translation error \\
Marine & UIEB \cite{li2020underwater} & Water-Net \cite{li2020underwater} & PSNR / SSIM \\
Marine & AQUALOC \cite{ferrera2019aqualoc} & visual-inertial-pressure SLAM & trajectory error \\
Retail & WILDTRACK \cite{chavdarova2018wildtrack} & MVDet \cite{hou2020multiview} & MODA \\
Retail & SKU-110K \cite{goldman2019precise} & dense detector with EM-merger \cite{goldman2019precise} & AP \\
\bottomrule
\end{tabular}
\end{table}



\section{Outlook}
\label{sec:part-4-section-14-outlook}

The five domains above, whose benchmarks are collected in Table~\ref{tab:p4_s14_benchmarks}, are a sample, not a census. 
Spatial intelligence is equally at work in fields this survey does not treat at length: virtual try-on and 3D e-commerce, immersive education and high-fidelity training simulators, security and surveillance, energy and infrastructure inspection, precision agriculture beyond the remote-sensing view of Section~\ref{sec:p4_s11_2}, and many others.
What unites them, and what unites them with the headline applications of the preceding chapters, is that each takes the perception, reasoning, and generation capabilities of Part~III and specializes them to the geometry, sensing modality, and constraints of a particular slice of the physical world.
